% Blueprint for "Fundamental Limitations on Subquadratic Alternatives to
% Transformers" by Josh Alman and Hantao Yu, ICLR 2025 / arXiv:2410.04271v2.
%
% Library facts are leaves with \mathlibok. Paper-specific CS definitions,
% reductions, hardness statements, and transformer constructions are novel nodes.

\chapter{Basic notation and transformer model}

\begin{definition}[Finite real vectors and matrices]
  \label{def:finite-real-arrays}
  \lean{Fin}
  \mathlibok
  For a positive integer $d$, a real vector of length $d$ is represented as a
  function $\{1,\ldots,d\}\to\R$, and an $n\times d$ real matrix is represented
  as a function $\{1,\ldots,n\}\times\{1,\ldots,d\}\to\R$.  The notation
  $v[i]$, $A_{i,:}$, and $A_{:,j}$ denotes coordinate projection, row projection,
  and column projection.
\end{definition}

\begin{definition}[Inner product and Euclidean norm]
  \label{def:inner-norm}
  \uses{def:finite-real-arrays}
  \lean{inner}
  \mathlibok
  For $v,w\in\R^d$, $\ip{v}{w}=\sum_i v[i]w[i]$ and
  $\|v\|=(\ip{v}{v})^{1/2}$.
\end{definition}

\begin{lemma}[Cauchy--Schwarz inequality]
  \label{lem:cauchy-schwarz}
  \uses{def:inner-norm}
  \lean{abs_inner_le_norm}
  \mathlibok
  For real finite-dimensional vectors $v,w$, $|\ip{v}{w}|\leq\|v\|\,\|w\|$.
\end{lemma}

\begin{definition}[Binary vector operations]
  \label{def:binary-vector-ops}
  \uses{def:finite-real-arrays}
  A binary vector is an element of $\{0,1\}^d$.  The all-one vector is
  $\mathbf{1}_d$.  If $v\in\{0,1\}^d$, then $\bar v\in\{0,1\}^d$ is defined by
  $\bar v[i]=1-v[i]$ for every coordinate $i$.
\end{definition}

\begin{lemma}[Complement flips binary inner products]
  \label{lem:binary-complement-inner}
  \uses{def:binary-vector-ops, def:inner-norm}
  If $v,w\in\{0,1\}^d$ and $\|v\|_1=\sum_i v[i]$, then
  $\ip{v}{\bar w}=\|v\|_1-\ip{v}{w}$.
\end{lemma}
\begin{proof}
  \uses{def:binary-vector-ops, def:inner-norm}
  By definition of the complement,
  $\ip{v}{\bar w}=\sum_i v[i](1-w[i])$.  Distributing the product and separating
  the finite sum gives $\sum_i v[i]-\sum_i v[i]w[i]=\|v\|_1-\ip{v}{w}$.
\end{proof}

\begin{definition}[Kronecker product]
  \label{def:kron-product}
  \lean{def_kron_product}
  \uses{def:finite-real-arrays}
  For $v\in\R^a$ and $w\in\R^b$, the Kronecker product $v\otimes w\in\R^{ab}$
  has coordinates indexed by pairs and satisfies $(v\otimes w)[i,j]=v[i]w[j]$.
\end{definition}

\begin{definition}[Kronecker power coordinate product]
  \label{def:kron-power-coordinate-product}
  \lean{def_kron_power_coordinate_product}
  \uses{def:kron-product, def:finite-real-arrays}
  For a vector $v\in\R^d$ and a positive integer $q$, the $q$-fold Kronecker
  power $v^{\otimes q}\in\R^{d^q}$ is the vector whose coordinates are indexed
  by functions $\alpha:\{1,\ldots,q\}\to\{1,\ldots,d\}$ and whose coordinate
  value is the concrete product
  \[
    v^{\otimes q}[\alpha]=\prod_{r=1}^{q} v[\alpha(r)].
  \]
  Equivalently, $v^{\otimes 1}=v$ under the one-coordinate identification and
  $v^{\otimes(q+1)}=v\otimes v^{\otimes q}$ with the coordinate equation of
  Definition~\ref{def:kron-product}.  This node defines only the power
  construction and its coordinate product relation.
\end{definition}

\begin{definition}[Vector concatenation]
  \label{def:kron-concat}
  \lean{def_vector_concat}
  \uses{def:finite-real-arrays}
  For $v\in\R^a$ and $w\in\R^b$, the concatenation
  $v\circ w\in\R^{a+b}$ is the vector with coordinates
  \[
    (v\circ w)[r]=v[r]\quad(1\leq r\leq a),
    \qquad
    (v\circ w)[a+s]=w[s]\quad(1\leq s\leq b).
  \]
  This node defines only concatenation; Kronecker products and Kronecker powers
  are Definition~\ref{def:kron-product} and
  Definition~\ref{def:kron-power-coordinate-product}.
\end{definition}

\begin{lemma}[Kronecker inner product]
  \label{lem:kron-inner}
  \uses{def:kron-power-coordinate-product, def:inner-norm}
  For finite real vectors $v,w$ and every positive integer $q$,
  \[
    \ip{v^{\otimes q}}{w^{\otimes q}}=\ip{v}{w}^{q}.
  \]
\end{lemma}
\begin{proof}
  \uses{def:kron-power-coordinate-product, def:inner-norm}
  For $q=1$ this is the definition.  The step from $q$ to $q+1$ expands the
  inner product over pairs of coordinates:
  $\sum_{i,\alpha} v[i]v^{\otimes q}[\alpha]w[i]w^{\otimes q}[\alpha]$.
  Separating the product of finite sums gives
  $\ip{v}{w}\ip{v^{\otimes q}}{w^{\otimes q}}$, which is
  $\ip{v}{w}^{q+1}$ by the induction hypothesis.
\end{proof}

\begin{lemma}[Kronecker norm]
  \label{lem:kron-norm}
  \uses{lem:kron-inner, def:inner-norm}
  For every finite real vector $v$ and positive integer $q$,
  $\|v^{\otimes q}\|=\|v\|^q$.
\end{lemma}
\begin{proof}
  \uses{lem:kron-inner, def:inner-norm}
  Squaring both sides reduces the claim to
  $\ip{v^{\otimes q}}{v^{\otimes q}}=\ip{v}{v}^{q}$, which is
  Lemma~\ref{lem:kron-inner} with $w=v$.  Both norms are nonnegative, so equality
  of squares implies equality of norms.
\end{proof}

\begin{definition}[Softmax on vectors]
  \label{def:softmax-vector}
  \uses{def:finite-real-arrays}
  For $x\in\R^n$,
  \[
    \softmax(x)=\frac{(\exp(x[1]),\ldots,\exp(x[n]))}{\sum_{i=1}^n\exp(x[i])}.
  \]
\end{definition}

\begin{definition}[Row-wise softmax on matrices]
  \label{def:softmax-matrix}
  \uses{def:softmax-vector, def:finite-real-arrays}
  For $A\in\R^{n\times n}$, $\softmax(A)$ is defined row-wise by
  $\softmax(A)_{i,:}=\softmax(A_{i,:})$.
\end{definition}

\begin{definition}[Document-keyword occurrence relation]
  \label{def:document-keyword-occurrence}
  \lean{def_document_keyword_occurrence}
  A document-similarity instance fixes a type ${\cal D}$ of documents, a type
  ${\cal K}$ of keywords, and a concrete occurrence predicate
  \[
    \Occurs(D,k)
  \]
  for $D\in{\cal D}$ and $k\in{\cal K}$, meaning that keyword $k$ occurs in
  document $D$.  The predicate is part of the input model as decidable data:
  for every pair $(D,k)$ there is a yes/no occurrence decision whose true
  branch is exactly $\Occurs(D,k)$ and whose false branch is exactly
  $\neg\Occurs(D,k)$.

  This occurrence datum contains no length parameter and no built-in keyword
  list.  The Lean-facing structure
  $\operatorname{def\_document\_keyword\_occurrence}$ has exactly the following
  four fields and field names:
  \[
    \mathrm{Document}:\mathrm{Type},\qquad
    \mathrm{Keyword}:\mathrm{Type},\qquad
    \mathrm{Occurs}:\mathrm{Document}\to\mathrm{Keyword}\to\mathrm{Prop},
  \]
  and
  \[
    \mathrm{occurs\_decidable}:
      \forall D:\mathrm{Document},\ \forall k:\mathrm{Keyword},\
      \operatorname{Decidable}(\mathrm{Occurs}(D,k)).
  \]
  Thus the tuple is
  $({\cal D},{\cal K},\Occurs,\operatorname{occurs\_decidable})$ with the
  displayed names.  The keyword vector used for a particular bag-of-words
  embedding is supplied separately in
  Definition~\ref{def:keyword-coordinate-vector}; it is not a field of this
  occurrence structure.
\end{definition}

\begin{definition}[Coordinate-indexed keyword vector]
  \label{def:keyword-coordinate-vector}
  \lean{def_keyword_coordinate_vector}
  \uses{def:document-keyword-occurrence, def:finite-real-arrays}
  For an occurrence datum
  $\mathsf{doc}:\operatorname{def\_document\_keyword\_occurrence}$ and a chosen
  bag-of-words length $\ell$, a keyword list is the concrete coordinate-indexed
  function
  \[
    K:\{1,\ldots,\ell\}\to\mathsf{doc}.\mathrm{Keyword}.
  \]
  Formally,
  \[
    \operatorname{def\_keyword\_coordinate\_vector}(\mathsf{doc},\ell)
      =(\{1,\ldots,\ell\}\to\mathsf{doc}.\mathrm{Keyword}).
  \]
  This definition has exactly the displayed inputs $(\mathsf{doc},\ell)$ and
  returns a type of keyword-vector parameters.  It does not construct a default
  keyword vector, erase the occurrence datum, use an ambient keyword family, or
  store the vector in $\mathsf{doc}$.  The $r$th coordinate of a document
  embedding refers to the concrete keyword $K(r)$.  Two instances with the same
  documents and occurrence predicate but different functions $K$ are different
  bag-of-words inputs and may produce different embedding matrices.
\end{definition}

\begin{definition}[Real-valued keyword occurrence indicator]
  \label{def:keyword-occurrence-indicator}
  \lean{def_keyword_occurrence_indicator}
  \uses{def:document-keyword-occurrence}
  For an occurrence datum
  $\mathsf{doc}:\operatorname{def\_document\_keyword\_occurrence}$, a document
  $D:\mathsf{doc}.\mathrm{Document}$, and a keyword
  $k:\mathsf{doc}.\mathrm{Keyword}$, define the real-valued occurrence
  indicator
  \[
    \operatorname{OccVal}_{\mathsf{doc}}(D,k)=
    \begin{cases}
      1\in\R, & \mathsf{doc}.\mathrm{Occurs}(D,k),\\
      0\in\R, & \neg\mathsf{doc}.\mathrm{Occurs}(D,k).
    \end{cases}
  \]
  Lean-facing, this is the function
  \[
    \operatorname{def\_keyword\_occurrence\_indicator}
      (\mathsf{doc})(D)(k):\R
      =
      \begin{cases}
        1, & \mathsf{doc}.\mathrm{Occurs}(D,k),\\
        0, & \neg\mathsf{doc}.\mathrm{Occurs}(D,k),
      \end{cases}
  \]
  where the branch is selected by the concrete
  $\mathsf{doc}.\mathrm{occurs\_decidable}(D,k)$ field.  No other document
  encoding or keyword semantics are implicit in this value.
\end{definition}

\begin{lemma}[Present keyword occurrence indicator]
  \label{lem:keyword-occurrence-indicator-present}
  \uses{def:keyword-occurrence-indicator}
  For every occurrence datum $\mathsf{doc}$, document
  $D:\mathsf{doc}.\mathrm{Document}$, and keyword
  $k:\mathsf{doc}.\mathrm{Keyword}$,
  \[
    \mathsf{doc}.\mathrm{Occurs}(D,k)\Rightarrow
    \operatorname{OccVal}_{\mathsf{doc}}(D,k)=1.
  \]
\end{lemma}
\begin{proof}
  \uses{def:keyword-occurrence-indicator}
  Under the hypothesis $\mathsf{doc}.\mathrm{Occurs}(D,k)$, the defining
  occurrence decision for $\operatorname{OccVal}_{\mathsf{doc}}(D,k)$ takes
  its true branch, whose value is $1$.
\end{proof}

\begin{lemma}[Absent keyword occurrence indicator]
  \label{lem:keyword-occurrence-indicator-absent}
  \uses{def:keyword-occurrence-indicator}
  For every occurrence datum $\mathsf{doc}$, document
  $D:\mathsf{doc}.\mathrm{Document}$, and keyword
  $k:\mathsf{doc}.\mathrm{Keyword}$,
  \[
    \neg\mathsf{doc}.\mathrm{Occurs}(D,k)\Rightarrow
    \operatorname{OccVal}_{\mathsf{doc}}(D,k)=0.
  \]
\end{lemma}
\begin{proof}
  \uses{def:keyword-occurrence-indicator}
  Under the hypothesis $\neg\mathsf{doc}.\mathrm{Occurs}(D,k)$, the defining
  occurrence decision for $\operatorname{OccVal}_{\mathsf{doc}}(D,k)$ takes
  its false branch, whose value is $0$.
\end{proof}

\begin{definition}[Bag-of-words embedding]
  \label{def:bag-of-words}
  \lean{def_bag_of_words}
  \uses{def:keyword-coordinate-vector, def:keyword-occurrence-indicator}
  For an occurrence datum
  $\mathsf{doc}:\operatorname{def\_document\_keyword\_occurrence}$ and a length
  $\ell$, the Lean-facing bag-of-words embedding is the function
  \[
    \operatorname{def\_bag\_of\_words}(\mathsf{doc},\ell):
      \mathsf{doc}.\mathrm{Document}\to
      \operatorname{def\_keyword\_coordinate\_vector}(\mathsf{doc},\ell)\to
      (\{1,\ldots,\ell\}\to\R)
  \]
  defined by the coordinate equation
  \[
    \operatorname{def\_bag\_of\_words}(\mathsf{doc},\ell)(D)(K)(r)
      =
      \operatorname{def\_keyword\_occurrence\_indicator}
        (\mathsf{doc})(D)(K(r)).
  \]
  We write this output as
  $\BOW_{\mathsf{doc}}(D,K):\{1,\ldots,\ell\}\to\R$.  The keyword vector $K$ is
  an explicit argument of this definition, not a field of the occurrence datum,
  and the codomain is the real vector space $\R^\ell$ used by the paper rather
  than a Boolean-valued coordinate type.
\end{definition}

\begin{lemma}[Bag-of-words coordinate is occurrence indicator]
  \label{lem:bow-coordinate-indicator}
  \uses{def:bag-of-words, def:keyword-occurrence-indicator}
  For every occurrence datum $\mathsf{doc}$, document
  $D:\mathsf{doc}.\mathrm{Document}$, coordinate-indexed keyword vector
  $K:\{1,\ldots,\ell\}\to\mathsf{doc}.\mathrm{Keyword}$, and coordinate $r$,
  \[
    \operatorname{def\_bag\_of\_words}(\mathsf{doc},\ell)(D)(K)(r)
      =\operatorname{OccVal}_{\mathsf{doc}}(D,K(r)).
  \]
\end{lemma}
\begin{proof}
  \uses{def:bag-of-words}
  This is the coordinate equality in the definition of
  the bag-of-words row
  $\operatorname{def\_bag\_of\_words}(\mathsf{doc},\ell)(D)(K)$ at coordinate
  $r$.  There is no separate fixed-keyword wrapper declaration: the keyword
  vector $K$ is the explicit second argument of
  $\operatorname{def\_bag\_of\_words}(\mathsf{doc},\ell)(D)$.
\end{proof}

\begin{lemma}[Bag-of-words present keyword coordinate]
  \label{lem:bow-coordinate-present}
  \uses{def:bag-of-words, lem:keyword-occurrence-indicator-present}
  For every occurrence datum $\mathsf{doc}$, document
  $D:\mathsf{doc}.\mathrm{Document}$, coordinate-indexed keyword vector
  $K:\{1,\ldots,\ell\}\to\mathsf{doc}.\mathrm{Keyword}$, and coordinate $r$,
  \[
    \mathsf{doc}.\mathrm{Occurs}(D,K(r))
      \Rightarrow
      \operatorname{def\_bag\_of\_words}(\mathsf{doc},\ell)(D)(K)(r)=1.
  \]
\end{lemma}
\begin{proof}
  \uses{def:bag-of-words, lem:keyword-occurrence-indicator-present}
  Expanding Definition~\ref{def:bag-of-words} rewrites the displayed
  bag-of-words coordinate as the occurrence indicator
  $\operatorname{OccVal}_{\mathsf{doc}}(D,K(r))$.  Applying
  Lemma~\ref{lem:keyword-occurrence-indicator-present} to the concrete keyword
  $K(r)$ and the hypothesis $\mathsf{doc}.\mathrm{Occurs}(D,K(r))$ gives that
  this indicator value is $1$.
\end{proof}

\begin{lemma}[Bag-of-words absent keyword coordinate]
  \label{lem:bow-coordinate-absent}
  \uses{def:bag-of-words, lem:keyword-occurrence-indicator-absent}
  For every occurrence datum $\mathsf{doc}$, document
  $D:\mathsf{doc}.\mathrm{Document}$, coordinate-indexed keyword vector
  $K:\{1,\ldots,\ell\}\to\mathsf{doc}.\mathrm{Keyword}$, and coordinate $r$,
  \[
    \neg\mathsf{doc}.\mathrm{Occurs}(D,K(r))
      \Rightarrow
      \operatorname{def\_bag\_of\_words}(\mathsf{doc},\ell)(D)(K)(r)=0.
  \]
\end{lemma}
\begin{proof}
  \uses{def:bag-of-words, lem:keyword-occurrence-indicator-absent}
  Expanding Definition~\ref{def:bag-of-words} rewrites the displayed
  bag-of-words coordinate as the occurrence indicator
  $\operatorname{OccVal}_{\mathsf{doc}}(D,K(r))$.  Applying
  Lemma~\ref{lem:keyword-occurrence-indicator-absent} to the concrete keyword
  $K(r)$ and the hypothesis $\neg\mathsf{doc}.\mathrm{Occurs}(D,K(r))$ gives
  that this indicator value is $0$.
\end{proof}

\begin{lemma}[Bag-of-words coordinates are binary]
  \label{lem:bow-coordinate-binary}
  \uses{def:bag-of-words, def:document-keyword-occurrence}
  For every occurrence datum $\mathsf{doc}$, document
  $D:\mathsf{doc}.\mathrm{Document}$, coordinate-indexed keyword vector
  $K:\{1,\ldots,\ell\}\to\mathsf{doc}.\mathrm{Keyword}$, and coordinate $r$,
  \[
    \operatorname{def\_bag\_of\_words}(\mathsf{doc},\ell)(D)(K)(r)=0
    \quad\text{or}\quad
    \operatorname{def\_bag\_of\_words}(\mathsf{doc},\ell)(D)(K)(r)=1.
  \]
\end{lemma}
\begin{proof}
  \uses{def:bag-of-words, def:keyword-occurrence-indicator, lem:keyword-occurrence-indicator-present, lem:keyword-occurrence-indicator-absent, def:document-keyword-occurrence}
  Expand Definition~\ref{def:bag-of-words} at the coordinate $r$, reducing the
  displayed value to the occurrence indicator
  $\operatorname{OccVal}_{\mathsf{doc}}(D,K(r))$.  The occurrence predicate is
  decidable for the pair $(D,K(r))$.  If
  $\mathsf{doc}.\mathrm{Occurs}(D,K(r))$, then
  Lemma~\ref{lem:keyword-occurrence-indicator-present} gives the value $1$.  If
  $\neg\mathsf{doc}.\mathrm{Occurs}(D,K(r))$, then
  Lemma~\ref{lem:keyword-occurrence-indicator-absent} gives the value $0$.
\end{proof}

\begin{definition}[Document row equals its bag-of-words embedding]
  \label{def:document-row-bow-equality}
  \uses{def:bag-of-words, def:keyword-coordinate-vector, def:finite-real-arrays}
  For an occurrence datum $\mathsf{doc}$, a document
  $D:\mathsf{doc}.\mathrm{Document}$, a coordinate-indexed keyword vector
  $K:\{1,\ldots,\ell\}\to\mathsf{doc}.\mathrm{Keyword}$, and a row vector
  $e\in\R^\ell$, the predicate
  \[
    \RowBOW_{\mathsf{doc}}(D,K,e)
  \]
  holds exactly when $e$ is the concrete bag-of-words embedding of $D$ with
  respect to $K$:
  \[
    \forall\,1\leq r\leq\ell,\qquad
    e[r]=\operatorname{def\_bag\_of\_words}(\mathsf{doc},\ell)(D)(K)(r).
  \]
  This predicate has exactly the displayed inputs $(\mathsf{doc},D,K,e)$.  It
  has no matrix index, no occurrence-value expansion, no nonzero-row promise,
  and no access to an ambient keyword vector stored in the occurrence datum.
\end{definition}

\begin{definition}[Matrix row as a vector]
  \label{def:matrix-row-vector}
  \uses{def:finite-real-arrays}
  For a matrix $E\in\R^{n\times\ell}$ and an index $i\in\{1,\ldots,n\}$,
  the row vector
  \[
    \operatorname{RowVec}(E,i):\{1,\ldots,\ell\}\to\R
  \]
  is defined coordinatewise by
  \[
    \operatorname{RowVec}(E,i)(r)=E[i,r]
    \qquad\text{for every }r\in\{1,\ldots,\ell\}.
  \]
  This is only the concrete row-projection function from a finite matrix to a
  finite vector.
\end{definition}

\begin{definition}[Valid document-similarity input matrix]
  \label{def:valid-document-input-matrix}
  \uses{def:bag-of-words, def:keyword-coordinate-vector, def:finite-real-arrays}
  For an occurrence datum $\mathsf{doc}$, a document family
  $D:\{1,\ldots,n\}\to\mathsf{doc}.\mathrm{Document}$, a coordinate-indexed
  keyword vector $K:\{1,\ldots,\ell\}\to\mathsf{doc}.\mathrm{Keyword}$, and a
  matrix $E\in\R^{n\times \ell}$, define
  \[
    \ValidBOW_{\mathsf{doc}}(D,K,E)
  \]
  to be the concrete matrix-entry relation
  \[
    \forall\,i\in\{1,\ldots,n\},\ \forall\,r\in\{1,\ldots,\ell\},\qquad
    E[i,r]=
      \operatorname{def\_bag\_of\_words}(\mathsf{doc},\ell)(D(i))(K)(r).
  \]
  Lean-facing, the document input is the single finite-indexed function $D$,
  not a variadic tuple of fields.  The predicate has exactly the displayed
  inputs $(\mathsf{doc},D,K,E)$, with $K$ explicit.  It contains no nonzero-row
  hypothesis and no occurrence-value expansion; those are separate consequences
  proved below.  Thus a valid document-similarity input is a real matrix whose
  entries are precisely the bag-of-words coordinates for the displayed document
  family and keyword vector.  Consequently, any later proof needing an entry
  equality specializes this definition directly at the relevant row and
  coordinate; Lemma~\ref{lem:valid-bow-entry-equality} records that
  specialization as a separate one-line formal interface.  There is no
  alternate rowwise validity predicate.
\end{definition}

\begin{lemma}[Valid document matrix entry is its bag-of-words coordinate]
  \label{lem:valid-bow-entry-equality}
  \uses{def:valid-document-input-matrix, def:bag-of-words}
  If $\ValidBOW_{\mathsf{doc}}(D,K,E)$ holds, then for every
  $i\in\{1,\ldots,n\}$ and $r\in\{1,\ldots,\ell\}$,
  \[
    E[i,r]=
      \operatorname{def\_bag\_of\_words}(\mathsf{doc},\ell)(D(i))(K)(r).
  \]
\end{lemma}
\begin{proof}
  \uses{def:valid-document-input-matrix}
  This is exactly Definition~\ref{def:valid-document-input-matrix}
  specialized to the displayed row $i$ and coordinate $r$.
\end{proof}

\begin{lemma}[Document row bag-of-words coordinate is occurrence value]
  \label{lem:document-row-bow-coordinate-occurrence-value}
  \uses{def:bag-of-words, def:keyword-occurrence-indicator}
  For all positive integers $n$ and $\ell$, every occurrence datum
  $\mathsf{doc}$, document family
  $D:\{1,\ldots,n\}\to\mathsf{doc}.\mathrm{Document}$,
  coordinate-indexed keyword vector
  $K:\{1,\ldots,\ell\}\to\mathsf{doc}.\mathrm{Keyword}$, row
  $i\in\{1,\ldots,n\}$, and coordinate $r\in\{1,\ldots,\ell\}$,
  \[
    \operatorname{def\_bag\_of\_words}(\mathsf{doc},\ell)(D(i))(K)(r)
      =
      \operatorname{def\_keyword\_occurrence\_indicator}
        (\mathsf{doc})(D(i))(K(r)).
  \]
\end{lemma}
\begin{proof}
  \uses{def:bag-of-words, def:keyword-occurrence-indicator}
  Expand Definition~\ref{def:bag-of-words} at the concrete document $D(i)$,
  keyword vector $K$, and coordinate $r$.  The right-hand side is exactly the
  Lean-facing occurrence-value function introduced in
  Definition~\ref{def:keyword-occurrence-indicator}.
\end{proof}

\begin{lemma}[Valid document matrix entries are occurrence values]
  \label{lem:valid-bow-entry-occurrence-value}
  \uses{lem:valid-bow-entry-equality, lem:document-row-bow-coordinate-occurrence-value}
  If $\ValidBOW_{\mathsf{doc}}(D,K,E)$ holds, then for every
  $i\in\{1,\ldots,n\}$ and $r\in\{1,\ldots,\ell\}$,
  \[
    E[i,r]=\operatorname{OccVal}_{\mathsf{doc}}(D(i),K(r)).
  \]
\end{lemma}
\begin{proof}
  \uses{lem:valid-bow-entry-equality, lem:document-row-bow-coordinate-occurrence-value}
  Lemma~\ref{lem:valid-bow-entry-equality} gives
  \[
    E[i,r]=
      \operatorname{def\_bag\_of\_words}(\mathsf{doc},\ell)(D(i))(K)(r).
  \]
  Lemma~\ref{lem:document-row-bow-coordinate-occurrence-value} gives
  \[
    \operatorname{def\_bag\_of\_words}(\mathsf{doc},\ell)(D(i))(K)(r)
      =\operatorname{OccVal}_{\mathsf{doc}}(D(i),K(r)).
  \]
  Transitivity of equality yields the displayed occurrence-value equation for
  $E[i,r]$.
\end{proof}

\begin{lemma}[Valid document matrix entry for present keyword]
  \label{lem:valid-bow-entry-occurrence-present}
  \uses{lem:valid-bow-entry-equality, lem:bow-coordinate-present}
  If $\ValidBOW_{\mathsf{doc}}(D,K,E)$ holds, then for every
  $i\in\{1,\ldots,n\}$ and $r\in\{1,\ldots,\ell\}$,
  \[
    \mathsf{doc}.\mathrm{Occurs}(D(i),K(r))\Rightarrow E[i,r]=1
  \]
\end{lemma}
\begin{proof}
  \uses{lem:valid-bow-entry-equality, lem:bow-coordinate-present}
  Lemma~\ref{lem:valid-bow-entry-equality} rewrites $E[i,r]$ as the
  bag-of-words coordinate of $D(i)$ at $r$.  Under the hypothesis
  $\mathsf{doc}.\mathrm{Occurs}(D(i),K(r))$,
  Lemma~\ref{lem:bow-coordinate-present} gives that same coordinate value as
  $1$.
\end{proof}

\begin{lemma}[Valid document matrix entry for absent keyword]
  \label{lem:valid-bow-entry-occurrence-absent}
  \uses{lem:valid-bow-entry-equality, lem:bow-coordinate-absent}
  If $\ValidBOW_{\mathsf{doc}}(D,K,E)$ holds, then for every
  $i\in\{1,\ldots,n\}$ and $r\in\{1,\ldots,\ell\}$,
  \[
    \neg\mathsf{doc}.\mathrm{Occurs}(D(i),K(r))\Rightarrow E[i,r]=0.
  \]
\end{lemma}
\begin{proof}
  \uses{lem:valid-bow-entry-equality, lem:bow-coordinate-absent}
  Lemma~\ref{lem:valid-bow-entry-equality} rewrites $E[i,r]$ as the
  bag-of-words coordinate of $D(i)$ at $r$.  Under the hypothesis
  $\neg\mathsf{doc}.\mathrm{Occurs}(D(i),K(r))$,
  Lemma~\ref{lem:bow-coordinate-absent} gives that same coordinate value as
  $0$.
\end{proof}

\begin{lemma}[Valid document matrix entry occurrence cases]
  \label{lem:valid-bow-entry-occurrence-cases}
  \uses{lem:valid-bow-entry-occurrence-present, lem:valid-bow-entry-occurrence-absent, def:document-keyword-occurrence}
  If $\ValidBOW_{\mathsf{doc}}(D,K,E)$ holds, then for every
  $i\in\{1,\ldots,n\}$ and $r\in\{1,\ldots,\ell\}$, exactly the concrete
  occurrence decision for $(D(i),K(r))$ determines the entry value:
  \[
    \bigl(\mathsf{doc}.\mathrm{Occurs}(D(i),K(r))\wedge E[i,r]=1\bigr)
    \quad\text{or}\quad
    \bigl(\neg\mathsf{doc}.\mathrm{Occurs}(D(i),K(r))\wedge E[i,r]=0\bigr).
  \]
\end{lemma}
\begin{proof}
  \uses{lem:valid-bow-entry-occurrence-present, lem:valid-bow-entry-occurrence-absent, def:document-keyword-occurrence}
  Use the decidability field of
  Definition~\ref{def:document-keyword-occurrence} for the proposition
  $\mathsf{doc}.\mathrm{Occurs}(D(i),K(r))$.  In the true case,
  Lemma~\ref{lem:valid-bow-entry-occurrence-present} supplies $E[i,r]=1$.
  In the false case, Lemma~\ref{lem:valid-bow-entry-occurrence-absent}
  supplies $E[i,r]=0$.
\end{proof}

\begin{definition}[Nonzero document-row promise]
  \label{def:nonzero-document-row-promise}
  \uses{def:bag-of-words, def:keyword-coordinate-vector, def:inner-norm}
  For an occurrence datum $\mathsf{doc}$, a document family
  $D:\{1,\ldots,n\}\to\mathsf{doc}.\mathrm{Document}$, keyword vector
  $K:\{1,\ldots,\ell\}\to\mathsf{doc}.\mathrm{Keyword}$, and matrix
  $E\in\R^{n\times\ell}$, the separate predicate
  \[
    \NonzeroRows_{\mathsf{doc}}(D,K,E)
  \]
  means that every document embedding row used in a cosine-similarity instance
  is nonzero:
  \[
    \forall\,i\in\{1,\ldots,n\},\qquad
    \|\operatorname{def\_bag\_of\_words}(\mathsf{doc},\ell)(D(i))(K)\|>0
  \]
  and no other condition.  The matrix argument $E$ is present only so that
  $\NonzeroRows_{\mathsf{doc}}(D,K,E)$ has the same displayed
  instance parameters as
  $\ValidBOW_{\mathsf{doc}}(D,K,E)$; the predicate itself does not
  inspect entries of $E$ and does not require a proof of $\ValidBOW$.  This is
  the standing promise for MSD and LSD instances, but it is not an argument of
  $\ValidBOW_{\mathsf{doc}}$.
\end{definition}

\begin{definition}[Cosine similarity]
  \label{def:cosine-similarity}
  \uses{def:inner-norm}
  For nonzero document embeddings $v,w\in\R^d$, their cosine similarity is
  \[
    \frac{\ip{v}{w}}{\|v\|\cdot\|w\|}.
  \]
\end{definition}

\begin{lemma}[Cosine similarity range for binary vectors]
  \label{lem:cosine-binary-range}
  \uses{def:cosine-similarity, lem:cauchy-schwarz}
  If $v,w\in\{0,1\}^d$ are nonzero binary vectors, then their cosine similarity
  lies in $[0,1]$.
\end{lemma}
\begin{proof}
  \uses{def:cosine-similarity, lem:cauchy-schwarz, def:inner-norm}
  The inner product is a sum of nonnegative products, so it is nonnegative.  The
  denominator is positive because both vectors are nonzero.  The upper bound is
  exactly Cauchy--Schwarz divided by this positive denominator.
\end{proof}

\begin{definition}[Attention, Definition 2.1]
  \label{def:attention}
  \uses{def:softmax-matrix, def:finite-real-arrays}
  For $d_{\rm in},d_{\rm out},m\in\N$, matrices
  $Q,K\in\R^{d_{\rm in}\times m}$ and $V\in\R^{d_{\rm in}\times d_{\rm out}}$
  define an attention map $A_{Q,K,V}:\R^{n\times d_{\rm in}}\to
  \R^{n\times d_{\rm out}}$ by
  \[
    A_{Q,K,V}(X)=\softmax(XQK^\top X^\top)XV.
  \]
\end{definition}

\begin{definition}[Family of attention maps]
  \label{def:attention-family}
  \uses{def:attention}
  The family $\mathcal A_{d_{\rm in},m,d_{\rm out}}$ consists of all maps
  $A_{Q,K,V}$ with $Q,K\in\R^{d_{\rm in}\times m}$ and
  $V\in\R^{d_{\rm in}\times d_{\rm out}}$.
\end{definition}

\begin{definition}[Multi-layer perceptron, Definition 2.2]
  \label{def:mlp}
  For positive integers $a$ and $b$, an $(a,b)$ multi-layer perceptron is a
  continuous function
  \[
    \varphi:(\{1,\ldots,a\}\to\R)\to(\{1,\ldots,b\}\to\R),
  \]
  equivalently a continuous map $\R^a\to\R^b$ on the corresponding
  finite-dimensional real vector spaces.  For every number of rows $n$, its
  row-wise matrix lift is the map
  \[
    \varphi^{[n]}:\R^{n\times a}\to\R^{n\times b}
  \]
  defined by the coordinate equality
  \[
    \bigl(\varphi^{[n]}(X)\bigr)_{i,:}=\varphi(X_{i,:})
    \qquad\text{for every }1\leq i\leq n.
  \]
  Thus the lifted output has entries
  $(\varphi^{[n]}(X))[i,j]=\varphi(X_{i,:})[j]$.
\end{definition}

\begin{definition}[Single-attention-unit transformer, Definition 2.3]
  \label{def:transformer}
  \uses{def:attention, def:mlp}
  Fix positive integers $n,d,d_{\rm in},d_{\rm out},m$.  A transformer in the
  paper's simplified single-attention-unit model is the following concrete
  data:
  an input MLP $\varphi_1:\R^d\to\R^{d_{\rm in}}$, matrices
  $Q,K\in\R^{d_{\rm in}\times m}$ and
  $V\in\R^{d_{\rm in}\times d_{\rm out}}$, and an output MLP
  \[
    \varphi_2:\R^{n d_{\rm out}}\to\R,
  \]
  where an $n\times d_{\rm out}$ matrix is identified with its flattened
  vector in $\R^{n d_{\rm out}}$.  Its value on a matrix input
  $E\in\R^{n\times d}$ is
  \[
    {\rm TF}(E)=
    \varphi_2\!\left(\operatorname{flat}
      \left(A_{Q,K,V}(\varphi_1^{[n]}(E))\right)\right).
  \]
  Equivalently, the computation first applies $\varphi_1$ row-wise, producing
  a matrix in $\R^{n\times d_{\rm in}}$; then applies the attention map
  $A_{Q,K,V}$ from Definition~\ref{def:attention}, producing a matrix in
  $\R^{n\times d_{\rm out}}$; then applies the continuous output MLP to the
  flattened finite-dimensional matrix.
\end{definition}

\begin{definition}[Document transformer input contract]
  \label{def:document-transformer-input-contract}
  \uses{def:transformer, def:valid-document-input-matrix, lem:valid-bow-entry-equality, lem:valid-bow-entry-occurrence-value, lem:valid-bow-entry-occurrence-present, lem:valid-bow-entry-occurrence-absent, lem:valid-bow-entry-occurrence-cases}
  For document-similarity tasks the transformer input dimension is exactly the
  bag-of-words length.  Concretely, the parameters satisfy $d=\ell$, the
  document data consist of an occurrence datum $\mathsf{doc}$, a document
  family $D:\{1,\ldots,n\}\to\mathsf{doc}.\mathrm{Document}$, and a coordinate-indexed
  keyword vector $K:\{1,\ldots,\ell\}\to\mathsf{doc}.\mathrm{Keyword}$.
  The Lean-facing contract is the single predicate
  \[
    \operatorname{def\_document\_transformer\_input\_contract}
      (\mathsf{doc},D,K,E) : \operatorname{Prop}
  \]
  defined by exactly the bag-of-words validity condition
  \[
    \operatorname{def\_document\_transformer\_input\_contract}
      (\mathsf{doc},D,K,E)
      \quad\Longleftrightarrow\quad
    \ValidBOW_{\mathsf{doc}}(D,K,E).
  \]
  It has precisely the displayed data arguments $(\mathsf{doc},D,K,E)$, with
  $n$ and $\ell$ supplied by the finite index types.  Its only input hypothesis
  is $\ValidBOW_{\mathsf{doc}}(D,K,E)$; it has no conjunct, premise, field, or
  hidden promise asserting that rows are nonzero.  In particular, the general
  transformer input contract must still hold for valid bag-of-words matrices
  with zero rows, because those are legitimate transformer inputs even though
  cosine similarity may be undefined for them.

  From this contract and Lemma~\ref{lem:valid-bow-entry-equality}, every
  entry $E[i,r]$ is the
  corresponding bag-of-words coordinate, and by
  Lemma~\ref{lem:valid-bow-entry-occurrence-value} it is the real-valued
  occurrence indicator
  $\operatorname{OccVal}_{\mathsf{doc}}(D(i),K(r))$.  By
  Lemma~\ref{lem:valid-bow-entry-occurrence-present}, present keywords give
  entry value $1$; by
  Lemma~\ref{lem:valid-bow-entry-occurrence-absent}, absent keywords give
  entry value $0$.  Lemma~\ref{lem:valid-bow-entry-occurrence-cases} records
  the complete decidable case split for each entry.
  The scalar answer to the document-similarity instance is the pure transformer
  value ${\rm TF}(E)$ from Definition~\ref{def:transformer} evaluated on this
  concretely determined bag-of-words matrix.  Thus the model never evaluates a
  document-similarity task on an unconstrained abstract row encoding, and the
  matrix column type of the transformer is the same finite coordinate type as
  the keyword list.
\end{definition}

\begin{definition}[Sentinel-augmented input matrix]
  \label{def:sentinel-augmented-input}
  \lean{def_sentinel_augmented_input}
  \uses{def:finite-real-arrays}
  For a base matrix $E\in\R^{n\times d}$ and a concrete sentinel row
  $s\in\R^d$, the sentinel augmentation data is the concrete finite object
  \[
    \operatorname{def\_sentinel\_augmented\_input}(E,s)
    =
    \bigl(\Aug_s(E),\operatorname{ProblemRow},
      \operatorname{AdmissiblePair},\operatorname{LiftDecision}\bigr)
  \]
  with the following four components.

  First, $\Aug_s(E)$ is the matrix
  \[
    \Aug_s(E)\in\R^{(n+1)\times d}
  \]
  defined by the row equalities
  \[
    \Aug_s(E)_{i,:}=E_{i,:}\quad(1\leq i\leq n),
    \qquad
    \Aug_s(E)_{n+1,:}=s.
  \]
  Second, the problem-row predicate on augmented row indices is
  \[
    \operatorname{ProblemRow}(r)\quad\Longleftrightarrow\quad
    1\leq r\leq n .
  \]
  Thus row $n+1$ is definitionally not a problem row.

  Third, the admissible witness-pair predicate on augmented row indices is
  \[
    \operatorname{AdmissiblePair}(p,q)
    \quad\Longleftrightarrow\quad
    1\leq p\leq n,\quad 1\leq q\leq n,\quad p\neq q .
  \]
  Consequently no pair containing the sentinel index $n+1$ is admissible.
  Explicitly, for every augmented index $r$,
  \[
    \neg\operatorname{AdmissiblePair}(n+1,r),
    \qquad
    \neg\operatorname{AdmissiblePair}(r,n+1).
  \]
  Every underlying document, OV, Max-IP, Min-IP, MSD, or LSD witness predicate
  is evaluated only on pairs satisfying this displayed predicate, so it ranges
  over exactly the original first $n$ rows.  In predicate form, for every base
  pair relation $W_E$ on indices of $E$,
  \[
    \exists p,q,\,
      \operatorname{AdmissiblePair}(p,q)\wedge W_E(p,q)
    \quad\Longleftrightarrow\quad
    \exists i,j\in\{1,\ldots,n\},\, i\neq j\wedge W_E(i,j),
  \]
  and
  \[
    \forall p,q,\,
      \operatorname{AdmissiblePair}(p,q)\Rightarrow W_E(p,q)
    \quad\Longleftrightarrow\quad
    \forall i,j\in\{1,\ldots,n\},\, i\neq j\Rightarrow W_E(i,j).
  \]

  Fourth, for any base matrix decision predicates
  $\Valid,\Yes,\No:\R^{n\times d}\to\operatorname{Prop}$, the lifted predicates
  attached to this augmentation are not new predicates on arbitrary
  $(n+1)\times d$ matrices.  They are the three concrete precompositions
  \[
    \Valid^{\Aug}(E,s)\Longleftrightarrow\Valid(E),\qquad
    \Yes^{\Aug}(E,s)\Longleftrightarrow\Yes(E),\qquad
    \No^{\Aug}(E,s)\Longleftrightarrow\No(E).
  \]
  Hence appending the sentinel preserves the original size-$n$ yes/no
  condition by definition: the augmented encoding is valid, yes, or no exactly
  when the base matrix $E$ is valid, yes, or no.  The sentinel row is only a
  fixed computational token supplied to the transformer construction; it is
  never part of the mathematical problem instance.
\end{definition}

\begin{definition}[Sentinel-augmented transformer solver]
  \label{def:sentinel-transformer-solves}
  \uses{def:transformer, def:sentinel-augmented-input}
  A sentinel-augmented transformer for a matrix decision problem with base
  input shape $n\times d$ consists of a fixed sentinel row $s\in\R^d$ and a
  single-attention-unit transformer whose sequence length parameter is
  $n+1$.  It solves predicates $\Valid(E)$, $\Yes(E)$, and $\No(E)$ on base
  matrices $E\in\R^{n\times d}$ when both concrete correctness implications
  hold after augmentation:
  \[
    \Valid(E)\wedge\Yes(E)\Rightarrow {\rm TF}(\Aug_s(E))=1,
    \qquad
    \Valid(E)\wedge\No(E)\Rightarrow {\rm TF}(\Aug_s(E))=0.
  \]
  Thus appending the sentinel is part of the fixed input encoding of the
  constructed transformer, while the mathematical decision problem remains the
  original problem on the first $n$ rows.
\end{definition}

\begin{definition}[Transformer solves a decision problem]
  \label{def:transformer-solves}
  \uses{def:transformer}
  A matrix decision problem with input shape $n\times d$ consists of predicates
  $\Valid(E)$, $\Yes(E)$, and $\No(E)$ on matrices $E\in\R^{n\times d}$ such
  that every valid instance is classified exactly once:
  \[
    \Valid(E)\Rightarrow
    \bigl((\Yes(E)\wedge\neg\No(E))\vee(\No(E)\wedge\neg\Yes(E))\bigr).
  \]
  A transformer ${\rm TF}:\R^{n\times d}\to\R$ solves this decision problem if
  both correctness implications hold for every matrix input:
  \[
    \Valid(E)\wedge\Yes(E)\Rightarrow {\rm TF}(E)=1,
    \qquad
    \Valid(E)\wedge\No(E)\Rightarrow {\rm TF}(E)=0.
  \]
  The no-instance implication is an independent requirement; it is not
  conditional on the existence of a yes-instance witness.
\end{definition}

\chapter{Fine-grained complexity and similarity problems}

\begin{definition}[Truly subquadratic time]
  \label{def:truly-subquadratic}
  A nonnegative running-time bound $T:\N\to\R$ is truly subquadratic in $n$ if
  there exist constants $\eps>0$, $C>0$, and $n_0\in\N$ such that for every
  $n\geq n_0$,
  \[
    T(n)\leq C n^{2-\eps}.
  \]
  When an auxiliary dimension $\ell$ is present, a bound
  $T:\N\times\N\to\R$ is truly subquadratic up to polynomial factors in $\ell$
  if there are $\eps>0$, $C>0$, an exponent $b\in\N$, and $n_0$ such that
  \[
    T(n,\ell)\leq C n^{2-\eps}(1+\ell)^b
  \]
  for all $n\geq n_0$ and all allowed $\ell$.
\end{definition}

\begin{definition}[Computational problem specification]
  \label{def:computational-problem-spec}
  A computational problem $P$ with principal size parameter $n$ and auxiliary
  dimension parameter $\ell$ consists of the following concrete data:
  an input type $I_P$, size functions
  \[
    N_P:I_P\to\N,\qquad L_P:I_P\to\N,
  \]
  a valid-input predicate $\Valid_P(x)$, an output type $O_P(x)$ for each
  input $x$, and a correctness relation
  \[
    \Correct_P(x,y)
  \]
  for $y\in O_P(x)$.  The relation $\Correct_P(x,y)$ is required to encode the
  real yes/no, optimization, or approximation condition of the problem, not
  merely that $y$ has the right syntactic type.  For a decision problem,
  $O_P(x)=\{0,1\}$ and correctness means output $1$ exactly on promised yes
  instances and output $0$ exactly on promised no instances.
\end{definition}

\begin{definition}[Deterministic operational semantics]
  \label{def:deterministic-operational-semantics}
  \uses{def:computational-problem-spec}
  Fix a computational problem $P$.  A deterministic computation model for
  algorithms solving $P$ consists of a type of program codes, and for each
  program code $A$ and input $x\in I_P$, a concrete configuration type
  $\operatorname{Config}_A(x)$ together with:
  an initial configuration
  $\operatorname{init}_A(x)\in\operatorname{Config}_A(x)$, a deterministic step
  map
  \[
    \operatorname{step}_A(x):
      \operatorname{Config}_A(x)\to\operatorname{Config}_A(x),
  \]
  a halting predicate $\operatorname{Halt}_A(x,c)$ on configurations, and an
  output relation
  \[
    \operatorname{Out}_A(x,c,y)
  \]
  between halted configurations $c$ and outputs $y\in O_P(x)$.  The run after
  $s$ steps is defined recursively by
  \[
    \operatorname{Run}_A(x,0)=\operatorname{init}_A(x),\qquad
    \operatorname{Run}_A(x,s+1)=
      \operatorname{step}_A(x)(\operatorname{Run}_A(x,s)).
  \]
  The statement
  \[
    \operatorname{HaltsWith}_A(x,s,y)
  \]
  means the run first halts at exactly step $s$ and returns $y$:
  \[
    \operatorname{Halt}_A(x,\operatorname{Run}_A(x,s))\wedge
    \operatorname{Out}_A(x,\operatorname{Run}_A(x,s),y)
  \]
  and, for every $r<s$,
  $\neg\operatorname{Halt}_A(x,\operatorname{Run}_A(x,r))$.  The output
  relation is functional on halted configurations, and the deterministic run
  has a unique first halting time:
  if $\operatorname{HaltsWith}_A(x,s,y)$ and
  $\operatorname{HaltsWith}_A(x,s',y')$, then
  $s=s'$ and $y=y'$.
  Every halted configuration has a unique output: if
  $\operatorname{Halt}_A(x,c)$, then there exists a unique
  $y\in O_P(x)$ such that $\operatorname{Out}_A(x,c,y)$.  Hence a halted
  configuration determines a concrete output value by unique description, not
  by an externally supplied oracle-answer parameter.

  This node is the operational semantics used below.  A step count is the
  finite number $s$ appearing in
  $\operatorname{HaltsWith}_A(x,s,y)$, not an arbitrary field attached to a
  total function.
\end{definition}

\begin{definition}[Correct algorithm for a problem]
  \label{def:problem-algorithm}
  \uses{def:computational-problem-spec, def:deterministic-operational-semantics}
  An algorithm $A$ for a computational problem $P$ is a program code in the
  deterministic operational semantics of
  Definition~\ref{def:deterministic-operational-semantics}.  Its execution on
  input $x$ is the concrete run $\operatorname{Run}_A(x,s)$ generated by the
  initial configuration and step map.  For a valid input $x$, the notation
  $A(x)$ denotes the unique output $y\in O_P(x)$ for which there exists a
  finite step count $s$ with $\operatorname{HaltsWith}_A(x,s,y)$; this notation
  is defined only from the halting relation and the uniqueness guaranteed by
  the operational semantics.

  The algorithm halts on promised inputs when, for every valid $x$, there
  exist $s\in\N$ and $y\in O_P(x)$ such that
  $\operatorname{HaltsWith}_A(x,s,y)$.  It is correct when
  \[
    \Valid_P(x)\Rightarrow
    \exists s\in\N,\ \exists y\in O_P(x),\
      \operatorname{HaltsWith}_A(x,s,y)\wedge \Correct_P(x,y)
  \]
  for every input $x$.  Its running-time witness is a function
  $T_A:\N\times\N\to\R_{\geq0}$ such that on every valid input $x$ there are
  $s\in\N$ and $y\in O_P(x)$ with
  \[
    \operatorname{HaltsWith}_A(x,s,y),\qquad
    \Correct_P(x,y),\qquad
    s\leq T_A(N_P(x),L_P(x)).
  \]
  Thus the claimed running time is semantically linked to the first halting
  time of the operational execution, and the returned value is linked to the
  problem's correctness relation.  The structure does not permit an
  unconstrained $\operatorname{computation\_steps}$ witness or an arbitrary
  total function unrelated to execution.
\end{definition}

\begin{definition}[Exact worst-case running-time bound]
  \label{def:exact-worst-case-running-time}
  \uses{def:problem-algorithm}
  A function $T:\N\times\N\to\R_{\geq0}$ is a worst-case running-time bound for
  an algorithm $A$ solving $P$ if for every valid input $x$ with
  $N_P(x)=n$ and $L_P(x)=\ell$, the execution of $A$ on $x$ halts within
  $T(n,\ell)$ steps and returns an output satisfying $\Correct_P(x,A(x))$.
  The bound is exact for a pair $(n,\ell)$ if it is the maximum of these
  running times over all valid inputs with those parameters, and it is an upper
  bound if it is at least that maximum.
\end{definition}

\begin{lemma}[Promised-input solvers halt correctly]
  \label{lem:solver-halts-on-promised-input-set}
  \uses{def:problem-algorithm, def:exact-worst-case-running-time}
  If $A$ is a correct algorithm for $P$ with worst-case running-time bound $T$,
  then for every valid promised input $x$, $A(x)$ halts within
  $T(N_P(x),L_P(x))$ steps and $\Correct_P(x,A(x))$ holds.
\end{lemma}
\begin{proof}
  \uses{def:problem-algorithm, def:exact-worst-case-running-time}
  This is the defining universal guarantee of a correct algorithm equipped with
  the stated worst-case running-time bound, specialized to the valid input
  $x$.
\end{proof}

\begin{definition}[Oracle reduction data]
  \label{def:oracle-reduction-data}
  \uses{def:computational-problem-spec}
  An oracle reduction datum $R$ from $P$ to $Q$ assigns to every valid
  $P$-input $x$ a finite number $r(x)$ of oracle calls, a sequence of concrete
  $Q$-inputs
  \[
    q_1(x),\ldots,q_{r(x)}(x)\in I_Q,
  \]
  and a postprocessing map which turns any sequence of oracle answers
  $z_s\in O_Q(q_s(x))$ into a final output
  \[
    y_R(x;z_1,\ldots,z_{r(x)})\in O_P(x).
  \]
  Adaptive reductions are represented by allowing the later query $q_s$ to be
  the query produced after the earlier answer values
  $z_1,\ldots,z_{s-1}$ have been fixed; the completed run is still the finite
  sequence of concrete query-answer pairs and the final output above.
\end{definition}

\begin{definition}[Oracle reduction trace]
  \label{def:oracle-reduction-trace}
  \uses{def:oracle-reduction-data}
  For an oracle reduction datum $R$ and a valid $P$-input $x$, an oracle
  reduction trace is the finite sequence
  \[
    (q_1,z_1),\ldots,(q_r,z_r),y.
  \]
  It is consistent with $R$ when $r=r(x)$, each $q_s$ is the query produced by
  $R$ at call $s$ on input $x$ after the earlier answers in the trace, each
  $z_s\in O_Q(q_s)$, and
  \[
    y=y_R(x;z_1,\ldots,z_r).
  \]
\end{definition}

\begin{definition}[Valid oracle queries]
  \label{def:valid-oracle-queries}
  \uses{def:oracle-reduction-trace}
  An oracle reduction datum $R$ sends valid $P$-inputs to valid $Q$-queries if
  for every valid $P$-input $x$ and every trace consistent with $R$ on $x$,
  each query in the trace satisfies $\Valid_Q(q_s)$.
\end{definition}

\begin{definition}[Correct oracle answers]
  \label{def:correct-oracle-answers}
  \uses{def:oracle-reduction-trace}
  In a trace for a reduction from $P$ to $Q$, the oracle answers are correct
  when every answer $z_s$ satisfies the target problem's correctness relation
  on the concrete query that produced it:
  \[
    \Correct_Q(q_s,z_s)
  \]
  for all oracle-call positions $s$.
\end{definition}

\begin{definition}[Oracle reduction validity and correctness obligations]
  \label{def:oracle-reduction-correctness}
  \uses{def:valid-oracle-queries, def:correct-oracle-answers}
  An oracle reduction datum $R$ from $P$ to $Q$ is correct if both of the
  following universal obligations hold.  First, it sends every valid $P$-input
  only to valid $Q$-queries in the sense of
  Definition~\ref{def:valid-oracle-queries}.  Second, for every valid
  $P$-input $x$ and every trace consistent with $R$ on $x$, if all oracle
  answers in that trace are correct for their concrete $Q$-queries, then the
  final output $y$ in the trace satisfies
  \[
    \Correct_P(x,y).
  \]
  This requirement is uniform over all possible correct oracle answers, so the
  reduced algorithm may not rely on a special choice among multiple valid
  answers returned by the oracle.
\end{definition}

\begin{definition}[Oracle reduction answer-history type]
  \label{def:oracle-reduction-answer-history}
  \uses{def:oracle-reduction-data}
  For an oracle reduction datum $R$, a valid $P$-input $x$, and a phase index
  $s\in\{0,\ldots,r(x)\}$, the answer-history type
  \[
    \Hist_R(x,s)
  \]
  is the concrete type of finite query-answer prefixes of length $s$:
  \[
    H=((q_1,z_1),\ldots,(q_s,z_s)).
  \]
  Such an $H$ is well formed when every $q_u$ is the $u$th query produced by
  the reduction datum $R$ on input $x$ after the earlier answers
  $z_1,\ldots,z_{u-1}$, and every $z_u\in O_Q(q_u)$.  Thus
  $\Hist_R(x,0)$ contains the empty history, while $\Hist_R(x,r(x))$ contains
  full oracle-answer histories for completed runs of $R$ on $x$.

  This history type is the only place in the local-phase presentation where
  raw oracle answers and query-answer pairs live.  It is passed as a separate
  argument to later query, update, and postprocessing functions.
\end{definition}

\begin{definition}[Oracle reduction non-oracle component types]
  \label{def:oracle-reduction-local-component-types}
  \uses{def:oracle-reduction-data}
  A local-phase presentation of an oracle reduction datum $R$ supplies one
  global, history-free record of component type families.  Lean-facing, the
  declaration
  \[
    \operatorname{def\_oracle\_reduction\_local\_component\_types}(P,Q,R)
  \]
  is a structure with exactly the following six named fields and no others:
  \[
  \begin{aligned}
    \operatorname{Control}&:
      \prod_{x:I_P}\prod_{h_x:\Valid_P(x)}
      \{0,\ldots,r(x)\}\to\operatorname{Type},\\
    \operatorname{CounterName}&:
      \prod_{x:I_P}\prod_{h_x:\Valid_P(x)}
      \{0,\ldots,r(x)\}\to\operatorname{Type},\\
    \operatorname{TapeName}&:
      \prod_{x:I_P}\prod_{h_x:\Valid_P(x)}
      \{0,\ldots,r(x)\}\to\operatorname{Type},\\
    \operatorname{TapeSymbol}&:
      \prod_{x:I_P}\prod_{h_x:\Valid_P(x)}
      \{0,\ldots,r(x)\}\to\operatorname{Type},\\
    \operatorname{RegisterName}&:
      \prod_{x:I_P}\prod_{h_x:\Valid_P(x)}
      \{0,\ldots,r(x)\}\to\operatorname{Type},\\
    \operatorname{RegisterValue}&:
      \prod_{x:I_P}\prod_{h_x:\Valid_P(x)}
      \{0,\ldots,r(x)\}\to\operatorname{Type}.
  \end{aligned}
  \]
  For a valid $P$-input $x$, proof $h_x:\Valid_P(x)$, and phase
  $s\in\{0,\ldots,r(x)\}$, the six instantiated history-free component types
  are therefore
  \[
    \operatorname{Control}_R(x,h_x,s),\quad
    \operatorname{CounterName}_R(x,h_x,s),\quad
    \operatorname{TapeName}_R(x,h_x,s),\quad
    \operatorname{TapeSymbol}_R(x,h_x,s),\quad
    \operatorname{RegisterName}_R(x,h_x,s),\quad
    \operatorname{RegisterValue}_R(x,h_x,s).
  \]
  The record itself is indexed only by $(P,Q,R)$.  The phase is an argument to
  each of the six fields, not an enclosing answer-history-dependent parameter.
  The structure has no field named $\operatorname{AnswerHistory}$,
  $\operatorname{History}$, $\operatorname{Trace}$, or any synonym, and it has
  no constructor argument whose type is an oracle-answer history.  In
  particular, the declaration must not include a parameter
  $H:\Hist_R(x,s)$, a default parameter
  $H:=\Hist_R(x,s)$, a let-bound reference to $\Hist_R(x,s)$, or any field that
  returns or stores $\Hist_R(x,s)$.

  These types are parameters of the local presentation at $(R,x,h_x,s)$ and
  are chosen before any oracle answers are supplied.  None of their
  declarations, indices, constructors, fields, subtype predicates, quotient
  relations, default arguments, local definitions, or field bodies may mention
  an answer-history type, a trace prefix, a prefix-history selector, a query
  produced from an answer prefix, or an oracle answer type $O_Q(q)$.

  The component meanings are fixed as follows.
  $\operatorname{Control}_R(x,h_x,s)$ is the control-state type.
  $\operatorname{CounterName}_R(x,h_x,s)\to\N$ is the counter store.
  $\operatorname{TapeName}_R(x,h_x,s)\to(\mathbb{Z}\to
  \operatorname{TapeSymbol}_R(x,h_x,s))$ is the work-tape content store, and
  $\operatorname{TapeName}_R(x,h_x,s)\to\mathbb{Z}$ is the
  work-head-position store.  $\operatorname{RegisterName}_R(x,h_x,s)\to
  \operatorname{RegisterValue}_R(x,h_x,s)$ is the store for all remaining
  deterministic local registers.  These store types are derived in the
  snapshot-record node from the six component fields above; they are not
  additional fields of
  $\operatorname{def\_oracle\_reduction\_local\_component\_types}$.  These
  six component fields are the only data types from which local snapshots are
  built.
\end{definition}

\begin{definition}[Oracle reduction local snapshot record]
  \label{def:oracle-reduction-local-snapshot-record}
  \uses{def:oracle-reduction-local-component-types}
  For a local-phase presentation of $R$, a valid input $x$, and a phase
  $s\in\{0,\ldots,r(x)\}$, the concrete local snapshot record is the product
  with exactly the following named projections:
  \[
  \begin{aligned}
    \operatorname{control}_R(x,s,u)&\in\operatorname{Control}_R(x,s),\\
    \operatorname{counters}_R(x,s,u)&\in
      \operatorname{CounterName}_R(x,s)\to\N,\\
    \operatorname{workTape}_R(x,s,u)&\in
      \operatorname{TapeName}_R(x,s)\to
      (\mathbb{Z}\to\operatorname{TapeSymbol}_R(x,s)),\\
    \operatorname{workHead}_R(x,s,u)&\in
      \operatorname{TapeName}_R(x,s)\to\mathbb{Z},\\
    \operatorname{registers}_R(x,s,u)&\in
      \operatorname{RegisterName}_R(x,s)\to
      \operatorname{RegisterValue}_R(x,s).
  \end{aligned}
  \]
  Equivalently, the record carrier is
  \[
  \begin{aligned}
    &\operatorname{Control}_R(x,s)
    \times(\operatorname{CounterName}_R(x,s)\to\N)
    \times(\operatorname{TapeName}_R(x,s)\to
      (\mathbb{Z}\to\operatorname{TapeSymbol}_R(x,s)))\\
    &\qquad\times(\operatorname{TapeName}_R(x,s)\to\mathbb{Z})
    \times(\operatorname{RegisterName}_R(x,s)\to
      \operatorname{RegisterValue}_R(x,s)),
  \end{aligned}
  \]
  with the displayed projection names.  This declaration is not an erased
  placeholder: any implementation of a local snapshot must provide exactly
  these non-oracle fields, even when some component type is empty or singleton
  for a particular reduction and phase.
\end{definition}

\begin{definition}[Oracle reduction non-oracle local snapshots]
  \label{def:oracle-reduction-local-snapshot}
  \uses{def:oracle-reduction-answer-history, def:oracle-reduction-local-snapshot-record}
  For an oracle reduction datum $R$, a valid $P$-input $x$, and a phase
  $s\in\{0,\ldots,r(x)\}$, the type
  \[
    \Local_R(x,s)
  \]
  is definitionally the local snapshot record of
  Definition~\ref{def:oracle-reduction-local-snapshot-record}.  Thus an element
  $u\in\Local_R(x,s)$ consists exactly of the non-oracle local data held by the
  reduction after $s$ oracle calls have been completed: its control state,
  counter store, work-tape contents, work-head positions, and deterministic
  register store.

  Lean-facing, $\Local_R(x,s)$ is indexed only by $(R,x,s)$ and is chosen
  before any particular answer history is supplied.  Its declaration is an
  answer-history-free local state type: no constructor argument, field,
  projection, invariant, subtype predicate, or field type may mention
  $\Hist_R(x,s)$, a prefix-history selector, a trace prefix, or an oracle
  answer type $O_Q(q)$ in any position.  In particular, a local snapshot may
  not contain a field such as
  $\Hist_R(x,s)\to T$, $T\to\Hist_R(x,s)$,
  $\forall H:\Hist_R(x,s),\,T(H)$,
  $\exists H:\Hist_R(x,s),\,T(H)$, a dependent pair over histories, a quotient
  by a relation using histories, or any analogous function/product/subtype
  whose domain, codomain, index, or proposition can encode oracle answers.
  A field named ``control state'', ``work tape'', or similar is permitted only
  when its type is built from non-oracle data already available from
  $(R,x,s)$ and does not contain histories or oracle answers under any type
  former.  In particular, the declaration may not replace the record by
  $\operatorname{Unit}$ unless the five displayed component fields themselves
  are still present with their named projections.

  Thus $\Local_R(x,s)$ stores only history-independent local data.  Any
  semantic fact saying that a local snapshot is the one reached after a
  particular oracle-answer prefix must therefore be stated as a relation with
  two separate arguments,
  \[
    u\in\Local_R(x,s),\qquad H\in\Hist_R(x,s),
  \]
  not by storing $H$ or the answers inside $u$.
\end{definition}

\begin{definition}[Oracle reduction local compatibility relation]
  \label{def:oracle-reduction-local-compatible}
  \uses{def:oracle-reduction-answer-history, def:oracle-reduction-local-snapshot}
  For a local-phase presentation of $R$, the external local compatibility
  relation is the proposition
  \[
    \operatorname{LocalCompatible}_R(x,s,u,H)
  \]
  on separate arguments
  \[
    u\in\Local_R(x,s),\qquad H\in\Hist_R(x,s).
  \]
  It holds exactly when $H$ is a well-formed answer history of length $s$ and
  the five projections of $u$ are the control state, counter store, work-tape
  contents, work-head positions, and deterministic register store obtained by
  running the non-oracle part of $R$ on input $x$ through the first $s$ oracle
  calls using the answer prefix $H$.

  This is a relation, not a field of $u$ and not a subtype condition in
  $\Local_R(x,s)$.  The history argument remains external, so compatibility
  ties a history-independent snapshot to a particular answer prefix without
  encoding that prefix or any oracle answer inside the snapshot.
\end{definition}

\begin{definition}[Oracle reduction raw state representation]
  \label{def:oracle-reduction-raw-state-representation}
  \uses{def:oracle-reduction-data, def:oracle-reduction-local-snapshot}
  For a local-phase presentation of $R$, the raw state representation on a
  valid $P$-input $x$ is the phase-indexed dependent sum
  \[
    \operatorname{RawState}_R(x)
    =
    \{\, (s,u)\mid s\in\{0,\ldots,r(x)\},\ u\in\Local_R(x,s)\,\}.
  \]
  It has the named projections
  \[
    \operatorname{rawPhase}_R(x,(s,u))=s,\qquad
    \operatorname{rawLocal}_R(x,(s,u))=u.
  \]
  The projection equations are part of the representation contract.  A state
  implementation therefore exposes the phase coordinate and local snapshot by
  name, rather than only through anonymous dependent-pair eliminators.
\end{definition}

\begin{definition}[Oracle reduction state compatibility relation]
  \label{def:oracle-reduction-state-compatible}
  \uses{def:oracle-reduction-answer-history, def:oracle-reduction-local-compatible, def:oracle-reduction-raw-state-representation}
  For a raw state $\sigma\in\operatorname{RawState}_R(x)$ and a separately
  supplied answer history $H$, the state-history compatibility relation is
  \[
    \operatorname{StateCompatible}_R(x,\sigma,H)
    \quad\Longleftrightarrow\quad
    \operatorname{LocalCompatible}_R
      \bigl(x,\operatorname{rawPhase}_R(x,\sigma),
        \operatorname{rawLocal}_R(x,\sigma),H\bigr).
  \]
  Hence compatibility simultaneously asserts that $H$ is a well-formed history
  prefix of the state's named phase and that the state's named local snapshot
  is the deterministic non-oracle snapshot reached after that prefix.  The
  relation is external to the state carrier and has the three displayed
  arguments $(x,\sigma,H)$.
\end{definition}

\begin{definition}[Oracle reduction state space]
  \label{def:oracle-reduction-state-space}
  \uses{def:oracle-reduction-data, def:oracle-reduction-answer-history, def:oracle-reduction-local-snapshot, def:oracle-reduction-raw-state-representation, def:oracle-reduction-state-compatible}
  For a local-phase presentation of an oracle reduction datum $R$, the state
  space on a valid $P$-input $x$ is the concrete raw state representation
  of Definition~\ref{def:oracle-reduction-raw-state-representation}:
  \[
    \State_R(x)=\operatorname{RawState}_R(x).
  \]
  Its projections are
  \[
    \operatorname{phase}_R(x,\sigma)
      =\operatorname{rawPhase}_R(x,\sigma),
    \qquad
    \operatorname{local}_R(x,\sigma)
      =\operatorname{rawLocal}_R(x,\sigma),
  \]
  with the displayed names.  The phase coordinate is bounded by the oracle-call
  count $r(x)$ from the reduction datum $R$, and the local projection has type
  $\Local_R(x,\operatorname{phase}_R(x,\sigma))$, so every state is explicitly
  tied to a definite point of an $R$-run and to the corresponding non-oracle
  snapshot.

  The state type contains only the phase and the non-oracle local snapshot.
  It does not contain an answer history and does not contain oracle answers.
  Compatibility with a particular prefix of oracle answers is instead the
  external relation
  \[
    \operatorname{StateCompatible}_R(x,\sigma,H)
  \]
  from Definition~\ref{def:oracle-reduction-state-compatible}, equivalently
  $\operatorname{LocalCompatible}_R(x,\operatorname{phase}_R(x,\sigma),
  \operatorname{local}_R(x,\sigma),H)$.  Thus the state-space contract consists
  of the carrier $\State_R(x)$, the named phase projection, the named local
  projection, and the named external compatibility relation.

  Consequently, later query, update, and postprocessing functions receive
  $\sigma\in\State_R(x)$ and the relevant history $H$ as separate arguments:
  $\sigma$ supplies the non-oracle local data, while $H$ supplies exactly the
  oracle query-answer prefix determined by $R$.
\end{definition}

\begin{definition}[Oracle reduction initial-state construction semantics]
  \label{def:oracle-reduction-initial-construction-semantics}
  \uses{def:oracle-reduction-state-space}
  For a local-phase presentation of $R$, the initialization computation on a
  valid $P$-input $x$ is specified by concrete non-oracle operational data,
  separate from all oracle histories.  It consists of a configuration type
  \[
    \operatorname{InitConfig}_R(x),
  \]
  an input-reading initial configuration
  \[
    \operatorname{initConfig}_R(x)\in\operatorname{InitConfig}_R(x),
  \]
  a deterministic local step map
  \[
    \operatorname{initStep}_R(x):
      \operatorname{InitConfig}_R(x)\to\operatorname{InitConfig}_R(x),
  \]
  a halting predicate $\operatorname{InitHalt}_R(x,c)$ on initialization
  configurations, and an output relation
  \[
    \operatorname{InitOut}_R(x,c,\sigma)
  \]
  between a halted initialization configuration $c$ and a state
  $\sigma\in\State_R(x)$.

  The initialization run is defined recursively by
  \[
    \operatorname{InitRun}_R(x,0)=\operatorname{initConfig}_R(x),
    \qquad
    \operatorname{InitRun}_R(x,s+1)=
      \operatorname{initStep}_R(x)(\operatorname{InitRun}_R(x,s)).
  \]
  The predicate
  \[
    \operatorname{InitHaltsWith}_R(x,s,\sigma)
  \]
  means that this non-oracle run first halts at exactly step $s$ and outputs
  exactly the state $\sigma$:
  \[
    \operatorname{InitHalt}_R(x,\operatorname{InitRun}_R(x,s))
    \wedge
    \operatorname{InitOut}_R(x,\operatorname{InitRun}_R(x,s),\sigma)
  \]
  and, for every $r<s$,
  \[
    \neg\operatorname{InitHalt}_R(x,\operatorname{InitRun}_R(x,r)).
  \]
  The output relation is functional on halted initialization configurations,
  and the deterministic initialization run has a unique first halting time:
  if
  $\operatorname{InitHaltsWith}_R(x,s,\sigma)$ and
  $\operatorname{InitHaltsWith}_R(x,s',\sigma')$, then
  $s=s'$ and $\sigma=\sigma'$.

  This node supplies the semantics for the initial-state construction only.
  It contains no oracle query, oracle answer, answer history, trace prefix, or
  target-problem output.  Its step count is the actual first halting time of
  this displayed initialization run.
\end{definition}

\begin{definition}[Oracle reduction initial state]
  \label{def:oracle-reduction-initial-state}
  \uses{def:oracle-reduction-state-space, def:oracle-reduction-initial-construction-semantics}
  For a local-phase presentation of $R$, the initial-state constructor and its
  local step count are not independent data.  They are a function
  \[
    \init_R(x)\in\State_R(x)
  \]
  and a natural number $\operatorname{InitSteps}_R(x)\in\N$ defined for every
  valid $P$-input $x$, together with the semantic certificate
  \[
    \operatorname{InitHaltsWith}_R
      \bigl(x,\operatorname{InitSteps}_R(x),\init_R(x)\bigr).
  \]
  Equivalently,
  $\operatorname{InitSteps}_R(x)$ is the exact first halting time of the
  concrete initialization run from
  Definition~\ref{def:oracle-reduction-initial-construction-semantics}, and
  $\init_R(x)$ is exactly the state output by that halted run.  The declaration
  therefore has the following invariant:
  \[
    \forall s,\sigma,\quad
    \operatorname{InitHaltsWith}_R(x,s,\sigma)
    \Rightarrow
    s=\operatorname{InitSteps}_R(x)\wedge \sigma=\init_R(x).
  \]
  Thus an implementation cannot certify this node by choosing an arbitrary
  natural-number field; the step count must be witnessed by the displayed
  non-oracle construction of exactly the displayed initial state.
\end{definition}

\begin{definition}[Oracle reduction trace prefix history]
  \label{def:oracle-reduction-trace-prefix-history}
  \uses{def:oracle-reduction-trace}
  Define
  $\operatorname{def\_oracle\_reduction\_trace\_prefix\_history}(P,Q,R)$ to be
  the type of prefix-history selector functions $H$.  Such an $H$ takes a
  valid $P$-input $x$, a trace
  $(q_1,z_1),\ldots,(q_r,z_r),y$ of $R$ on $x$, and an index
  $s\in\{1,\ldots,r+1\}$, and returns the finite list
  \[
    H(x,t,s)=((q_1,z_1),\ldots,(q_{s-1},z_{s-1})).
  \]
  In particular $H(x,t,r+1)$ is the full query-answer history.

  This declaration is a type of selectors, not a recursively defined function
  with a hidden body.  The equality saying that the selected list is the first
  $s-1$ entries of the trace is part of the local trace-consistency predicate
  below, so later nodes use a concrete selector argument $H$ rather than an
  undefined global prefix-history value.
\end{definition}

\begin{definition}[Oracle reduction state-dependent query construction]
  \label{def:oracle-reduction-query-construction}
  \uses{def:oracle-reduction-state-space, def:oracle-reduction-trace-prefix-history}
  At oracle-call position $s$, the query constructor of a local-phase
  presentation is a state-dependent function
  \[
    \query_R(x,s,\sigma,H_s)\in I_Q
  \]
  where $\sigma\in\State_R(x)$ is the previous local state and $H_s$ is the
  prefix history before call $s$.  Its local step count is the function
  \[
    \operatorname{QuerySteps}_R(x,s,\sigma,H_s)\in\N .
  \]
  Thus the query and its construction cost are both allowed to depend on the
  previous state, not only on $x$, the call index, and the earlier
  query-answer pairs.  In a locally generated trace, the concrete query at
  position $s$ satisfies
  \[
    q_s=\query_R(x,s,\sigma_{s-1},H_s).
  \]
\end{definition}

\begin{definition}[Oracle reduction state update]
  \label{def:oracle-reduction-state-update}
  \uses{def:oracle-reduction-query-construction}
  After oracle-call position $s$, the state-transition function is
  \[
    \update_R(x,s,\sigma_{s-1},H_s,q_s,z_s)\in\State_R(x).
  \]
  Its local step count
  \[
    \operatorname{UpdateSteps}_R(x,s,\sigma_{s-1},H_s,q_s,z_s)\in\N
  \]
  counts exactly the non-oracle work needed to update the local state after the
  answer $z_s$ to query $q_s$ is available.  The next state in the run is
  \[
    \sigma_s=\update_R(x,s,\sigma_{s-1},H_s,q_s,z_s).
  \]
\end{definition}

\begin{definition}[Oracle reduction generated state sequence]
  \label{def:oracle-reduction-generated-state-sequence}
  \uses{def:oracle-reduction-initial-state, def:oracle-reduction-state-update, def:oracle-reduction-trace-prefix-history}
  For a concrete trace
  $t=(q_1,z_1),\ldots,(q_r,z_r),y$ of $R$ on a valid input $x$, define
  \[
    \operatorname{def\_oracle\_reduction\_generated\_state\_sequence}
  \]
  to be the predicate on candidate finite state sequences
  \[
    G:\{0,\ldots,r\}\to\State_R(x).
  \]
  The predicate holds exactly when $G$ satisfies the two recursive equations
  \[
    G(0)=\init_R(x),
  \]
  and for every $1\leq s\leq r$,
  \[
    G(s)=\update_R(x,s,G(s-1),H(x,t,s),q_s,z_s).
  \]
  This node is a proposition relating an explicit candidate sequence $G$ to
  the update equations.  It does not introduce a separate globally named
  recursive function returning $G$, and later nodes quantify over such a
  concrete $G$.
\end{definition}

\begin{definition}[Oracle reduction final postprocessing]
  \label{def:oracle-reduction-final-postprocessing}
  \uses{def:oracle-reduction-generated-state-sequence}
  The final postprocessing function of a local-phase presentation is
  \[
    \finish_R(x,\sigma,H)\in O_P(x),
  \]
  where $\sigma$ is an explicitly supplied final local state and $H$ is an
  explicitly supplied full query-answer history.  Its local step count
  \[
    \operatorname{FinishSteps}_R(x,\sigma,H)\in\N
  \]
  counts exactly the non-oracle work needed to produce the final output from
  those arguments.  In a locally generated trace with generated state sequence
  $G$, the final state argument is $G(r)$, the history argument is
  $H(x,t,r+1)$, and the final output satisfies
  \[
    y=\finish_R(x,G(r),H(x,t,r+1)).
  \]
\end{definition}

\begin{definition}[Oracle reduction local trace consistency]
  \label{def:oracle-reduction-local-trace-consistency}
  \uses{def:oracle-reduction-query-construction, def:oracle-reduction-generated-state-sequence, def:oracle-reduction-final-postprocessing}
  A trace
  \[
    (q_1,z_1),\ldots,(q_r,z_r),y
  \]
  is locally generated by the local-phase presentation of $R$ on $x$ when
  there exists an explicit state sequence
  $G:\{0,\ldots,r\}\to\State_R(x)$ such that
  Definition~\ref{def:oracle-reduction-generated-state-sequence} holds for
  $G$, the selector $H(x,t,s)$ is the first $s-1$ query-answer entries of the
  trace for every $s$, and both displayed equations hold:
  \[
    \forall\,1\leq s\leq r,\qquad
    q_s=\query_R(x,s,G(s-1),H(x,t,s)),
  \]
  and
  \[
    y=\finish_R(x,G(r),H(x,t,r+1)).
  \]
  This predicate is the formal link between an oracle reduction trace and the
  state-dependent local operational data.  In particular, the query at call
  $s$ is certified from the explicit generated previous state $G(s-1)$.
\end{definition}

\begin{definition}[Oracle reduction local phases]
  \label{def:oracle-reduction-local-phases}
  \uses{def:oracle-reduction-local-trace-consistency}
  For a trace
  $t=(q_1,z_1),\ldots,(q_r,z_r),y$ of $R$ on $x$, define
  $\operatorname{def\_oracle\_reduction\_local\_phases}$ to be the relation
  between a candidate generated state sequence
  $G:\{0,\ldots,r\}\to\State_R(x)$ and a natural number $L$.  It holds exactly
  when the trace is locally generated using $G$ in the sense of
  Definition~\ref{def:oracle-reduction-local-trace-consistency} and $L$ is the
  local running-time sum
  \[
    L
    =
    \operatorname{InitSteps}_R(x)
    +\sum_{s=1}^{r}
      \Bigl(
        \operatorname{QuerySteps}_R(x,s,G(s-1),H(x,t,s))
        +
        \operatorname{UpdateSteps}_R(x,s,G(s-1),H(x,t,s),q_s,z_s)
      \Bigr)
    +\operatorname{FinishSteps}_R(x,G(r),H(x,t,r+1)).
  \]
  The query term uses the previous generated state $G(s-1)$ explicitly, and
  the update term is tied to the state equation defining $G(s)$ in
  Definition~\ref{def:oracle-reduction-generated-state-sequence}.  The local
  running-time value $L$
  includes exactly the initial-state construction, every state-dependent query
  construction, every state update after an oracle answer, and final
  postprocessing.  It does not include the time spent solving the $Q$-queries.
\end{definition}

\begin{definition}[Oracle reduction cost accounting]
  \label{def:oracle-reduction-cost}
  \uses{def:oracle-reduction-local-phases, def:exact-worst-case-running-time}
  Fix a correct $Q$-algorithm with running-time bound $T_Q$.  For a trace
  $t=(q_1,z_1),\ldots,(q_r,z_r),y$ of $R$ on input $x$, a candidate generated
  state sequence $G$, and a natural number $L$ satisfying
  $\operatorname{def\_oracle\_reduction\_local\_phases}(x,t,G,L)$, the induced
  running time of the simulated reduction run is bounded by
  \[
    L
    +\sum_{s=1}^{r} T_Q(N_Q(q_s),L_Q(q_s)).
  \]
  If no such $G$ and $L$ exist, the trace is not a locally generated trace of
  the reduction and contributes no valid simulated execution.
  The oracle-cost term uses the actual principal size and auxiliary dimension
  of each constructed query, not the parameters of the original $P$-input.
\end{definition}

\begin{definition}[Subquadratic oracle-query budget]
  \label{def:subquadratic-oracle-query-budget}
  \uses{def:oracle-reduction-cost, def:truly-subquadratic}
  An oracle reduction datum $R$ has a subquadratic transfer budget from $P$ to
  $Q$ if the following implication holds for every valid $P$-input $x$ with
  $N_P(x)=n$ and $L_P(x)=\ell$.  Whenever the assumed $Q$-algorithm has running
  time
  \[
    T_Q(n',\ell')\leq C(n')^{2-\eps}(1+\ell')^b,
  \]
  every trace of $R$ on $x$ satisfies a bound
  \[
    L+\sum_s T_Q(N_Q(q_s),L_Q(q_s))
      \leq C' n^{2-\eps'}(1+\ell)^{b'}
  \]
  for every locally generated trace and every local-cost witness
  $\operatorname{def\_oracle\_reduction\_local\_phases}(x,t,G,L)$,
  for some constants $\eps'>0$, $C'>0$, and $b'\in\N$ that depend only on
  $\eps,C,b$ and the reduction, not on the particular input $x$ or oracle
  answers.  Thus the number of calls, query sizes, and all local construction,
  update, and postprocessing costs are explicitly included in the transfer.
\end{definition}

\begin{definition}[Quadratic-time baseline algorithm]
  \label{def:quadratic-baseline-algorithm}
  \uses{def:computational-problem-spec, def:problem-algorithm, def:exact-worst-case-running-time}
  A computational problem $P$ has a quadratic-time baseline algorithm if there
  is a concrete deterministic algorithm $B_P$ for $P$, a worst-case running-time
  bound $T_{B_P}:\N\times\N\to\R_{\geq0}$ for the actual operational execution
  of $B_P$, and constants $C>0$, $b\in\N$, and $n_0\in\N$ such that $B_P$
  halts correctly on every valid input and
  \[
    T_{B_P}(n,\ell)\leq C n^2(1+\ell)^b
  \]
  for every $n\geq n_0$ and every auxiliary dimension $\ell$ in the domain of
  $P$.

  Lean-facing, the proposition
  \[
    \operatorname{def\_quadratic\_baseline\_algorithm}(P)
  \]
  contains the displayed witnesses $(B_P,T_{B_P},C,b,n_0)$ and the two
  displayed obligations: semantic correctness of $B_P$ in the sense of
  Definition~\ref{def:problem-algorithm}, and the quadratic polynomial upper
  bound above.  It is not a name tag on $P$ and it is not automatically true
  for an arbitrary computational-problem specification.
\end{definition}

\begin{definition}[Composed oracle-reduction query phase machine]
  \label{def:composed-oracle-query-phase-machine}
  \uses{def:oracle-reduction-query-construction}
  For fixed data $(R,x,s,\sigma,H_s)$ with
  $\sigma\in\State_R(x)$ and $H_s$ the prefix history before oracle call $s$,
  the query phase in the composed program is the canonical countdown machine
  \[
    \operatorname{QueryConfig}_{R}(x,s,\sigma,H_s)
      =\{0,\ldots,\operatorname{QuerySteps}_R(x,s,\sigma,H_s)\}.
  \]
  Its initial configuration is
  $\operatorname{QuerySteps}_R(x,s,\sigma,H_s)$, its step map sends a positive
  counter $k+1$ to $k$ and fixes $0$, its halt predicate is exactly
  $k=0$, and its output at the halted configuration is exactly
  \[
    \query_R(x,s,\sigma,H_s)\in I_Q .
  \]
  Thus the query phase has no externally supplied query output: the only query
  it can emit is the state-dependent query constructor from
  Definition~\ref{def:oracle-reduction-query-construction}, after exactly the
  displayed number of local countdown steps.
\end{definition}

\begin{definition}[Composed oracle-reduction update phase machine]
  \label{def:composed-oracle-update-phase-machine}
  \uses{def:oracle-reduction-state-update}
  For fixed data $(R,x,s,\sigma,H_s,q,z)$, the update phase in the composed
  program is the canonical countdown machine
  \[
    \operatorname{UpdateConfig}_{R}(x,s,\sigma,H_s,q,z)
      =\{0,\ldots,\operatorname{UpdateSteps}_R(x,s,\sigma,H_s,q,z)\}.
  \]
  Its initial configuration is
  $\operatorname{UpdateSteps}_R(x,s,\sigma,H_s,q,z)$, its step map decrements a
  positive counter and fixes $0$, its halt predicate is exactly $k=0$, and its
  output at the halted configuration is exactly the next local state
  \[
    \update_R(x,s,\sigma,H_s,q,z)\in\State_R(x).
  \]
  Hence the update phase cannot choose an independent next state; its only
  halted output is the state-transition value from
  Definition~\ref{def:oracle-reduction-state-update}.
\end{definition}

\begin{definition}[Composed oracle-reduction finish phase machine]
  \label{def:composed-oracle-finish-phase-machine}
  \uses{def:oracle-reduction-final-postprocessing}
  For fixed data $(R,x,\sigma,H)$ with $\sigma\in\State_R(x)$ and $H$ a full
  answer history, the finish phase in the composed program is the canonical
  countdown machine
  \[
    \operatorname{FinishConfig}_{R}(x,\sigma,H)
      =\{0,\ldots,\operatorname{FinishSteps}_R(x,\sigma,H)\}.
  \]
  Its initial configuration is
  $\operatorname{FinishSteps}_R(x,\sigma,H)$, its step map decrements a positive
  counter and fixes $0$, its halt predicate is exactly $k=0$, and its output at
  the halted configuration is exactly
  \[
    \finish_R(x,\sigma,H)\in O_P(x).
  \]
  Thus the final $P$-output of the composed program is produced by the
  displayed final-postprocessing function, not by a separate output parameter.
\end{definition}

\begin{definition}[Composed oracle-reduction phase records]
  \label{def:composed-oracle-phase-records}
  \uses{def:problem-algorithm, def:oracle-reduction-initial-state, def:composed-oracle-query-phase-machine, def:composed-oracle-update-phase-machine, def:composed-oracle-finish-phase-machine}
  Fix a local-phase oracle reduction datum $R$ from $P$ to $Q$ and a concrete
  deterministic $Q$-algorithm $A_Q$.  The configuration carrier of the composed
  program $R[A_Q]$ on a valid input $x$ is the disjoint union of the following
  concrete phase records, with no additional output-producing case:
  \[
    \operatorname{InitPhase}(c_{\rm init}),
    \quad
    \operatorname{QueryPhase}(s,\sigma_{s-1},H_s,c_{\rm qry}),
  \]
  \[
    \operatorname{OraclePhase}(s,\sigma_{s-1},H_s,q_s,c_Q),
    \quad
    \operatorname{UpdatePhase}(s,\sigma_{s-1},H_s,q_s,z_s,c_{\rm upd}),
  \]
  and
  \[
    \operatorname{FinishPhase}(\sigma_r,H_{r+1},c_{\rm fin}).
  \]
  Here $c_{\rm init}$ is an initialization configuration of
  Definition~\ref{def:oracle-reduction-initial-state}.  The configurations
  $c_{\rm qry}$, $c_{\rm upd}$, and $c_{\rm fin}$ are respectively elements of
  the canonical countdown machines of
  Definitions~\ref{def:composed-oracle-query-phase-machine},
  \ref{def:composed-oracle-update-phase-machine}, and
  \ref{def:composed-oracle-finish-phase-machine}.  The configuration $c_Q$ is
  an actual operational configuration of $A_Q$ on the displayed concrete query.
  The phase record stores the previous state and prefix history as separate
  arguments whenever they are needed; it does not store arbitrary oracle
  answers or an arbitrary final $P$-output.
\end{definition}

\begin{definition}[Composed oracle-reduction boundary equations]
  \label{def:composed-oracle-boundary-equations}
  \uses{def:oracle-reduction-initial-state, def:oracle-reduction-final-postprocessing, def:oracle-reduction-trace-prefix-history}
  For a valid $P$-input $x$, a trace
  $t=(q_1,z_1),\ldots,(q_{r(x)},z_{r(x)}),y$, and a candidate generated state
  sequence $G:\{0,\ldots,r(x)\}\to\State_R(x)$, the boundary-equation
  predicate is the conjunction of exactly the following two equalities:
  \[
    G(0)=\init_R(x)
  \]
  and
  \[
    y=\finish_R(x,G(r(x)),H(x,t,r(x)+1)).
  \]
  Lean-facing, the declaration
  $\operatorname{def\_composed\_oracle\_boundary\_equations}(x,t,G)$ is a
  proposition with these equalities as eliminable conjuncts.  It cannot be
  proved by constructing an initialization phase record or a finish phase
  record unless the displayed equality tying that record to the supplied
  sequence $G$ and output $y$ is also proved.
\end{definition}

\begin{definition}[Composed oracle-reduction per-call equations]
  \label{def:composed-oracle-call-equations}
  \uses{def:oracle-reduction-query-construction, def:oracle-reduction-state-update, lem:solver-halts-on-promised-input-set}
  For a valid $P$-input $x$, a trace
  $t=(q_1,z_1),\ldots,(q_{r(x)},z_{r(x)}),y$, a candidate generated state
  sequence $G:\{0,\ldots,r(x)\}\to\State_R(x)$, and oracle first-halting-time
  data $U:\{1,\ldots,r(x)\}\to\N$, the per-call-equation predicate holds when
  the following three facts hold for every oracle-call position
  $1\leq s\leq r(x)$:
  \[
    q_s=\query_R(x,s,G(s-1),H(x,t,s)),
  \]
  \[
    \operatorname{HaltsWith}_{A_Q}(q_s,U(s),z_s),
  \]
  and
  \[
    G(s)=\update_R(x,s,G(s-1),H(x,t,s),q_s,z_s).
  \]
  The same preceding state $G(s-1)$, prefix history $H(x,t,s)$, query $q_s$,
  and oracle answer $z_s$ must appear in all three conjuncts.  In particular,
  the update equation is not a local statement about an independently built
  update-phase configuration: it is an equality whose left-hand side is the
  supplied next state $G(s)$ from the candidate state sequence.
\end{definition}

\begin{definition}[Composed oracle-reduction ordered phase trace]
  \label{def:composed-oracle-ordered-phase-trace}
  \uses{def:composed-oracle-phase-records, def:composed-oracle-boundary-equations, def:composed-oracle-call-equations}
  For the same data $(x,t,G,U)$, an ordered phase trace is a finite sequence
  $\Pi$ of concrete phase records from
  Definition~\ref{def:composed-oracle-phase-records}, indexed by the ordered
  positions
  \[
    0,1,\ldots,3r(x)+1.
  \]
  It satisfies the following positional equations.  Position $0$ is an
  initialization phase whose halted output state is exactly $G(0)$, and this
  same state satisfies the equality $G(0)=\init_R(x)$.  For every
  $1\leq s\leq r(x)$, the three consecutive positions
  $3s-2,3s-1,3s$ are respectively
  \[
    \operatorname{QueryPhase}(s,G(s-1),H(x,t,s),c_{\rm qry,s}),
  \]
  \[
    \operatorname{OraclePhase}(s,G(s-1),H(x,t,s),q_s,c_{Q,s}),
  \]
  and
  \[
    \operatorname{UpdatePhase}(s,G(s-1),H(x,t,s),q_s,z_s,c_{\rm upd,s}),
  \]
  where the query value, oracle answer, and next state obey the per-call
  equations of Definition~\ref{def:composed-oracle-call-equations}.  The final
  position $3r(x)+1$ is
  \[
    \operatorname{FinishPhase}(G(r(x)),H(x,t,r(x)+1),c_{\rm fin}),
  \]
  and its output is exactly the trace output
  $y=\finish_R(x,G(r(x)),H(x,t,r(x)+1))$.

  Thus
  $\operatorname{def\_composed\_oracle\_ordered\_phase\_trace}(x,t,G,U)$ is a
  proposition asserting the existence of this single ordered sequence $\Pi$
  together with the displayed positional equalities.  It is not satisfied by a
  collection of phase records constructed independently of $\Pi$, and it is not
  satisfied by records whose state arguments differ from the entries of $G$.
  The ordered trace is the bridge between the syntactic phase records and the
  semantic equations in the trace certificate.
\end{definition}

\begin{definition}[Composed oracle-reduction run equations]
  \label{def:composed-oracle-run-equations}
  \uses{def:composed-oracle-boundary-equations, def:composed-oracle-call-equations, def:composed-oracle-ordered-phase-trace}
  The Lean-facing predicate
  \[
    \operatorname{def\_composed\_oracle\_run\_equations}(x,t,G,U)
  \]
  is the concrete trace certificate generated by the step relation of the
  composed program $R[A_Q]$.  Its arguments are a valid $P$-input $x$, a trace
  \[
    t=(q_1,z_1),\ldots,(q_{r(x)},z_{r(x)}),y,
  \]
  a candidate generated state sequence
  $G:\{0,\ldots,r(x)\}\to\State_R(x)$, and oracle first-halting-time data
  $U:\{1,\ldots,r(x)\}\to\N$.  The predicate is definitionally the conjunction
  of the following three subpredicates:
  \[
    \operatorname{def\_composed\_oracle\_boundary\_equations}(x,t,G),
  \]
  \[
    \operatorname{def\_composed\_oracle\_call\_equations}(x,t,G,U),
  \]
  and
  \[
    \operatorname{def\_composed\_oracle\_ordered\_phase\_trace}(x,t,G,U).
  \]
  Consequently, eliminating a proof of
  $\operatorname{def\_composed\_oracle\_run\_equations}(x,t,G,U)$ yields all of
  the following concrete obligations for the same trace $t$, state sequence
  $G$, and oracle time data $U$:
  \[
    G(0)=\init_R(x),
  \]
  \[
    \forall\,1\leq s\leq r(x),\quad
      q_s=\query_R(x,s,G(s-1),H(x,t,s)),
  \]
  \[
    \forall\,1\leq s\leq r(x),\quad
      \operatorname{HaltsWith}_{A_Q}(q_s,U(s),z_s),
  \]
  \[
    \forall\,1\leq s\leq r(x),\quad
      G(s)=\update_R(x,s,G(s-1),H(x,t,s),q_s,z_s),
  \]
  \[
    y=\finish_R(x,G(r(x)),H(x,t,r(x)+1)),
  \]
  and a single ordered phase-record sequence
  $\Pi=(\operatorname{InitPhase},\operatorname{QueryPhase},
  \operatorname{OraclePhase},\operatorname{UpdatePhase},\ldots,
  \operatorname{FinishPhase})$ whose records carry exactly these same states,
  histories, queries, oracle answers, and final output in the order specified
  by Definition~\ref{def:composed-oracle-ordered-phase-trace}.

  This definition is intentionally not a mere constructor for phase records.
  A Lean declaration for this node must expose the three displayed conjuncts
  (or an equivalent proposition from which each displayed equation and the
  ordered phase trace can be projected).  It must not replace the initialization
  equality, the update equalities, or the ordered run/trace relation by local
  let-bound phase records, by an existentially chosen unrelated state sequence,
  or by a predicate that only says the phase records are well typed.
\end{definition}

\begin{definition}[Composed oracle-reduction program carrier]
  \label{def:composed-oracle-program-carrier}
  \uses{def:composed-oracle-phase-records}
  For fixed $(P,Q,R,A_Q)$ and a valid $P$-input $x$, define
  \[
    \operatorname{ComposedConfig}_{R,A_Q}(x)
  \]
  to be exactly the disjoint union of phase records from
  Definition~\ref{def:composed-oracle-phase-records}.  The five constructors
  are
  \[
    \operatorname{InitPhase},\quad \operatorname{QueryPhase},\quad
    \operatorname{OraclePhase},\quad \operatorname{UpdatePhase},\quad
    \operatorname{FinishPhase},
  \]
  with the arguments and component configuration types specified in that
  definition.  There is no sixth constructor, no externally supplied halted
  output constructor, and no constructor that packages a complete trace
  certificate.
\end{definition}

\begin{definition}[Composed oracle-reduction initial configuration]
  \label{def:composed-oracle-initial-configuration}
  \uses{def:composed-oracle-program-carrier, def:oracle-reduction-initial-construction-semantics}
  The initial configuration of the composed program on a valid $P$-input $x$ is
  the concrete phase record
  \[
    \operatorname{ComposedInit}_{R,A_Q}(x)
      =
    \operatorname{InitPhase}(\operatorname{initConfig}_R(x))
    \in \operatorname{ComposedConfig}_{R,A_Q}(x).
  \]
  This definition has no argument for a starting state, trace, answer history,
  or precomputed output.  The initialization phase begins with the actual
  initial configuration of the initialization machine of $R$.
\end{definition}

\begin{definition}[Composed oracle-reduction history operations]
  \label{def:composed-oracle-history-operations}
  \uses{def:oracle-reduction-answer-history}
  For a valid $P$-input $x$, the composed program uses concrete prefix
  histories, independent of any completed trace selector.  The empty history
  is the unique element
  \[
    \operatorname{EmptyHist}_R(x)\in\Hist_R(x,0)
  \]
  whose finite query-answer list has length $0$.

  For an oracle-call index $1\leq s\leq r(x)$, a prefix history
  $H_s\in\Hist_R(x,s-1)$ before call $s$, a query $q\in I_Q$, and an answer
  $z\in O_Q(q)$, the append operation is the concrete finite-list append
  \[
    \operatorname{AppendHist}_R(x,s,H_s,q,z)\in\Hist_R(x,s).
  \]
  Its first $s-1$ query-answer entries are exactly the entries of $H_s$ in the
  same order, and its $s$th entry is exactly $(q,z)$.  Equivalently, for every
  position $u<s$,
  \[
    \operatorname{AppendHist}_R(x,s,H_s,q,z)[u]=H_s[u],
  \]
  and
  \[
    \operatorname{AppendHist}_R(x,s,H_s,q,z)[s]=(q,z).
  \]
  These equations are definitional branch data for the composed step map.  The
  implementation must not replace the empty history or appended history by an
  existentially chosen history satisfying a later trace predicate.
\end{definition}

\begin{definition}[Oracle halted-output selector]
  \label{def:oracle-halted-output-selector}
  \uses{def:deterministic-operational-semantics, def:problem-algorithm}
  Let $A_Q$ be the concrete deterministic $Q$-program used by the composed
  reduction, let $q\in I_Q$, and let
  $c_Q\in\operatorname{Config}_{A_Q}(q)$ be an oracle configuration satisfying
  $\operatorname{Halt}_{A_Q}(q,c_Q)$.  The halted-output selector is the
  uniquely determined value
  \[
    \operatorname{HaltedOut}_{A_Q}(q,c_Q)\in O_Q(q)
  \]
  such that
  \[
    \operatorname{Out}_{A_Q}
      \bigl(q,c_Q,\operatorname{HaltedOut}_{A_Q}(q,c_Q)\bigr).
  \]
  Moreover, if $\operatorname{Out}_{A_Q}(q,c_Q,z')$, then
  $z'=\operatorname{HaltedOut}_{A_Q}(q,c_Q)$.  This selector is the only
  oracle answer used by the halted-oracle branch of
  $\operatorname{ComposedStep}_{R,A_Q}$.  It is not an argument of that branch
  and it is not chosen by an existential proof of trace consistency.
\end{definition}

\begin{definition}[Composed oracle-reduction step branch targets]
  \label{def:composed-oracle-step-branch-targets}
  \uses{def:composed-oracle-program-carrier, def:oracle-reduction-initial-state, def:composed-oracle-query-phase-machine, def:composed-oracle-update-phase-machine, def:composed-oracle-finish-phase-machine, def:composed-oracle-history-operations, def:oracle-halted-output-selector}
  For fixed $(P,Q,R,A_Q)$, valid input $x$, and a current configuration
  $c\in\operatorname{ComposedConfig}_{R,A_Q}(x)$, the branch target of the
  composed step is the following total case split.

  If $c=\operatorname{InitPhase}(c_{\rm init})$ and
  $\operatorname{InitHalt}_R(x,c_{\rm init})$ is false, the target is
  \[
    \operatorname{InitPhase}
      \bigl(\operatorname{initStep}_R(x)(c_{\rm init})\bigr).
  \]
  If $c=\operatorname{InitPhase}(c_{\rm init})$ and
  $\operatorname{InitHalt}_R(x,c_{\rm init})$ is true, the target ignores any
  output value carried by the halted initialization configuration and uses the
  definitionally displayed state $\init_R(x)$.  When $r(x)=0$ it is
  \[
    \operatorname{FinishPhase}
      \bigl(\init_R(x),\operatorname{EmptyHist}_R(x),
        \operatorname{FinishSteps}_R
          (x,\init_R(x),\operatorname{EmptyHist}_R(x))\bigr),
  \]
  and when $r(x)>0$ it is
  \[
    \operatorname{QueryPhase}
      \bigl(1,\init_R(x),\operatorname{EmptyHist}_R(x),
        \operatorname{QuerySteps}_R
          (x,1,\init_R(x),\operatorname{EmptyHist}_R(x))\bigr).
  \]

  If $c=\operatorname{QueryPhase}(s,\sigma,H_s,c_{\rm qry})$ and the query
  countdown value $c_{\rm qry}$ is positive, the target is the same query phase
  with countdown predecessor $c_{\rm qry}-1$.  If $c_{\rm qry}=0$, define
  \[
    q=\query_R(x,s,\sigma,H_s).
  \]
  The target is exactly
  \[
    \operatorname{OraclePhase}
      \bigl(s,\sigma,H_s,q,\operatorname{init}_{A_Q}(q)\bigr).
  \]

  If $c=\operatorname{OraclePhase}(s,\sigma,H_s,q,c_Q)$ and
  $\operatorname{Halt}_{A_Q}(q,c_Q)$ is false, the target is
  \[
    \operatorname{OraclePhase}
      \bigl(s,\sigma,H_s,q,\operatorname{step}_{A_Q}(q)(c_Q)\bigr).
  \]
  If $\operatorname{Halt}_{A_Q}(q,c_Q)$ is true, define
  \[
    z=\operatorname{HaltedOut}_{A_Q}(q,c_Q).
  \]
  The target is exactly
  \[
    \operatorname{UpdatePhase}
      \bigl(s,\sigma,H_s,q,z,
        \operatorname{UpdateSteps}_R(x,s,\sigma,H_s,q,z)\bigr).
  \]

  If $c=\operatorname{UpdatePhase}(s,\sigma,H_s,q,z,c_{\rm upd})$ and the
  update countdown value $c_{\rm upd}$ is positive, the target is the same
  update phase with countdown predecessor $c_{\rm upd}-1$.  If
  $c_{\rm upd}=0$, define
  \[
    \sigma'=\update_R(x,s,\sigma,H_s,q,z),
    \qquad
    H_{s+1}=\operatorname{AppendHist}_R(x,s,H_s,q,z).
  \]
  When $s<r(x)$ the target is exactly
  \[
    \operatorname{QueryPhase}
      \bigl(s+1,\sigma',H_{s+1},
        \operatorname{QuerySteps}_R(x,s+1,\sigma',H_{s+1})\bigr),
  \]
  and when $s=r(x)$ the target is exactly
  \[
    \operatorname{FinishPhase}
      \bigl(\sigma',H_{s+1},
        \operatorname{FinishSteps}_R(x,\sigma',H_{s+1})\bigr).
  \]

  If $c=\operatorname{FinishPhase}(\sigma,H,c_{\rm fin})$ and
  $c_{\rm fin}$ is positive, the target is the same finish phase with
  countdown predecessor $c_{\rm fin}-1$.  If $c_{\rm fin}=0$, the target is
  exactly the unchanged configuration
  $\operatorname{FinishPhase}(\sigma,H,0)$.

  The phrase ``the target is exactly'' is part of the formal contract: each
  branch returns the displayed constructor with the displayed arguments.  No
  branch target may be replaced by an existentially chosen phase record, by a
  field supplied in a structure of step-map data, or by a configuration known
  only to satisfy a later trace-characterization predicate.
\end{definition}

\begin{definition}[Composed oracle-reduction step map]
  \label{def:composed-oracle-step-map}
  \uses{def:composed-oracle-program-carrier, def:composed-oracle-query-phase-machine, def:composed-oracle-update-phase-machine, def:composed-oracle-finish-phase-machine, def:problem-algorithm, def:composed-oracle-history-operations, def:oracle-halted-output-selector, def:composed-oracle-step-branch-targets}
  The deterministic step map
  \[
    \operatorname{ComposedStep}_{R,A_Q}(x):
      \operatorname{ComposedConfig}_{R,A_Q}(x)\to
      \operatorname{ComposedConfig}_{R,A_Q}(x)
  \]
  is the concrete total function whose value at every configuration is the
  branch target of
  Definition~\ref{def:composed-oracle-step-branch-targets}.  Equivalently, it
  is defined by case analysis on the five phase constructors as follows.

  On $\operatorname{InitPhase}(c)$, if $c$ has not halted for the initialization
  machine of $R$, the next configuration is
  $\operatorname{InitPhase}(\operatorname{initStep}_R(x)(c))$.  If $c$ has
  halted, the next state is the displayed initial state $\init_R(x)$.  If
  $r(x)=0$ the step enters the finish countdown phase
  \[
    \operatorname{FinishPhase}
      \bigl(\init_R(x),\operatorname{EmptyHist}_R(x),
        \operatorname{FinishSteps}_R
          (x,\init_R(x),\operatorname{EmptyHist}_R(x))\bigr),
  \]
  and otherwise it enters the first query countdown phase
  \[
    \operatorname{QueryPhase}
      \bigl(1,\init_R(x),\operatorname{EmptyHist}_R(x),
        \operatorname{QuerySteps}_R
          (x,1,\init_R(x),\operatorname{EmptyHist}_R(x))\bigr).
  \]

  On a query phase
  $\operatorname{QueryPhase}(s,\sigma,H_s,c_{\rm qry})$, if the query countdown
  has not halted, the step decrements $c_{\rm qry}$ inside the same phase.  If
  it has halted, the emitted query is definitionally
  \[
    q=\query_R(x,s,\sigma,H_s),
  \]
  and the next configuration is
  \[
    \operatorname{OraclePhase}(s,\sigma,H_s,q,\operatorname{init}_{A_Q}(q)).
  \]

  On an oracle phase
  $\operatorname{OraclePhase}(s,\sigma,H_s,q,c_Q)$, if $c_Q$ has not halted for
  the operational semantics of $A_Q$ on $q$, the next configuration is
  \[
    \operatorname{OraclePhase}
      (s,\sigma,H_s,q,\operatorname{step}_{A_Q}(q)(c_Q)).
  \]
  If $c_Q$ has halted, let $z$ be the unique output selected by the output
  relation of $A_Q$ at this halted configuration, namely
  $z=\operatorname{HaltedOut}_{A_Q}(q,c_Q)$.  The next configuration is
  \[
    \operatorname{UpdatePhase}
      \bigl(s,\sigma,H_s,q,z,
        \operatorname{UpdateSteps}_R(x,s,\sigma,H_s,q,z)\bigr).
  \]

  On an update phase
  $\operatorname{UpdatePhase}(s,\sigma,H_s,q,z,c_{\rm upd})$, if the update
  countdown has not halted, the step decrements $c_{\rm upd}$ inside the same
  phase.  If it has halted, define
  \[
    \sigma'=\update_R(x,s,\sigma,H_s,q,z),
    \qquad
    H_{s+1}=\operatorname{AppendHist}_R(x,s,H_s,q,z).
  \]
  If $s<r(x)$ the next configuration is the query phase for call $s+1$ with
  state $\sigma'$ and history $H_{s+1}$; if $s=r(x)$ it is the finish countdown
  phase with state $\sigma'$ and full history $H_{s+1}$.

  On a finish phase
  $\operatorname{FinishPhase}(\sigma,H,c_{\rm fin})$, if the finish countdown
  has not halted, the step decrements $c_{\rm fin}$ inside the same phase; if
  it has halted, the step map fixes that halted finish configuration.

  These clauses are the complete definition of the step map.  The oracle-phase
  clause invokes the actual step map and output relation of the supplied
  program $A_Q$ on the actual query $q$; no oracle answer is provided as an
  external parameter to the composed program.  Lean-facing, the declaration
  $\operatorname{def\_composed\_oracle\_step\_map}$ must return the function
  $\operatorname{ComposedStep}_{R,A_Q}(x)$ itself.  It is not a structure with
  a field named $\operatorname{ComposedStep}$, $\operatorname{step}$,
  $\operatorname{next}$, or any synonym, and it is not a proposition asserting
  the existence of a function satisfying branch equations.  The five branch
  equations above are definitional equalities or projection lemmas about this
  single returned function; they are not caller-supplied axioms and they do not
  authorize existential branch choices.
\end{definition}

\begin{definition}[Composed oracle-reduction halt and output]
  \label{def:composed-oracle-halt-output}
  \uses{def:composed-oracle-step-map, def:composed-oracle-finish-phase-machine}
  The halt predicate of the composed program is true exactly on halted finish
  countdown records:
  \[
    \operatorname{ComposedHalt}_{R,A_Q}
      \bigl(x,\operatorname{FinishPhase}(\sigma,H,c_{\rm fin})\bigr)
    \quad\Longleftrightarrow\quad
    c_{\rm fin}=0,
  \]
  and it is false on $\operatorname{InitPhase}$, $\operatorname{QueryPhase}$,
  $\operatorname{OraclePhase}$, and $\operatorname{UpdatePhase}$ records.
  The output relation at a halted finish record is
  \[
    \operatorname{ComposedOut}_{R,A_Q}
      \bigl(x,\operatorname{FinishPhase}(\sigma,H,0),y\bigr)
    \quad\Longleftrightarrow\quad
    y=\finish_R(x,\sigma,H).
  \]
  Thus the only possible returned $P$-output is the final-postprocessing value
  for the state and full history stored in the halted finish phase.
\end{definition}

\begin{definition}[Composed oracle-reduction operational semantics]
  \label{def:composed-oracle-operational-semantics}
  \uses{def:problem-algorithm, def:composed-oracle-program-carrier, def:composed-oracle-initial-configuration, def:composed-oracle-step-map, def:composed-oracle-halt-output}
  \lean{def_composed_oracle_operational_semantics}
  The deterministic operational semantics used by the composed substitution
  $R[A_Q]$ is a concrete semantics for the source problem $P$, denoted
  \[
    \mathsf{Sem}_{R,A_Q}.
  \]
  For each valid $P$-input $x$, its configuration carrier, initial
  configuration, step map, halt predicate, and output relation are exactly
  \[
  \begin{aligned}
    &\operatorname{Config}_{\mathsf{Sem}_{R,A_Q}}(x)
      =\operatorname{ComposedConfig}_{R,A_Q}(x),\\
    &\operatorname{init}_{\mathsf{Sem}_{R,A_Q}}(x)
      =\operatorname{ComposedInit}_{R,A_Q}(x),\\
    &\operatorname{step}_{\mathsf{Sem}_{R,A_Q}}(x)
      =\operatorname{ComposedStep}_{R,A_Q}(x),\\
    &\operatorname{Halt}_{\mathsf{Sem}_{R,A_Q}}(x,c)
      =\operatorname{ComposedHalt}_{R,A_Q}(x,c),\\
    &\operatorname{Out}_{\mathsf{Sem}_{R,A_Q}}(x,c,y)
      =\operatorname{ComposedOut}_{R,A_Q}(x,c,y).
  \end{aligned}
  \]
  Lean-facing, this node is the declaration
  $\operatorname{def\_composed\_oracle\_operational\_semantics}$ and returns a
  value of the existing deterministic-operational-semantics type for $P$ from
  Definition~\ref{def:deterministic-operational-semantics}.  It is not a new
  structure packaging a semantics together with a program code, and it does not
  introduce any declaration named
  $\operatorname{def\_composed\_oracle\_program\_code}$.  The five displayed
  component equations are the complete semantics of the composed run.
\end{definition}

\begin{definition}[Composed oracle-reduction program code]
  \label{def:composed-oracle-program-code}
  \uses{def:problem-algorithm, def:composed-oracle-operational-semantics, def:composed-oracle-program-carrier, def:composed-oracle-initial-configuration, def:composed-oracle-step-map, def:composed-oracle-halt-output}
  \lean{def_composed_oracle_program_code}
  The deterministic $P$-program code underlying the substitution
  $R[A_Q]$ is the single concrete program code inside the operational semantics
  $\mathsf{Sem}_{R,A_Q}$ of
  Definition~\ref{def:composed-oracle-operational-semantics}.  It is denoted
  \[
    \mathsf{Code}_{R,A_Q}
    :=\operatorname{code}(\mathsf{Sem}_{R,A_Q}).
  \]
  Expanding the code inside that semantics gives exactly the tuple
  \[
  \begin{aligned}
    &\operatorname{Config}_{R[A_Q]}(x)
      =\operatorname{ComposedConfig}_{R,A_Q}(x),\\
    &\operatorname{init}_{R[A_Q]}(x)
      =\operatorname{ComposedInit}_{R,A_Q}(x),\\
    &\operatorname{step}_{R[A_Q]}(x)
      =\operatorname{ComposedStep}_{R,A_Q}(x),\\
    &\operatorname{Halt}_{R[A_Q]}(x,c)
      =\operatorname{ComposedHalt}_{R,A_Q}(x,c),\\
    &\operatorname{Out}_{R[A_Q]}(x,c,y)
      =\operatorname{ComposedOut}_{R,A_Q}(x,c,y).
  \end{aligned}
  \]
  Lean-facing, this node is exactly the declaration
  $\operatorname{def\_composed\_oracle\_program\_code}$, and it returns only
  this program-code value.  It is not a structure, record, or proposition with
  fields such as $\operatorname{Semantics}$, $\operatorname{code}$, or a
  ``concrete'' compatibility proof.  No later node may redeclare
  $\operatorname{def\_composed\_oracle\_program\_code}$ as a type of program
  packages; later nodes must refer to the value defined here, or to the
  separate semantics value
  $\operatorname{def\_composed\_oracle\_operational\_semantics}$.
  This declaration contains no correctness proof, no running-time proof, and
  no externally supplied trace or accounting predicate; those are separate
  lemmas about this constructed code.
\end{definition}

\begin{lemma}[Composed oracle-reduction trace characterization]
  \label{lem:composed-oracle-trace-characterization}
  \uses{def:deterministic-operational-semantics, def:composed-oracle-operational-semantics, def:composed-oracle-program-code, def:composed-oracle-run-equations, def:oracle-reduction-local-phases}
  \lean{lem_composed_oracle_trace_characterization}
  For every valid input $x$, step count $S$, and output $y$, the operational
  halting fact
  \[
    \operatorname{HaltsWith}_{R[A_Q]}(x,S,y)
  \]
  holds if and only if there exist a trace
  \[
    t=(q_1,z_1),\ldots,(q_r,z_r),y,
  \]
  a generated state sequence $G:\{0,\ldots,r\}\to\State_R(x)$, oracle
  first-halting-time data $U:\{1,\ldots,r\}\to\N$, and a local cost $L$ such
  that
  \[
    \operatorname{def\_composed\_oracle\_run\_equations}(x,t,G,U),
  \]
  \[
    \operatorname{def\_oracle\_reduction\_local\_phases}(x,t,G,L),
  \]
  and
  \[
    S=L+\sum_{s=1}^{r}U(s).
  \]
  In particular, any such certificate includes the displayed local-phase
  equations
  \[
    G(0)=\init_R(x),\qquad
    G(s)=\update_R(x,s,G(s-1),H(x,t,s),q_s,z_s),
  \]
  the final-output equation
  \[
    y=\finish_R(x,G(r),H(x,t,r+1)),
  \]
  and the actual oracle-run facts
  \[
    \operatorname{HaltsWith}_{A_Q}(q_s,U(s),z_s)
    \qquad(1\leq s\leq r).
  \]
  Equivalently, every halting execution of the composed program is represented
  by exactly such a trace certificate, and any trace certificate satisfying
  $\operatorname{def\_composed\_oracle\_run\_equations}(x,t,G,U)$,
  $\operatorname{def\_oracle\_reduction\_local\_phases}(x,t,G,L)$, with
  step count $S=L+\sum_s U(s)$ is a halting execution of the same composed
  program $R[A_Q]$ with output $y$.
  Here $\operatorname{HaltsWith}_{R[A_Q]}$ is the concrete first-halting
  predicate from Definition~\ref{def:deterministic-operational-semantics}
  applied to the semantics
  $\operatorname{def\_composed\_oracle\_operational\_semantics}$ and to its
  code value
  $\operatorname{def\_composed\_oracle\_program\_code}$.  This lemma introduces
  no new type, structure, or compatibility predicate named
  $\operatorname{def\_composed\_oracle\_program\_code}$; that name is already
  the concrete code value of Definition~\ref{def:composed-oracle-program-code}.
  It is invalid to formalize this trace-characterization node by replacing
  the existing code value with a fresh package whose fields are
  $\operatorname{Semantics}$ and $\operatorname{code}$, or by adding a second
  declaration named $\operatorname{def\_composed\_oracle\_program\_code}$.
  The only formal declaration for this node is the lemma
  $\operatorname{lem\_composed\_oracle\_trace\_characterization}$.
\end{lemma}
\begin{proof}
  \uses{def:deterministic-operational-semantics, def:composed-oracle-operational-semantics, def:composed-oracle-program-code, def:composed-oracle-run-equations, def:oracle-reduction-local-phases}
  Unfold the phase-record transition relation and the certificate predicate of
  Definition~\ref{def:composed-oracle-run-equations}.  Reading a halting
  execution from left to right records exactly the query-answer pairs created
  during the oracle phases, the oracle first-halting times $U(s)$, the states
  generated by the update equations, and the final postprocessing output.  The
  non-oracle phase lengths give the local-cost witness
  $L$ from Definition~\ref{def:oracle-reduction-local-phases}, while the
  oracle phases contribute precisely $\sum_s U(s)$ steps, giving
  $S=L+\sum_s U(s)$.
  Conversely, a trace satisfying
  $\operatorname{def\_composed\_oracle\_run\_equations}(x,t,G,U)$ and the
  local-phase relation determines the same finite sequence of phase records:
  initialization, query construction, the actual first-halting run of $A_Q$ on
  each displayed query, update, and final postprocessing.  Therefore the
  deterministic composed program halts with the displayed output $y$ at the
  displayed step count $S$.
\end{proof}

\begin{lemma}[Composed oracle-reduction step-count accounting]
  \label{lem:composed-oracle-step-count-accounting}
  \uses{def:deterministic-operational-semantics, def:composed-oracle-operational-semantics, def:composed-oracle-program-code, lem:composed-oracle-trace-characterization, def:oracle-reduction-cost, lem:solver-halts-on-promised-input-set}
  \lean{lem_composed_oracle_step_count_accounting}
  In the trace characterization of a halting execution of $R[A_Q]$, if the
  simulated oracle calls halt at first halting times $u_1,\ldots,u_r$, then the
  first halting step count $s$ of the composed execution is exactly
  \[
    s=L+\sum_{i=1}^{r}u_i .
  \]
  If $A_Q$ has worst-case running-time bound $T_Q$, then the same execution
  satisfies the induced cost bound
  \[
    s\leq L+\sum_{i=1}^{r}T_Q(N_Q(q_i),L_Q(q_i)).
  \]
  The step count $s$ is the first-halting time supplied by the
  $\operatorname{HaltsWith}$ predicate of
  Definition~\ref{def:deterministic-operational-semantics} for the composed
  operational semantics and code value of
  Definitions~\ref{def:composed-oracle-operational-semantics} and
  \ref{def:composed-oracle-program-code}.  This accounting lemma does not
  introduce or require any separate program-code wrapper; it uses the same
  concrete code value as
  Lemma~\ref{lem:composed-oracle-trace-characterization}.
\end{lemma}
\begin{proof}
  \uses{def:deterministic-operational-semantics, def:composed-oracle-operational-semantics, def:composed-oracle-program-code, lem:composed-oracle-trace-characterization, def:oracle-reduction-cost, lem:solver-halts-on-promised-input-set}
  Definition~\ref{def:oracle-reduction-local-phases} counts exactly the
  non-oracle initialization, query-construction, update, and postprocessing
  transitions.  The remaining transitions are precisely the simulated
  executions of $A_Q$ on the constructed queries, whose first halting times are
  $u_i$.  Summing the disjoint phase lengths gives the equality for $s$.
  Applying Lemma~\ref{lem:solver-halts-on-promised-input-set} to each valid
  constructed query bounds each $u_i$ by $T_Q(N_Q(q_i),L_Q(q_i))$, yielding the
  displayed inequality.
\end{proof}

\begin{definition}[Composed oracle-reduction algorithm]
  \label{def:composed-oracle-reduction-algorithm}
  \uses{def:composed-oracle-operational-semantics, def:composed-oracle-program-code, def:oracle-reduction-correctness, def:composed-oracle-run-equations, lem:composed-oracle-trace-characterization, lem:composed-oracle-step-count-accounting}
  \lean{def_composed_oracle_reduction_algorithm}
  Let $R$ be a correct oracle reduction datum from $P$ to $Q$ with a local-phase
  presentation, and let $A_Q$ be a concrete correct deterministic algorithm for
  $Q$ with running-time bound $T_Q$.  The composed algorithm
  \[
    R[A_Q]
  \]
  is definitionally the deterministic $P$-program code of
  Definition~\ref{def:composed-oracle-program-code}.  Lean-facing, the
  declaration
  \[
    \operatorname{def\_composed\_oracle\_reduction\_algorithm}(P,Q,R,A_Q)
  \]
  has exactly the displayed construction arguments and returns the program-code
  value $\operatorname{def\_composed\_oracle\_program\_code}$ inside the
  semantics value
  $\operatorname{def\_composed\_oracle\_operational\_semantics}$.  The
  configuration carrier, initial configuration, step map, halt predicate, and
  output relation are respectively the five component equations named in
  Definitions~\ref{def:composed-oracle-operational-semantics} and
  \ref{def:composed-oracle-program-code}.  Consequently the construction has
  no parameter for an alternative $P$-semantics, no parameter for an arbitrary
  $P$-algorithm, no wrapper type of program packages, and no assumptions
  asserting that some externally supplied algorithm is equivalent to the
  composition.

  The separate trace characterization for this constructed code is
  Lemma~\ref{lem:composed-oracle-trace-characterization}.  For every valid
  input $x$, natural number $S$, and output $y$, that lemma states that the
  halting predicate of the returned program holds exactly when there exist one
  trace
  $t=(q_1,z_1),\ldots,(q_{r(x)},z_{r(x)}),y$, one generated state sequence $G$,
  one oracle first-halting-time function $U$, and one local cost $L$ satisfying
  the concrete run-equation certificate
  $\operatorname{def\_composed\_oracle\_run\_equations}(x,t,G,U)$, the local
  phase relation
  $\operatorname{def\_oracle\_reduction\_local\_phases}(x,t,G,L)$, and the
  exact step-count equation
  \[
    S=L+\sum_{s=1}^{r(x)}U(s).
  \]
  Each $U(s)$ is the first halting time of the actual run
  $\operatorname{HaltsWith}_{A_Q}(q_s,U(s),z_s)$ on the actual query $q_s$ from
  the same trace.  Thus oracle answers are not assumed as free data: they are
  outputs of $A_Q$ inside the returned program's step relation.

  Therefore a Lean declaration for this node is the code alias above, with its
  operational projections supplied by
  Definitions~\ref{def:composed-oracle-program-carrier},
  \ref{def:composed-oracle-initial-configuration},
  \ref{def:composed-oracle-step-map},
  \ref{def:composed-oracle-halt-output}, and
  \ref{def:composed-oracle-program-code}.  The trace and accounting statements
  remain the separate lemmas named in the \uses list.  It is invalid to prove
  this node by taking arbitrary parameters named like
  $\operatorname{SP}$, $\operatorname{AP}$,
  $\operatorname{\_configuration\_carrier}$,
  $\operatorname{\_halting\_output}$,
  $\operatorname{\_trace\_semantics}$, or
  $\operatorname{\_step\_count\_governance}$, or by requiring the caller to
  supply proof hypotheses that an unrelated program has the desired carrier,
  trace semantics, output relation, or running-time accounting.  Those
  components are part of the constructed program code above, not external
  obligations.
\end{definition}

\begin{definition}[Composed witness halting fact]
  \label{def:fg-composed-halting-fact}
  \uses{def:composed-oracle-reduction-algorithm}
  Fix $P,Q,R,A_Q$, a valid $P$-input $x$, a step count $s\in\N$, and an output
  $y\in O_P(x)$.  The predicate
  \[
    \operatorname{def\_fg\_composed\_halting\_fact}(P,Q,R,A_Q,x,s,y)
  \]
  is exactly the single proposition
  \[
    \operatorname{HaltsWith}_{R[A_Q]}(x,s,y).
  \]
  The program in this halting fact is the single composed program $R[A_Q]$ of
  Definition~\ref{def:composed-oracle-reduction-algorithm}.  This predicate
  has no argument for an alternative $P$-algorithm, no trace certificate
  argument, and no field that can replace the displayed halting fact.
\end{definition}

\begin{definition}[Composed witness correctness fact]
  \label{def:fg-composed-correctness-fact}
  \uses{def:oracle-reduction-correctness}
  Fix $P,Q,R,A_Q$, a valid $P$-input $x$, a step count $s\in\N$, and an output
  $y\in O_P(x)$.  The predicate
  \[
    \operatorname{def\_fg\_composed\_correctness\_fact}(P,Q,R,A_Q,x,s,y)
  \]
  is exactly the single proposition
  \[
    \Correct_P(x,y).
  \]
  The parameters $R$, $A_Q$, and $s$ are retained only so this predicate has
  the same witness arguments as the corresponding halting fact; they do not
  add any alternative program, trace data, or replaceable correctness field.
\end{definition}

\begin{definition}[Composed witness halting and correctness]
  \label{def:fg-composed-halting-correctness}
  \uses{def:fg-composed-halting-fact, def:fg-composed-correctness-fact}
  Fix $P,Q,R,A_Q$, a valid $P$-input $x$, a step count $s\in\N$, and an output
  $y\in O_P(x)$.  The predicate
  \[
    \operatorname{def\_fg\_composed\_halting\_correctness}(P,Q,R,A_Q,x,s,y)
  \]
  is exactly the two-conjunct proposition
  \[
    \operatorname{def\_fg\_composed\_halting\_fact}(P,Q,R,A_Q,x,s,y)
    \quad\text{and}\quad
    \operatorname{def\_fg\_composed\_correctness\_fact}(P,Q,R,A_Q,x,s,y).
  \]
  Expanding the two helper predicates gives exactly
  \[
    \operatorname{HaltsWith}_{R[A_Q]}(x,s,y)
    \quad\text{and}\quad
    \Correct_P(x,y).
  \]
  The program in the halting conjunct is the single composed program $R[A_Q]$
  of Definition~\ref{def:composed-oracle-reduction-algorithm}.  This predicate
  has no argument for an alternative $P$-algorithm and no field that can
  replace either displayed fact.
\end{definition}

\begin{definition}[Composed witness trace certificate]
  \label{def:fg-composed-trace-certificate}
  \uses{def:composed-oracle-run-equations, lem:composed-oracle-trace-characterization}
  Fix $P,Q,R,A_Q$, a valid $P$-input $x$, an output $y\in O_P(x)$, a trace
  \[
    t=((q_1,z_1),\ldots,(q_{r(x)},z_{r(x)}),y),
  \]
  a generated state sequence $G:\{0,\ldots,r(x)\}\to\State_R(x)$, and oracle
  first-halting-time data $U:\{1,\ldots,r(x)\}\to\N$.  The predicate
  \[
    \operatorname{def\_fg\_composed\_trace\_certificate}(P,Q,R,A_Q,x,y,t,G,U)
  \]
  is exactly the conjunction of
  \[
    \operatorname{def\_composed\_oracle\_run\_equations}(x,t,G,U),
  \]
  \[
    \forall\,1\leq i\leq r(x),\qquad
      \operatorname{HaltsWith}_{A_Q}(q_i,U(i),z_i),
  \]
  and
  \[
    y=\finish_R(x,G(r(x)),H(x,t,r(x)+1)).
  \]
  The entries $q_i,z_i$ are the query and answer entries of the displayed trace
  $t$, and the prefix-history selector $H$ is the one from the fixed
  local-phase presentation of $R$.  The oracle halting facts and the final
  postprocessing equality are conjuncts of this predicate, so they are
  eliminable from a proof for the same witnesses $(x,y,t,G,U)$ rather than
  consequences quoted only through a lemma name.
\end{definition}

\begin{definition}[Composed witness local accounting]
  \label{def:fg-composed-local-accounting}
  \uses{def:oracle-reduction-cost, lem:composed-oracle-step-count-accounting}
  Fix $P,Q,R,A_Q,T_Q$, a valid $P$-input $x$, a step count $s\in\N$, a trace
  $t=((q_1,z_1),\ldots,(q_{r(x)},z_{r(x)}),y)$, a generated state sequence
  $G:\{0,\ldots,r(x)\}\to\State_R(x)$, oracle first-halting-time data
  $U:\{1,\ldots,r(x)\}\to\N$, and a local cost $L\in\N$.  The predicate
  \[
    \operatorname{def\_fg\_composed\_local\_accounting}
      (P,Q,R,A_Q,T_Q,x,s,t,G,U,L)
  \]
  is exactly the conjunction of
  \[
    \operatorname{def\_oracle\_reduction\_local\_phases}(x,t,G,L),
  \]
  \[
    s=L+\sum_{i=1}^{r(x)}U(i),
  \]
  and
  \[
    s\leq
      L+\sum_{i=1}^{r(x)}T_Q(N_Q(q_i),L_Q(q_i)).
  \]
  Thus the local-phase certificate, the exact first-halting step-count
  equation, and the induced oracle-cost upper bound are separate eliminable
  conjuncts for the same trace $t$, state sequence $G$, oracle time data $U$,
  and local cost $L$.
\end{definition}

\begin{definition}[Composed witness subquadratic transfer bound]
  \label{def:fg-composed-transfer-bound}
  \uses{def:subquadratic-oracle-query-budget}
  Fix $P,Q,R$, constants $\eps'>0$, $C'>0$, and $b'\in\N$, a valid $P$-input
  $x$, and a step count $s\in\N$.  The predicate
  \[
    \operatorname{def\_fg\_composed\_transfer\_bound}(P,Q,R,\eps',C',b',x,s)
  \]
  is the single inequality
  \[
    s\leq C' N_P(x)^{2-\eps'}(1+L_P(x))^{b'}.
  \]
  This is only the final transferred $P$-running-time bound.  It contains no
  trace data, no local-cost data, and no oracle-run data; those are the
  separate obligations of
  Definitions~\ref{def:fg-composed-trace-certificate} and
  \ref{def:fg-composed-local-accounting}.
\end{definition}

\begin{definition}[Fine-grained composed witness obligations]
  \label{def:fine-grained-composed-witness-obligations}
  \uses{def:fg-composed-halting-correctness, def:fg-composed-trace-certificate, def:fg-composed-local-accounting, def:fg-composed-transfer-bound}
  Fix $P,Q,R,A_Q,T_Q$ and constants $\eps'>0$, $C'>0$, and $b'\in\N$.  The
  predicate
  \[
    \operatorname{def\_fine\_grained\_composed\_witness\_obligations}
      (P,Q,R,A_Q,T_Q,\eps',C',b')
  \]
  is a proposition with exactly these displayed arguments and no additional
  program, semantics, trace-characterization, accounting, or helper-proof
  parameters.  It holds exactly when, for every valid $P$-input $x$, there
  exist concrete data
  \[
    s\in\N,\quad y\in O_P(x),\quad
    t=((q_1,z_1),\ldots,(q_{r(x)},z_{r(x)}),y),
  \]
  a generated state sequence $G:\{0,\ldots,r(x)\}\to\State_R(x)$, oracle
  first-halting-time data $U:\{1,\ldots,r(x)\}\to\N$, and a local cost
  $L\in\N$ such that all four one-to-one witness obligations hold for those
  same witnesses:
  \[
    \operatorname{def\_fg\_composed\_halting\_correctness}
      (P,Q,R,A_Q,x,s,y),
  \]
  \[
    \operatorname{def\_fg\_composed\_trace\_certificate}
      (P,Q,R,A_Q,x,y,t,G,U),
  \]
  \[
    \operatorname{def\_fg\_composed\_local\_accounting}
      (P,Q,R,A_Q,T_Q,x,s,t,G,U,L),
  \]
  and
  \[
    \operatorname{def\_fg\_composed\_transfer\_bound}
      (P,Q,R,\eps',C',b',x,s).
  \]
  Expanding these four helper predicates gives exactly the original displayed
  universal-existential proposition: the composed program halts with output
  $y$, that output satisfies $\Correct_P(x,y)$, the trace satisfies
  $\operatorname{def\_composed\_oracle\_run\_equations}(x,t,G,U)$, the same
  witnesses satisfy
  $\operatorname{def\_oracle\_reduction\_local\_phases}(x,t,G,L)$,
  \[
    s=L+\sum_{i=1}^{r(x)}U(i),
  \]
  each oracle answer satisfies
  $\operatorname{HaltsWith}_{A_Q}(q_i,U(i),z_i)$, the final output satisfies
  \[
    y=\finish_R(x,G(r(x)),H(x,t,r(x)+1)),
  \]
  the induced oracle-cost bound
  \[
    s\leq
      L+\sum_{i=1}^{r(x)}T_Q(N_Q(q_i),L_Q(q_i))
  \]
  holds, and the transferred subquadratic bound
  \[
    s\leq C' N_P(x)^{2-\eps'}(1+L_P(x))^{b'}
  \]
  holds.  The helper predicates do not add a new program, trace, state
  sequence, or local-cost witness; they only name four individually
  formalizable conjunct groups of the same proposition.
\end{definition}

\begin{definition}[Fine-grained reduction]
  \label{def:fine-grained-reduction}
  \uses{def:problem-algorithm, def:quadratic-baseline-algorithm, def:oracle-reduction-correctness, def:subquadratic-oracle-query-budget, def:oracle-reduction-cost, def:composed-oracle-reduction-algorithm, def:fine-grained-composed-witness-obligations}
  The Lean-facing fine-grained reduction predicate has the explicit arguments
  \[
    \operatorname{def\_fine\_grained\_reduction}
      (P,Q,h_P,h_Q),
  \]
  where $P$ and $Q$ are computational problem specifications and
  \[
    h_P:\operatorname{def\_quadratic\_baseline\_algorithm}(P),
    \qquad
    h_Q:\operatorname{def\_quadratic\_baseline\_algorithm}(Q)
  \]
  are required baseline-algorithm witnesses from
  Definition~\ref{def:quadratic-baseline-algorithm}.  The predicate is not
  defined for bare problem specifications with no baseline witnesses.

  With those baseline witnesses fixed, $P$ reduces to $Q$ in the fine-grained
  subquadratic sense if, for every constant $\eps>0$ and every particular
  correct deterministic $Q$-algorithm $A_Q$ whose actual operational
  running-time bound satisfies
  \[
    T_Q(n',\ell')\leq C(n')^{2-\eps}(1+\ell')^b
  \]
  on valid $Q$-instances, there exist an oracle reduction datum $R$ from
  $P$ to $Q$, a local-phase presentation of $R$, and constants
  $\eps'>0$, $C'>0$, and $b'\in\N$ such that all of the following obligations
  hold for this same displayed tuple $(R,A_Q,T_Q,\eps',C',b')$.

  First, $R$ is a correct oracle reduction: it sends valid $P$-inputs to valid
  $Q$-queries and every locally generated trace with correct oracle answers
  has a final output satisfying $\Correct_P$, exactly as in
  Definition~\ref{def:oracle-reduction-correctness}.  Second, $R$ has the
  subquadratic transfer budget of
  Definition~\ref{def:subquadratic-oracle-query-budget} against the displayed
  running-time bound $T_Q$.  Third, the $P$-algorithm witnessing the reduction
  is definitionally the composed algorithm
  \[
    A_P=R[A_Q]
  \]
  from Definition~\ref{def:composed-oracle-reduction-algorithm}.  Its program
  code, run relation, output, and step count are the concrete substitution
  simulation of $R$ using the actual runs of $A_Q$ on the actual constructed
  oracle queries.  The reduction is not allowed to existentially choose any
  other $P$-algorithm $A_P$.

  Fourth, the composed witness obligations of
  Definition~\ref{def:fine-grained-composed-witness-obligations} hold for that
  same $R[A_Q]$.  Thus the final subquadratic correctness and time guarantee is
  not a separate existential consequence: it is the halting, correctness,
  final-postprocessing, oracle-answer, and step-count accounting certificate of
  the composed algorithm constructed in the third obligation.  Problems are
  subquadratic equivalent if the same baseline-witnessed reduction predicate
  holds in both directions.
\end{definition}

\begin{definition}[$k$SAT decision problem]
  \label{def:ksat}
  For positive integers $k$ and $n$, a $k$CNF formula on variables
  $x_1,\ldots,x_n$ is a finite conjunction of clauses, where each clause is a
  disjunction of at most $k$ literals and each literal is either $x_i$ or
  $\neg x_i$ for some $i$.  An assignment is a function
  $\sigma:\{1,\ldots,n\}\to\{0,1\}$.  It satisfies a positive literal $x_i$ when
  $\sigma(i)=1$, satisfies a negative literal $\neg x_i$ when $\sigma(i)=0$,
  satisfies a clause when it satisfies at least one literal in that clause, and
  satisfies the formula when it satisfies every clause.  The decision problem
  $k\SAT$ asks whether a given $k$CNF formula has a satisfying assignment.
\end{definition}

\begin{conjecture}[Strong Exponential Time Hypothesis, Definition 2.4/A.1]
  \label{conj:seth}
  \uses{def:ksat, def:problem-algorithm, def:exact-worst-case-running-time}
  For every $\eps>0$, there is a positive integer $k$ such that no correct
  deterministic decision algorithm for $k\SAT$ decides all valid $k$CNF
  formulas with $n$ variables in time
  $O(2^{(1-\eps)n}\poly(M))$, uniformly for every encoded-size parameter
  $M\in\N$, for all sufficiently large $n$.

  Lean-facing, this conjecture quantifies only over program codes $A$ for the
  $k\SAT$ computational-problem specification together with semantic
  worst-case running-time bounds $T$ in the sense of
  Definitions~\ref{def:problem-algorithm} and
  \ref{def:exact-worst-case-running-time}.  The forbidden fast-algorithm
  statement is:
  there do not exist $A,T,C,b,n_0$ such that
  $A$ halts on every valid $k\SAT$ input, returns $1$ exactly when the formula
  is satisfiable and $0$ exactly when it is unsatisfiable, $T$ bounds the actual
  first halting time of that correct execution on every valid input, and
  \[
    T(n,M)\leq C\,2^{(1-\eps)n}(1+M)^b
  \]
  for every $n\geq n_0$ and every natural number $M\in\N$, including size
  parameters for which no valid encoded $k$CNF formula with $n$ variables
  exists.  For valid inputs, the relevant value of $M$ is the input's actual
  encoded formula size; the quantified running-time bound itself is a uniform
  bound on the full size-parameter function $T(n,M)$, not only on the subset of
  parameter pairs realized by valid formulas.
  Incorrect programs, partial programs, and arbitrary exact-time functions
  unattached to a correct halting execution are outside the quantified
  obstruction.
\end{conjecture}

\begin{definition}[Orthogonal Vectors, Definition 2.5/A.2]
  \label{def:ov}
  Given binary vectors $v_1,\ldots,v_n\in\{0,1\}^{\ell}$, the problem
  $\OV_{n,\ell}$ asks whether there are indices $i\neq j$ with
  $\ip{v_i}{v_j}=0$.
\end{definition}

\begin{conjecture}[Orthogonal Vectors Conjecture, Conjecture 2.6/A.3]
  \label{conj:ovc}
  \uses{def:ov, def:problem-algorithm, def:exact-worst-case-running-time}
  For every $\eps>0$ there exists a constant $c>0$ such that
  $\OV_{n,\ell_c(n)}$ cannot be solved in $O(n^{2-\eps})$ time, where $c$ is a
  positive real constant and the integer dimension schedule is
  \[
    \ell_c(n)=\max\{1,\lceil c\log_2(\max\{n,2\})\rceil\}.
  \]
  Equivalently, the paper's notation $c\log n$ means this real-constant
  logarithmic dimension up to the displayed integer rounding, not an
  integer-valued constant multiplied by an integer logarithm.

  Lean-facing, this conjecture is the direct negation of the existence of a
  correct deterministic $\OV$ decision program with the displayed semantic
  worst-case bound.  For each $\eps>0$ there is a real number $c$ with
  $c>0$ such that there do not exist a program code $A$, a worst-case
  running-time bound
  $T:\N\times\N\to\R_{\geq0}$, constants $C>0$ and $n_0$, and the correctness
  and halting witnesses from Definitions~\ref{def:problem-algorithm} and
  \ref{def:exact-worst-case-running-time}, satisfying all of the following:
  for every valid $\OV$ input with $n$ vectors in the rounded dimension
  $\ell_c(n)$ above, $A$ halts and returns $1$ exactly when some distinct pair
  has inner product $0$ and returns $0$ exactly otherwise, $T$ bounds the
  actual first halting time of that execution, and
  \[
    T(n,\ell_c(n))\leq C n^{2-\eps}
  \]
  for every $n\geq n_0$.  The constant $c$ is not a natural number, and
  $\ell_c(n)$ is not the narrower expression
  $c_{\N}\cdot\operatorname{Nat.log2}(n)$ for some
  $c_{\N}\in\N$; any formalization must retain the positive real constant and
  the explicit ceiling/max conversion to a positive integer dimension.
  No additional premise such as
  $\operatorname{def\_truly\_subquadratic}(T)$ is part of the forbidden
  fast-algorithm package.  The predicate
  $\operatorname{def\_truly\_subquadratic}$ may be derived from the displayed
  big-$O$ bound when needed, but it is not a separate hypothesis that weakens
  the conjecture.
\end{conjecture}

\begin{theorem}[SETH implies OVC]
  \label{thm:seth-implies-ovc}
  \uses{conj:seth, conj:ovc}
  \notready
  If SETH holds, then OVC holds.
\end{theorem}
% TODO: The paper cites Williams (2005) for this implication and gives no proof.

\begin{definition}[Bichromatic Orthogonal Vectors]
  \label{def:bichromatic-ov}
  \uses{def:ov}
  In bichromatic $\OV_{n,\ell}$ one is given
  $A=\{a_1,\ldots,a_n\}$ and $B=\{b_1,\ldots,b_n\}$ in $\{0,1\}^{\ell}$ and
  asks whether some pair $(a_i,b_j)$ satisfies $\ip{a_i}{b_j}=0$.
\end{definition}

\begin{definition}[Monochromatic-to-bichromatic OV reduction datum]
  \label{def:mono-to-bi-ov-reduction-datum}
  \uses{def:ov, def:bichromatic-ov, def:oracle-reduction-data}
  The reduction datum $R_{\rm mono\to bi}$ from monochromatic $\OV$ to
  bichromatic $\OV$ acts on a monochromatic input
  $V=(v_1,\ldots,v_n)\subseteq\{0,1\}^{\ell}$ as the following concrete
  one-query-or-zero-query local-phase reduction.  The fine-grained regime is
  $n\geq2$; the finitely many inputs with $n<2$ are decided by the constant
  zero-query branch by directly returning no, since no distinct pair of indices
  exists.

  First it computes the Boolean value
  \[
    Z(V)\quad\Longleftrightarrow\quad
    \exists i\in\{1,\ldots,n\},\ \forall r\in\{1,\ldots,\ell\},\ v_i[r]=0
  \]
  by the deterministic nested scan over all entries of the input matrix.  The
  scan step count is exactly bounded by a fixed constant multiple of
  $n(1+\ell)$, because each of the $n\ell$ Boolean coordinates is inspected at
  most once and the per-row accumulator is updated in constant time.

  If $n\geq2$ and $Z(V)$ is true, then the oracle-call count is
  $r(V)=0$, there is no bichromatic query, and the final postprocessing output
  is the yes answer for monochromatic $\OV$.  If $Z(V)$ is false, then the
  oracle-call count is $r(V)=1$ and the unique query is the concrete
  bichromatic instance
  \[
    q_1(V)=(A(V),B(V)),\qquad A(V)=V,\qquad B(V)=V,
  \]
  with principal size $N_Q(q_1(V))=n$ and auxiliary dimension
  $L_Q(q_1(V))=\ell$.  The copy phase writes the two lists $A(V)$ and $B(V)$
  coordinatewise, so its step count is also bounded by a fixed constant
  multiple of $n(1+\ell)$.  The update phase merely stores the oracle answer
  $z_1\in\{0,1\}$, and the final postprocessing returns exactly $z_1$ as the
  monochromatic $\OV$ answer.  Thus every locally generated trace has local
  non-oracle cost
  \[
    L\leq C_0\,n(1+\ell)
  \]
  for an absolute constant $C_0$, and the oracle-cost sum is either $0$ in the
  zero-vector branch or exactly
  $T_Q(n,\ell)$ in the one-query branch.  These scan, query-size,
  postprocessing, and cost equalities are part of the datum; they are not
  deferred to a later opaque reduction predicate.
\end{definition}

\begin{lemma}[Monochromatic-to-bichromatic OV query correctness]
  \label{lem:mono-to-bi-ov-query-correctness}
  \uses{def:mono-to-bi-ov-reduction-datum, def:ov, def:bichromatic-ov, def:inner-norm}
  For every monochromatic OV input $V=(v_1,\ldots,v_n)$ with $n\geq2$, the two
  branches of $R_{\rm mono\to bi}$ preserve the yes/no condition as follows.
  If $Z(V)$ is true, then $V$ is a yes instance of monochromatic $\OV$.  If
  $Z(V)$ is false, then the unique query $(V,V)$ is a yes instance of
  bichromatic $\OV$ if and only if $V$ is a yes instance of monochromatic
  $\OV$.
\end{lemma}
\begin{proof}
  \uses{def:mono-to-bi-ov-reduction-datum, def:ov, def:bichromatic-ov, def:inner-norm}
  If $Z(V)$ is true, choose a zero row $v_i$.  Since $n\geq2$, choose an index
  $j\neq i$.  Then $\ip{v_i}{v_j}=0$, so the monochromatic instance is yes.
  If $Z(V)$ is false, then every vector is nonzero.  For a binary vector,
  nonzero means it has some coordinate equal to $1$, and hence
  $\ip{v_i}{v_i}\geq1$.  Thus no self-pair in the bichromatic query $(V,V)$ is
  orthogonal.  A bichromatic orthogonal pair must therefore use two distinct
  original rows, giving a monochromatic OV witness.  Conversely any distinct
  monochromatic OV witness is exactly the same cross-pair in the query
  $(V,V)$.
\end{proof}

\begin{lemma}[Monochromatic-to-bichromatic OV branch answer transfer]
  \label{lem:mono-to-bi-ov-branch-answer-transfer}
  \uses{def:mono-to-bi-ov-reduction-datum, lem:mono-to-bi-ov-query-correctness, def:oracle-reduction-correctness}
  For every valid monochromatic OV input $V=(v_1,\ldots,v_n)$, the
  oracle-answer correctness obligation for $R_{\rm mono\to bi}$ is exactly
  branch-sensitive.

  If $n<2$, then $r(V)=0$, the final postprocessing output is $0$, and this is
  the correct monochromatic OV answer because no distinct pair of input rows
  exists.  If $n\geq2$ and $Z(V)$ is true, then $r(V)=0$, there is no query and
  no oracle answer to quantify over, the final postprocessing output is $1$,
  and this is the correct monochromatic OV answer.

  If $n\geq2$ and $Z(V)$ is false, then $r(V)=1$ and the unique query is
  \[
    q_1(V)=(V,V).
  \]
  In this one-query branch, and only in this branch, every correct yes/no
  oracle answer $z_1$ for the bichromatic query $q_1(V)$ is the correct
  yes/no answer for the original monochromatic instance $V$.  Equivalently,
  in the nonzero branch the final postprocessing value
  $\finish_{R_{\rm mono\to bi}}(V,\sigma_1,H_2)=z_1$ satisfies the
  monochromatic OV correctness relation whenever $z_1$ satisfies the
  bichromatic OV correctness relation on $q_1(V)$.

  The zero-vector branch has no proposition of the form ``every correct oracle
  answer on $q_1(V)$ gives the correct original answer'', because $q_1(V)$ is
  not constructed there.  In particular, the correctness proof must not use
  the self-query $(V,V)$ in the zero-vector branch, where that query could
  contain self-pair artefacts and is irrelevant to the reduction run.
\end{lemma}
\begin{proof}
  \uses{def:mono-to-bi-ov-reduction-datum, lem:mono-to-bi-ov-query-correctness}
  The branch equations for $r(V)$, $q_1(V)$, and final postprocessing are part
  of Definition~\ref{def:mono-to-bi-ov-reduction-datum}.  For $n<2$, the
  monochromatic OV witness predicate is empty because it requires distinct
  indices.  For $n\geq2$ and $Z(V)$ true,
  Lemma~\ref{lem:mono-to-bi-ov-query-correctness} proves that $V$ is a yes
  instance, so the direct output $1$ is correct without any oracle call.  For
  $n\geq2$ and $Z(V)$ false, the same lemma states that the unique query
  $(V,V)$ is a yes instance exactly when the original instance is a yes
  instance.  Therefore a correct Boolean answer to that unique query is exactly
  the correct Boolean answer for $V$, and the update/finish clauses return
  that same answer.
\end{proof}

\begin{lemma}[Monochromatic-to-bichromatic OV local cost transfer]
  \label{lem:mono-to-bi-ov-local-cost-transfer}
  \uses{def:mono-to-bi-ov-reduction-datum, def:subquadratic-oracle-query-budget, def:oracle-reduction-cost, def:truly-subquadratic}
  Suppose a bichromatic $\OV$ algorithm has a semantic running-time bound
  \[
    T_Q(n',\ell')\leq C(n')^{2-\eps}(1+\ell')^b
  \]
  on valid bichromatic inputs.  For every locally generated trace of
  $R_{\rm mono\to bi}$ on a monochromatic input of size $(n,\ell)$, the
  composed reduction run has total time
  \[
    L+\sum_s T_Q(N_Q(q_s),L_Q(q_s))
    \leq
    C_0n(1+\ell)+C n^{2-\eps}(1+\ell)^b.
  \]
  Consequently, for every fixed $\eps>0$ there are constants
  $\eps'>0$, $C'>0$, and $b'\in\N$ such that the same total time is at most
  \[
    C'n^{2-\eps'}(1+\ell)^{b'}
  \]
  for all sufficiently large $n$.  In particular, a truly subquadratic
  bichromatic OV solver transfers to a truly subquadratic monochromatic OV
  solver with only polynomial overhead in $\ell$.
\end{lemma}
\begin{proof}
  \uses{def:mono-to-bi-ov-reduction-datum, def:oracle-reduction-cost, def:truly-subquadratic}
  The datum gives $L\leq C_0n(1+\ell)$.  In the zero-vector branch the oracle
  sum is empty.  In the nonzero branch there is exactly one query, and its
  size and dimension are exactly $(n,\ell)$, so the oracle contribution is at
  most $T_Q(n,\ell)\leq Cn^{2-\eps}(1+\ell)^b$.  This proves the first
  displayed inequality.  Taking for example $\eps'=\eps/2$ and
  $b'=\max(b,1)$, the linear local term is absorbed into
  $C'n^{2-\eps'}(1+\ell)^{b'}$ for sufficiently large $n$, and the oracle term
  is bounded by the same expression after increasing $C'$.
\end{proof}

\begin{lemma}[Monochromatic-to-bichromatic OV reduction correctness and cost]
  \label{lem:mono-to-bi-ov-fine-grained-reduction}
  \uses{def:mono-to-bi-ov-reduction-datum, lem:mono-to-bi-ov-query-correctness, lem:mono-to-bi-ov-branch-answer-transfer, lem:mono-to-bi-ov-local-cost-transfer, def:fine-grained-reduction}
  For any fixed baseline witnesses for monochromatic $\OV$ and bichromatic
  $\OV$, the datum $R_{\rm mono\to bi}$ satisfies
  \[
    \operatorname{def\_fine\_grained\_reduction}
      (\OV,\operatorname{BiOV},h_{\rm mono},h_{\rm bi}).
  \]
  In particular, substituting any truly subquadratic bichromatic $\OV$
  algorithm into the composed program
  $R_{\rm mono\to bi}[A_{\rm bi}]$ gives the concrete monochromatic $\OV$
  algorithm certified by Definition~\ref{def:fine-grained-reduction}.  More
  explicitly, the witness for this reduction includes the following eliminable
  obligations:
  \[
    r(V)=0\ \text{and final output }1\quad\text{when }n\geq2\text{ and }Z(V),
  \]
  \[
    r(V)=1,\quad q_1(V)=(V,V),\quad
    N_Q(q_1(V))=n,\quad L_Q(q_1(V))=\ell
    \quad\text{when }\neg Z(V),
  \]
  the zero-vector scan and copy construction have local cost at most
  $C_0n(1+\ell)$, and the oracle-answer correctness obligation is exactly the
  branch-sensitive obligation of
  Lemma~\ref{lem:mono-to-bi-ov-branch-answer-transfer}: the $n<2$ and
  zero-vector branches have no query and their direct final outputs are
  correct, while in the nonzero one-query branch every correct oracle answer
  on the displayed query $q_1(V)=(V,V)$ gives the correct monochromatic
  $\OV$ answer.  There is no universal correctness obligation for an oracle
  answer on $(V,V)$ in the zero-vector branch.  Finally, any oracle bound
  $T_Q(n',\ell')\leq C(n')^{2-\eps}(1+\ell')^b$ transfers to a composed
  running-time bound
  $C'n^{2-\eps'}(1+\ell)^{b'}$ for fixed constants
  $\eps'>0,C'>0,b'$.
\end{lemma}
\begin{proof}
  \uses{def:mono-to-bi-ov-reduction-datum, lem:mono-to-bi-ov-query-correctness, lem:mono-to-bi-ov-branch-answer-transfer, lem:mono-to-bi-ov-local-cost-transfer, def:fine-grained-reduction}
  Lemma~\ref{lem:mono-to-bi-ov-query-correctness} proves the two correctness
  branches: the zero-vector branch is a direct yes instance, and in the
  nonzero branch the query $(V,V)$ has a bichromatic orthogonal pair exactly
  when the original instance has a distinct orthogonal pair.  The datum itself
  supplies the branch equations $r(V)=0$ or $r(V)=1$, the query equation
  $q_1(V)=(V,V)$, and the size equations
  $N_Q(q_1(V))=n$, $L_Q(q_1(V))=\ell$.
  Lemma~\ref{lem:mono-to-bi-ov-branch-answer-transfer} supplies the formal
  oracle-answer transfer in precisely the one-query branch and supplies the
  direct-output correctness in the two zero-query branches.
  Lemma~\ref{lem:mono-to-bi-ov-local-cost-transfer} supplies the local scan and
  copy cost, the one-query oracle-cost accounting, and the preservation of a
  truly subquadratic exponent after composition.  Therefore the oracle
  reduction correctness, subquadratic transfer budget, and composed witness
  obligations required by Definition~\ref{def:fine-grained-reduction} all hold
  for the concrete composed program $R_{\rm mono\to bi}[A_{\rm bi}]$.
\end{proof}

\begin{definition}[Bichromatic-to-monochromatic OV reduction datum]
  \label{def:bi-to-mono-ov-reduction-datum}
  \uses{def:ov, def:bichromatic-ov, def:oracle-reduction-data}
  The reduction datum $R_{\rm bi\to mono}$ from bichromatic $\OV$ to
  monochromatic $\OV$ maps a bichromatic input
  $A=\{a_1,\ldots,a_n\}$ and $B=\{b_1,\ldots,b_n\}$ in $\{0,1\}^{\ell}$ to the
  single monochromatic instance in dimension $\ell+2$ consisting of
  \[
    a'_i=(a_i,1,0),\qquad b'_j=(b_j,0,1).
  \]
  It queries monochromatic $\OV$ on the multiset
  $A'\cup B'$ and returns exactly the oracle's yes/no answer.  The construction
  doubles the number of vectors, increases the dimension by $2$, and has only
  linear local overhead in the input encoding size.
\end{definition}

\begin{lemma}[Bichromatic-to-monochromatic OV query equivalence]
  \label{lem:bi-to-mono-ov-query-equivalence}
  \uses{def:bi-to-mono-ov-reduction-datum, def:ov, def:bichromatic-ov, def:inner-norm}
  Let $A=\{a_1,\ldots,a_n\}$ and $B=\{b_1,\ldots,b_n\}$ be a valid bichromatic
  OV input in dimension $\ell$, and let
  \[
    A'=\{(a_i,1,0):1\leq i\leq n\},\qquad
    B'=\{(b_j,0,1):1\leq j\leq n\}
  \]
  be the transformed rows in the single monochromatic query
  $Q(A,B)=A'\cup B'$.  Then the query has an orthogonal pair of distinct rows
  if and only if the original bichromatic input has a cross orthogonal pair:
  \[
    \exists u\neq v\in Q(A,B),\ \ip{u}{v}=0
    \quad\Longleftrightarrow\quad
    \exists i,j,\ \ip{a_i}{b_j}=0.
  \]
  More granularly, for all $i,i'$ and $j,j'$,
  \[
    \ip{(a_i,1,0)}{(a_{i'},1,0)}=\ip{a_i}{a_{i'}}+1\geq1,
  \]
  \[
    \ip{(b_j,0,1)}{(b_{j'},0,1)}=\ip{b_j}{b_{j'}}+1\geq1,
  \]
  and
  \[
    \ip{(a_i,1,0)}{(b_j,0,1)}=\ip{a_i}{b_j}.
  \]
  Thus every orthogonal witness in the monochromatic query is necessarily a
  cross pair with one row from $A'$ and one row from $B'$.
\end{lemma}
\begin{proof}
  \uses{def:bi-to-mono-ov-reduction-datum, def:ov, def:bichromatic-ov, def:inner-norm}
  Expanding the three appended coordinates gives the displayed inner-product
  equalities.  Same-side transformed pairs receive an additional contribution
  $1$ from the side marker coordinate, so they cannot be orthogonal.  Cross
  pairs receive no contribution from the two marker coordinates, so their
  transformed inner product is exactly the original inner product.  Therefore
  any monochromatic orthogonal pair must be cross and gives a bichromatic
  witness, while any bichromatic witness gives the corresponding cross
  orthogonal pair in the query.
\end{proof}

\begin{lemma}[Bichromatic-to-monochromatic OV oracle answer transfer]
  \label{lem:bi-to-mono-ov-oracle-answer-transfer}
  \uses{def:bi-to-mono-ov-reduction-datum, lem:bi-to-mono-ov-query-equivalence, def:oracle-reduction-correctness}
  For the reduction $R_{\rm bi\to mono}$ on a valid bichromatic input
  $(A,B)$, the oracle-call count is $1$, the unique oracle query is the
  monochromatic instance $Q(A,B)=A'\cup B'$ from
  Definition~\ref{def:bi-to-mono-ov-reduction-datum}, and final
  postprocessing returns exactly the oracle answer $z_1$.

  Consequently, every correct yes/no answer $z_1$ for the monochromatic query
  $Q(A,B)$ is the correct yes/no answer for the original bichromatic OV input
  $(A,B)$.  Equivalently, if $z_1=1$, then
  $Q(A,B)$ has an orthogonal pair and hence $(A,B)$ has a cross orthogonal
  pair; if $z_1=0$, then $Q(A,B)$ has no orthogonal pair and hence $(A,B)$ has
  no cross orthogonal pair.
\end{lemma}
\begin{proof}
  \uses{def:bi-to-mono-ov-reduction-datum, lem:bi-to-mono-ov-query-equivalence}
  The query-count, query-construction, and final-postprocessing equations are
  the concrete clauses of Definition~\ref{def:bi-to-mono-ov-reduction-datum}.
  Lemma~\ref{lem:bi-to-mono-ov-query-equivalence} identifies the yes condition
  of the query with the yes condition of the original bichromatic input.  A
  correct Boolean oracle answer is therefore true exactly in the original yes
  case and false exactly in the original no case, and the final
  postprocessing returns that same Boolean value.
\end{proof}

\begin{lemma}[Bichromatic-to-monochromatic OV reduction correctness and cost]
  \label{lem:bi-to-mono-ov-fine-grained-reduction}
  \uses{def:bi-to-mono-ov-reduction-datum, lem:bi-to-mono-ov-query-equivalence, lem:bi-to-mono-ov-oracle-answer-transfer, def:fine-grained-reduction}
  For any fixed baseline witnesses for bichromatic $\OV$ and monochromatic
  $\OV$, the datum $R_{\rm bi\to mono}$ satisfies
  \[
    \operatorname{def\_fine\_grained\_reduction}
      (\operatorname{BiOV},\OV,h_{\rm bi},h_{\rm mono}).
  \]
  In particular, substituting any truly subquadratic monochromatic $\OV$
  algorithm into the composed program
  $R_{\rm bi\to mono}[A_{\rm mono}]$ gives the concrete bichromatic $\OV$
  algorithm certified by Definition~\ref{def:fine-grained-reduction}.  More
  explicitly, the witness for this reduction includes the following eliminable
  obligations:
  \[
    r(A,B)=1,\qquad q_1(A,B)=A'\cup B',
  \]
  \[
    N_Q(q_1(A,B))=2n,\qquad L_Q(q_1(A,B))=\ell+2,
  \]
  the iff
  \[
    q_1(A,B)\text{ is a yes instance of monochromatic }\OV
    \quad\Longleftrightarrow\quad
    (A,B)\text{ is a yes instance of bichromatic }\OV,
  \]
  and the answer-transfer statement that every correct oracle answer on
  $q_1(A,B)$ is the correct final answer for $(A,B)$ because final
  postprocessing returns that same answer.
\end{lemma}
\begin{proof}
  \uses{def:bi-to-mono-ov-reduction-datum, lem:bi-to-mono-ov-query-equivalence, lem:bi-to-mono-ov-oracle-answer-transfer, def:fine-grained-reduction}
  Lemma~\ref{lem:bi-to-mono-ov-query-equivalence} gives the central
  correctness equivalence: the constructed monochromatic query has an
  orthogonal pair exactly when the original bichromatic instance has a cross
  orthogonal pair.  Lemma~\ref{lem:bi-to-mono-ov-oracle-answer-transfer}
  turns that equivalence into the required oracle-answer correctness
  obligation for the single query and the final postprocessing step.  The size
  changes
  from $n$ to $2n$, the dimension changes from $\ell$ to $\ell+2$, and the
  encoding construction is linear, so the subquadratic exponent and the
  polynomial-in-$\ell$ factor are preserved.  The final algorithm is the
  composed program $R_{\rm bi\to mono}[A_{\rm mono}]$ named in
  Definition~\ref{def:fine-grained-reduction}, not a separately chosen
  existential solver.
\end{proof}

\begin{lemma}[Monochromatic and bichromatic OV are equivalent, Lemma A.4]
  \label{lem:mono-bi-ov-equiv}
  \uses{def:ov, def:bichromatic-ov, def:fine-grained-reduction, lem:mono-to-bi-ov-fine-grained-reduction, lem:bi-to-mono-ov-fine-grained-reduction}
  With fixed quadratic-baseline witnesses for the monochromatic and
  bichromatic OV problem specifications, there are fine-grained reductions in
  both directions:
  \[
    \operatorname{def\_fine\_grained\_reduction}
      (\OV,\operatorname{BiOV},h_{\rm mono},h_{\rm bi})
  \]
  and
  \[
    \operatorname{def\_fine\_grained\_reduction}
      (\operatorname{BiOV},\OV,h_{\rm bi},h_{\rm mono}).
  \]
  Consequently a truly subquadratic algorithm for $\OV_{n,\ell}$ exists if and
  only if a truly subquadratic algorithm for bichromatic $\OV_{n,\ell}$ exists,
  up to polynomial factors in $\ell$.
\end{lemma}
\begin{proof}
  \uses{lem:mono-to-bi-ov-fine-grained-reduction, lem:bi-to-mono-ov-fine-grained-reduction, def:fine-grained-reduction}
  Lemma~\ref{lem:mono-to-bi-ov-fine-grained-reduction} supplies the reduction
  from monochromatic OV to bichromatic OV by the concrete composed algorithm
  $R_{\rm mono\to bi}[A_{\rm bi}]$.  Therefore any truly subquadratic
  bichromatic solver transfers to a truly subquadratic monochromatic solver
  with only polynomial overhead in $\ell$.  Lemma~\ref{lem:bi-to-mono-ov-fine-grained-reduction}
  supplies the reverse reduction by the concrete composed algorithm
  $R_{\rm bi\to mono}[A_{\rm mono}]$, so any truly subquadratic monochromatic
  solver transfers to a truly subquadratic bichromatic solver with the same
  allowed overhead.  These are exactly the two implications in the displayed
  equivalence.
\end{proof}

\begin{definition}[Minimum Inner Product, Definition 2.7/A.5]
  \label{def:min-ip}
  Given binary vectors $v_1,\ldots,v_n\in\{0,1\}^{\ell}$,
  $\MinIP_{n,\ell}$ asks to find a pair $i\neq j$ minimizing $\ip{v_i}{v_j}$.
\end{definition}

\begin{definition}[Approximate Minimum Inner Product, Definition 2.8/A.6]
  \label{def:gamma-min-ip}
  \uses{def:min-ip}
  For $\gamma\geq 1$, $\gamma$-$\MinIP_{n,\ell}$ asks to find a pair whose inner
  product is a $\gamma$-approximation of the minimum inner product.
\end{definition}

\begin{definition}[Minimum Inner Product decision version, Definition 2.9/A.7]
  \label{def:min-ip-decision}
  \uses{def:min-ip}
  Given a threshold $0\leq t\leq\ell$, $\MinIP_{n,\ell,t}$ asks whether some
  pair $i\neq j$ has $\ip{v_i}{v_j}\leq t$.
\end{definition}

\begin{definition}[Bichromatic Minimum Inner Product]
  \label{def:bichromatic-min-ip}
  \uses{def:min-ip}
  A rectangular bichromatic Min-IP input consists of two nonempty finite
  indexed lists
  \[
    A:\{1,\ldots,n_A\}\to\{0,1\}^{\ell},\qquad
    B:\{1,\ldots,n_B\}\to\{0,1\}^{\ell}.
  \]
  Its principal size is $N=\max\{n_A,n_B\}$ and its auxiliary dimension is
  $\ell$.  The exact optimization version asks for a pair
  $(i,j)\in\{1,\ldots,n_A\}\times\{1,\ldots,n_B\}$ minimizing
  $\ip{A_i}{B_j}$.  The $\gamma$-approximate version, for $\gamma\geq1$, asks
  for a pair $(i^*,j^*)$ satisfying
  \[
    \min_{i,j}\ip{A_i}{B_j}
    \leq \ip{A_{i^*}}{B_{j^*}}
    \leq \gamma\min_{i,j}\ip{A_i}{B_j}.
  \]
  The threshold version with $0\leq t\leq\ell$ asks for the Boolean answer
  \[
    \exists i,j,\quad \ip{A_i}{B_j}\leq t.
  \]
  The equal-size paper notation is the special case $n_A=n_B=n$.  The
  divide-and-recurse reductions below use the rectangular formulation so that
  canonical balanced splits of an arbitrary input length can be queried
  without adding artificial vectors that might change the minimum.
\end{definition}

\begin{definition}[Concrete OV-from-Min-IP reduction datum]
  \label{def:ov-from-minip-reduction-datum}
  \uses{def:ov, def:min-ip, def:gamma-min-ip, def:problem-algorithm}
  Fix $\gamma\geq1$.  The reduction datum is not a single ambiguous
  ``Min-IP solver'' interface.  It has a source branch
  \[
    \operatorname{source}\in
      \operatorname{ExactSource}(A_{\rm exact})
      \;\sqcup\;
      \operatorname{GammaSource}(A_{\gamma}),
  \]
  where $A_{\rm exact}$ is a correct exact $\MinIP_{n,\ell}$ program and
  $A_{\gamma}$ is a correct $\gamma$-$\MinIP_{n,\ell}$ program.  The two
  source branches have separate returned-pair selectors and separate halting
  facts.

  Lean-facing, this source is a genuine disjoint sum with branch-dependent
  payload.  An element of
  $\operatorname{ExactSource}(A_{\rm exact})$ contains only the exact
  $\MinIP$ program, its exact correctness/halting data, and the exact output
  selector below.  It contains no $\gamma$-$\MinIP$ program, no approximate
  correctness proof, no approximate returned-pair selector, and no
  $\gamma$-branch running-time bound.  An element of
  $\operatorname{GammaSource}(A_{\gamma})$ contains only the
  $\gamma$-$\MinIP$ program, its approximate correctness/halting data, and the
  approximate output selector below.  It contains no exact $\MinIP$ program, no
  exact correctness proof, no exact returned-pair selector, and no exact-branch
  running-time bound.  Later declarations quantify over one value
  $\operatorname{source}$ of this disjoint sum and eliminate it by cases; they
  do not take a pair of exact-and-approximate algorithms or a pair of
  exact-and-approximate runtime hypotheses.

  In the exact branch, for every input
  $V=(v_1,\ldots,v_n)\subseteq\{0,1\}^{\ell}$ with $n\geq2$, the selector
  \[
    \operatorname{ExactOutPair}_R(V)=(i^*_{\rm ex},j^*_{\rm ex})
  \]
  is exactly the halted output pair of $A_{\rm exact}$ on the exact Min-IP
  input $V$, at that program's first halting time.  Its correctness relation is
  the exact equality
  \[
    \ip{v_{i^*_{\rm ex}}}{v_{j^*_{\rm ex}}}
      =\min_{i\neq j}\ip{v_i}{v_j}.
  \]
  In the $\gamma$ branch, for every such $V$, the selector
  \[
    \operatorname{GammaOutPair}_R(V)=(i^*_{\gamma},j^*_{\gamma})
  \]
  is exactly the halted output pair of $A_{\gamma}$ on the
  $\gamma$-$\MinIP$ input $(V,\gamma)$, at that program's first halting time.
  Its correctness relation is the approximation inequality
  \[
    \min_{i\neq j}\ip{v_i}{v_j}
      \leq
    \ip{v_{i^*_{\gamma}}}{v_{j^*_{\gamma}}}
      \leq
    \gamma\min_{i\neq j}\ip{v_i}{v_j}.
  \]
  The two selectors are branch-specific data.  A formalization must not use
  $\operatorname{ExactOutPair}_R(V)$, a field named
  $\operatorname{exactPair}$, or any exact-output witness as the halted output
  of the $\gamma$-$\MinIP$ program.  Conversely,
  $\operatorname{GammaOutPair}_R(V)$ is not available in the exact branch.

  The induced OV decision program
  $R_{\MinIP\to\OV}[A_{\min}]$ has the following concrete operational phases.
  On input $V$, if $n<2$ it halts immediately and outputs $0$.  If
  $n\geq2$ and the source is exact, it runs $A_{\rm exact}$ on $V$ until the
  first halting output $\operatorname{ExactOutPair}_R(V)$ is produced.  If
  $n\geq2$ and the source is $\gamma$-approximate, it runs $A_{\gamma}$ on
  $(V,\gamma)$ until the first halting output
  $\operatorname{GammaOutPair}_R(V)$ is produced.  In either branch it then
  performs the deterministic coordinate scan of the branch's own returned pair:
  \[
    \beta_{\rm ex}(V)
      =\sum_{r=1}^{\ell}v_{i^*_{\rm ex}}[r]v_{j^*_{\rm ex}}[r],
    \qquad
    \beta_{\gamma}(V)
      =\sum_{r=1}^{\ell}v_{i^*_{\gamma}}[r]v_{j^*_{\gamma}}[r].
  \]
  The scan step count is the concrete value
  $\operatorname{LocalScanSteps}_R(V)$ and satisfies the displayed datum
  inequality
  \[
    \operatorname{LocalScanSteps}_R(V)\leq C_0(1+\ell)
  \]
  for an absolute constant $C_0$ supplied by the datum.  The branch/output
  overhead is bounded by an absolute constant $C_1$.  The final Boolean output
  is $1$ exactly when the scanned value for the active branch is $0$, and is
  $0$ otherwise.

  Thus the only subroutine call is the actual run of the active branch's
  program on the original Min-IP input, the returned pair is the active
  branch's own first halted output, and the final Boolean is the equality test
  on the concrete inner product of that same pair.  There is no externally
  supplied OV answer, no abstract correctness-equivalence field, and no
  shared output-pair field that can be used for both exact and approximate
  sources.
\end{definition}

\begin{lemma}[Min-IP optimum is zero exactly on OV yes instances]
  \label{lem:minip-optimum-zero-iff-ov}
  \uses{def:ov, def:min-ip}
  For every binary-vector input $V=(v_1,\ldots,v_n)$ with $n\geq2$,
  \[
    \min_{i\neq j}\ip{v_i}{v_j}=0
    \quad\Longleftrightarrow\quad
    V\text{ is a yes instance of }\OV_{n,\ell}.
  \]
\end{lemma}
\begin{proof}
  \uses{def:ov, def:min-ip}
  Inner products of binary vectors are sums of nonnegative terms, so every
  feasible pair has inner product at least $0$.  If an OV witness
  $i\neq j$ exists, then the minimum is at most $\ip{v_i}{v_j}=0$, hence it is
  exactly $0$.  Conversely, if the minimum is $0$, any pair attaining that
  minimum is a distinct orthogonal pair and therefore an OV witness.
\end{proof}

\begin{lemma}[Exact Min-IP output gives correct OV answer]
  \label{lem:exact-minip-output-ov-correct}
  \uses{def:ov-from-minip-reduction-datum, lem:minip-optimum-zero-iff-ov}
  Suppose $n\geq2$ and the exact $\MinIP$ algorithm in
  Definition~\ref{def:ov-from-minip-reduction-datum} halts on $V$ with a pair
  $(i^*,j^*)$ satisfying
  \[
    \ip{v_{i^*}}{v_{j^*}}
      =\min_{i\neq j}\ip{v_i}{v_j}.
  \]
  Then the final test
  $\ip{v_{i^*}}{v_{j^*}}=0$ outputs $1$ exactly on OV yes instances and
  outputs $0$ exactly on OV no instances.
\end{lemma}
\begin{proof}
  \uses{def:ov-from-minip-reduction-datum, lem:minip-optimum-zero-iff-ov}
  By the exact-output hypothesis, the tested value is the global minimum inner
  product over distinct pairs.  Lemma~\ref{lem:minip-optimum-zero-iff-ov}
  identifies the event that this minimum is $0$ with the OV yes condition.
  The final branch of Definition~\ref{def:ov-from-minip-reduction-datum}
  returns precisely the Boolean value of that equality test.
\end{proof}

\begin{lemma}[Approximate Min-IP output gives correct OV answer]
  \label{lem:approx-minip-output-ov-correct}
  \uses{def:ov-from-minip-reduction-datum, lem:minip-optimum-zero-iff-ov, def:gamma-min-ip}
  Suppose $n\geq2$, $\gamma\geq1$, and the $\gamma$-$\MinIP$ algorithm in
  Definition~\ref{def:ov-from-minip-reduction-datum} halts on $V$ with a pair
  $(i^*,j^*)$ satisfying
  \[
    \min_{i\neq j}\ip{v_i}{v_j}
      \leq \ip{v_{i^*}}{v_{j^*}}
      \leq \gamma\min_{i\neq j}\ip{v_i}{v_j}.
  \]
  Then the final test
  $\ip{v_{i^*}}{v_{j^*}}=0$ outputs $1$ exactly on OV yes instances and
  outputs $0$ exactly on OV no instances.
\end{lemma}
\begin{proof}
  \uses{def:ov-from-minip-reduction-datum, lem:minip-optimum-zero-iff-ov, def:gamma-min-ip}
  If $V$ is an OV yes instance, Lemma~\ref{lem:minip-optimum-zero-iff-ov} gives
  that the minimum is $0$.  The approximation inequalities become
  $0\leq\ip{v_{i^*}}{v_{j^*}}\leq0$, so the tested inner product is $0$ and the
  final output is $1$.  If $V$ is an OV no instance, the same lemma says the
  minimum is not $0$; since all binary inner products are nonnegative integers,
  the minimum is at least $1$, and the lower approximation inequality gives
  $\ip{v_{i^*}}{v_{j^*}}\geq1$.  The equality test is then false and the final
  output is $0$.
\end{proof}

\begin{lemma}[Exact Min-IP detects OV: branch correctness certificate]
  \label{lem:exact-minip-detects-ov-correctness-certificate}
  \uses{def:ov-from-minip-reduction-datum, lem:exact-minip-output-ov-correct}
  In the exact source case
  $\operatorname{source}=\operatorname{ExactSource}(A_{\rm exact})$, the
  concrete program $R_{\MinIP\to\OV}[A_{\rm exact}]$ has a correctness
  certificate
  \[
    \operatorname{ExactMinIPDetectsOVCorrect}_R(A_{\rm exact})
  \]
  with exactly the following content.  If $n<2$, the program makes no Min-IP
  call and outputs $0$.  If $n\geq2$, it runs only the exact program
  $A_{\rm exact}$ on the exact $\MinIP$ input $V$ until the first halted output
  $\operatorname{ExactOutPair}_R(V)$, scans that same pair, and outputs $1$
  exactly when the scanned inner product is $0$:
  \[
    R_{\MinIP\to\OV}[A_{\rm exact}](V)=1
      \quad\Longleftrightarrow\quad
    \exists i\neq j,\ \ip{v_i}{v_j}=0.
  \]
  It outputs $0$ exactly in the negation of this condition, and its local
  first-halting-time equation is
  \[
    s_{\OV}(V)=s_{\rm exact}(V)
      +\operatorname{LocalScanSteps}_R(V)+C_1.
  \]

  This node has no parameters or hypotheses for a $\gamma$-$\MinIP$ program,
  no $\gamma$-branch returned pair, and no runtime bound for either branch.
  It is a correctness and local first-halting certificate for the exact branch
  alone.
\end{lemma}
\begin{proof}
  \uses{def:ov-from-minip-reduction-datum, lem:exact-minip-output-ov-correct}
  Unfold the exact branch of
  Definition~\ref{def:ov-from-minip-reduction-datum}.  The small-input branch
  is correct because no distinct pair exists.  For $n\geq2$, the scanned pair
  is exactly $\operatorname{ExactOutPair}_R(V)$, the first halted output of the
  exact Min-IP program.  Lemma~\ref{lem:exact-minip-output-ov-correct} proves
  that testing this pair for inner product zero gives precisely the OV yes/no
  answer.  The local first-halting-time equation is the exact-branch phase
  decomposition from the same reduction datum.
\end{proof}

\begin{lemma}[$\gamma$-Min-IP detects OV: branch correctness certificate]
  \label{lem:gamma-minip-detects-ov-correctness-certificate}
  \uses{def:ov-from-minip-reduction-datum, lem:approx-minip-output-ov-correct}
  In the approximate source case
  $\operatorname{source}=\operatorname{GammaSource}(A_{\gamma})$, the concrete
  program $R_{\MinIP\to\OV}[A_{\gamma}]$ has a correctness certificate
  \[
    \operatorname{GammaMinIPDetectsOVCorrect}_R(A_{\gamma})
  \]
  with exactly the following content.  If $n<2$, the program makes no Min-IP
  call and outputs $0$.  If $n\geq2$, it runs only the $\gamma$-$\MinIP$
  program $A_{\gamma}$ on the approximate input $(V,\gamma)$ until the first
  halted output $\operatorname{GammaOutPair}_R(V)$, scans that same approximate
  output pair, and outputs $1$ exactly when the scanned inner product is $0$:
  \[
    R_{\MinIP\to\OV}[A_{\gamma}](V)=1
      \quad\Longleftrightarrow\quad
    \exists i\neq j,\ \ip{v_i}{v_j}=0.
  \]
  It outputs $0$ exactly in the negation of this condition, and its local
  first-halting-time equation is
  \[
    s_{\OV}(V)=s_{\gamma}(V,\gamma)
      +\operatorname{LocalScanSteps}_R(V)+C_1.
  \]

  This node has no parameters or hypotheses for an exact $\MinIP$ program, no
  exact-branch returned pair, and no runtime bound for either branch.  It is a
  correctness and local first-halting certificate for the $\gamma$ branch
  alone.  Every occurrence of the returned pair in this statement is
  $\operatorname{GammaOutPair}_R(V)$.
\end{lemma}
\begin{proof}
  \uses{def:ov-from-minip-reduction-datum, lem:approx-minip-output-ov-correct}
  Unfold the approximate branch of
  Definition~\ref{def:ov-from-minip-reduction-datum}.  The small-input branch
  is correct because no distinct pair exists.  For $n\geq2$, the scanned pair
  is exactly $\operatorname{GammaOutPair}_R(V)$, the first halted output of the
  $\gamma$-$\MinIP$ program.  Lemma~\ref{lem:approx-minip-output-ov-correct}
  proves that testing this pair for inner product zero gives precisely the OV
  yes/no answer.  The local first-halting-time equation is the $\gamma$-branch
  phase decomposition from the same reduction datum.
\end{proof}

\begin{lemma}[OV-from-Min-IP running time]
  \label{lem:ov-from-minip-running-time}
  \uses{def:ov-from-minip-reduction-datum, def:exact-worst-case-running-time, def:truly-subquadratic}
  Let $R$ be the concrete reduction datum of
  Definition~\ref{def:ov-from-minip-reduction-datum}.  The running-time
  statement is a branch-by-branch theorem about the induced OV program
  $R_{\MinIP\to\OV}[A_{\min}]$, not an assumption about an externally supplied
  $T_{\OV}$.

  In the exact branch, suppose $A_{\rm exact}$ has semantic worst-case bound
  $T_{\min}$ in the sense of
  Definition~\ref{def:exact-worst-case-running-time}.  Then for every valid OV
  input $V$ with $n\geq2$, the first halting time of the induced OV program on
  $V$ is exactly
  \[
    s_{\OV}(V)
      =
    s_{\rm exact}(V)
      +\operatorname{LocalScanSteps}_R(V)+C_1,
  \]
  where $s_{\rm exact}(V)$ is the first halting time of $A_{\rm exact}$ on the
  exact Min-IP input $V$, and its halted output is
  $\operatorname{ExactOutPair}_R(V)$.  Therefore
  \[
    s_{\OV}(V)
      \leq
    T_{\min}(n,\ell)+C_0(1+\ell)+C_1 .
  \]

  In the $\gamma$ branch, suppose $A_{\gamma}$ has semantic worst-case bound
  $T_{\min}$ on $\gamma$-$\MinIP$ inputs.  Then for every valid OV input $V$
  with $n\geq2$, the first halting time of the induced OV program on $V$ is
  exactly
  \[
    s_{\OV}(V)
      =
    s_{\gamma}(V,\gamma)
      +\operatorname{LocalScanSteps}_R(V)+C_1,
  \]
  where $s_{\gamma}(V,\gamma)$ is the first halting time of $A_{\gamma}$ on
  the approximate input $(V,\gamma)$, and its halted output is
  $\operatorname{GammaOutPair}_R(V)$.  Therefore
  \[
    s_{\OV}(V)
      \leq
    T_{\min}(n,\ell)+C_0(1+\ell)+C_1 .
  \]

  In the small-input branch $n<2$, the induced OV program makes no Min-IP
  call and halts within the same absolute overhead constant $C_1$.  Hence the
  total worst-case bound for the constructed OV program is the concrete
  function
  \[
    T_{\OV}^R(n,\ell)=T_{\min}(n,\ell)+C_0(1+\ell)+C_1
  \]
  for $n\geq2$, with the constant branch bound for $n<2$.  If
  \[
    T_{\min}(n,\ell)\leq C n^{2-\eps}(1+\ell)^b
  \]
  for some $\eps>0$, then this displayed constructed bound
  $T_{\OV}^R$ is truly subquadratic in $n$ up to polynomial factors in $\ell$,
  after decreasing the exponent and increasing the polynomial factor if
  necessary.

  Lean-facing, this lemma must expose the two branch conclusions above.  It
  must derive the $T_{\OV}^R$ bound from the first-halting time of the active
  Min-IP program plus $\operatorname{LocalScanSteps}_R(V)$ and $C_1$.  It
  must not take as hypotheses a predicate named
  $\operatorname{runtime\_accounting}$, an assumed inequality for an arbitrary
  $T_{\OV}$, or a small-input bound for an externally supplied OV solver.
  In the $\gamma$ branch, every halting fact and every scan equation must use
  $\operatorname{GammaOutPair}_R(V)$, not
  $\operatorname{ExactOutPair}_R(V)$ or any field named
  $\operatorname{exactPair}$.
\end{lemma}
\begin{proof}
  \uses{def:ov-from-minip-reduction-datum, def:exact-worst-case-running-time, def:truly-subquadratic}
  Unfold Definition~\ref{def:ov-from-minip-reduction-datum}.  In the exact
  branch, the constructed program first executes the exact Min-IP program on
  $V$ until its first halted output
  $\operatorname{ExactOutPair}_R(V)$ appears, then performs exactly the
  displayed coordinate scan and final Boolean branch.  In the $\gamma$ branch,
  the same phase decomposition uses the first halted output
  $\operatorname{GammaOutPair}_R(V)$ of the $\gamma$-$\MinIP$ program on
  $(V,\gamma)$.  The phase lengths are disjoint, so their sum gives the two
  displayed first-halting-time equations.  The semantic worst-case bound
  controls the active Min-IP phase, and the datum inequality
  $\operatorname{LocalScanSteps}_R(V)\leq C_0(1+\ell)$ controls the scan phase.
  The $n<2$ branch is the direct constant-time output branch.  The added
  linear-in-$\ell$ and constant terms are absorbed into the allowed
  polynomial-in-$\ell$ factor after reducing the subquadratic exponent, for
  example from $\eps$ to $\eps/2$.
\end{proof}

\begin{lemma}[OV-from-Min-IP inherited running-time certificate]
  \label{lem:ov-from-minip-inherited-subquadratic-certificate}
  \uses{def:ov-from-minip-reduction-datum, lem:ov-from-minip-running-time, def:problem-algorithm, def:exact-worst-case-running-time, def:truly-subquadratic}
  Fix $\gamma\geq1$ and one active source branch of
  Definition~\ref{def:ov-from-minip-reduction-datum}.  Suppose the active
  exact or $\gamma$-approximate Min-IP program is a correct algorithm with
  semantic worst-case bound $T_{\min}$, and suppose there are constants
  $\eps>0$, $C>0$, $b\in\N$, and $n_0$ such that
  \[
    T_{\min}(n,\ell)\leq C n^{2-\eps}(1+\ell)^b
  \]
  for every $n\geq n_0$ and every allowed dimension $\ell$.

  Then the concrete induced OV program
  \[
    A_{\OV}^R=R_{\MinIP\to\OV}[A_{\min}]
  \]
  from Definition~\ref{def:ov-from-minip-reduction-datum} has the explicit
  inherited running-time certificate
  \[
    \operatorname{InheritedOVTimeCert}_R(A_{\min},T_{\min})
  \]
  consisting of exactly the following data and facts:
  \[
    T_{\OV}^R(n,\ell)=
    \begin{cases}
      C_1, & n<2,\\
      T_{\min}(n,\ell)+C_0(1+\ell)+C_1, & n\geq2,
    \end{cases}
  \]
  a proof that $T_{\OV}^R$ is a semantic worst-case running-time bound for the
  actual first-halting execution of the same program $A_{\OV}^R$, and
  constants
  \[
    \eps'=\eps/2>0,\qquad b'=\max\{b,1\},
  \]
  together with a constant $C'>0$ and threshold $n_0'$ such that
  \[
    T_{\OV}^R(n,\ell)\leq C'n^{2-\eps'}(1+\ell)^{b'}
  \]
  for every $n\geq n_0'$ and every allowed dimension $\ell$.

  This certificate is part of the conclusion, not a hypothesis.  The
  worst-case-bound proof is derived from the branch-specific first-halting-time
  equations of Lemma~\ref{lem:ov-from-minip-running-time}; the subquadratic
  inequality is derived from the displayed bound on $T_{\min}$.  The same
  constructed program $A_{\OV}^R$ appears in the correctness certificate, the
  worst-case-bound certificate, and the truly subquadratic certificate.
  Lean-facing, a declaration for this node must return the tuple
  \[
    (T_{\OV}^R,\ h_{\rm worst},\ \eps',\ C',\ b',\ n_0',
      h_{\rm subquad})
  \]
  with the displayed meanings.  It must not assume
  $\operatorname{def\_truly\_subquadratic}(T_{\OV})$ for an arbitrary
  externally supplied $T_{\OV}$, and it must not prove only branch-local
  inequalities without producing the global worst-case bound $T_{\OV}^R$ for
  $A_{\OV}^R$.
\end{lemma}
\begin{proof}
  \uses{lem:ov-from-minip-running-time, def:exact-worst-case-running-time, def:truly-subquadratic}
  Define $T_{\OV}^R$ by the displayed small-input and $n\geq2$ cases.  In each
  source branch, Lemma~\ref{lem:ov-from-minip-running-time} identifies the
  first halting time of the induced OV program with the active Min-IP
  first-halting time plus the scan and constant overhead, and bounds that sum
  by $T_{\min}(n,\ell)+C_0(1+\ell)+C_1$.  This proves that $T_{\OV}^R$ is a
  worst-case bound for the actual induced program.  Substituting the assumed
  bound on $T_{\min}$ gives
  $T_{\OV}^R(n,\ell)\leq C n^{2-\eps}(1+\ell)^b+C_0(1+\ell)+C_1$ for $n\geq2$.
  With $\eps'=\eps/2$ and $b'=\max\{b,1\}$, the linear and constant overheads
  are absorbed into $C'n^{2-\eps'}(1+\ell)^{b'}$ after increasing $C'$ and
  $n_0'$.  The $n<2$ branch is covered by the constant case of
  $T_{\OV}^R$.
\end{proof}

\begin{lemma}[Minimum inner product detects OV: correctness certificate]
  \label{lem:min-ip-detects-ov-correctness-certificate}
  \uses{def:ov-from-minip-reduction-datum, lem:exact-minip-detects-ov-correctness-certificate, lem:gamma-minip-detects-ov-correctness-certificate}
  For every $\gamma\geq1$ and every active source branch
  \[
    \operatorname{source}\in
      \operatorname{ExactSource}(A_{\rm exact})
      \sqcup
      \operatorname{GammaSource}(A_{\gamma})
  \]
  from Definition~\ref{def:ov-from-minip-reduction-datum}, the concrete
  run-and-test program for that one active branch returns the OV answer on
  every valid OV input.  This node contains only the correctness and local
  first-halting behavior of the constructed program.  It has no worst-case
  running-time hypotheses and no inactive-branch hypotheses.

  More precisely, it returns the certificate
  \[
    \operatorname{MinIPDetectsOVCorrect}_R(\operatorname{source})
  \]
  defined by case analysis on the displayed disjoint sum:
  \[
    \operatorname{MinIPDetectsOVCorrect}_R
      (\operatorname{ExactSource}(A_{\rm exact}))
    =
    \operatorname{ExactMinIPDetectsOVCorrect}_R(A_{\rm exact}),
  \]
  \[
    \operatorname{MinIPDetectsOVCorrect}_R
      (\operatorname{GammaSource}(A_{\gamma}))
    =
    \operatorname{GammaMinIPDetectsOVCorrect}_R(A_{\gamma}).
  \]
  Eliminating the exact branch yields exactly
  Lemma~\ref{lem:exact-minip-detects-ov-correctness-certificate}; eliminating
  the $\gamma$ branch yields exactly
  Lemma~\ref{lem:gamma-minip-detects-ov-correctness-certificate}.  A Lean
  statement for this label must therefore have one active source argument and
  must not require hypotheses named like $h_{\rm Exact}$ and $h_{\rm Gamma}$
  simultaneously.  In particular, it must not ask for a semantic
  worst-case-running-time bound for the inactive exact branch in the
  $\gamma$ case, nor for a semantic worst-case-running-time bound for the
  inactive $\gamma$ branch in the exact case.
\end{lemma}
\begin{proof}
  \uses{def:ov-from-minip-reduction-datum, lem:exact-minip-detects-ov-correctness-certificate, lem:gamma-minip-detects-ov-correctness-certificate}
  Cases on the disjoint-sum source.  In the exact case, return
  Lemma~\ref{lem:exact-minip-detects-ov-correctness-certificate}.  In the
  $\gamma$ case, return
  Lemma~\ref{lem:gamma-minip-detects-ov-correctness-certificate}.  Since the
  case split exposes only the active branch payload, no inactive-branch
  program, correctness proof, or running-time bound is available or needed.
\end{proof}

\begin{lemma}[Minimum inner product detects OV: inherited runtime component]
  \label{lem:min-ip-detects-ov-inherited-runtime-component}
  \uses{def:ov-from-minip-reduction-datum, lem:ov-from-minip-inherited-subquadratic-certificate}
  Fix $\gamma\geq1$ and one active source branch
  \[
    \operatorname{source}\in
      \operatorname{ExactSource}(A_{\rm exact})
      \sqcup
      \operatorname{GammaSource}(A_{\gamma})
  \]
  from Definition~\ref{def:ov-from-minip-reduction-datum}.  Let
  $A_{\min}(\operatorname{source})$ be the exact program in the exact case and
  the $\gamma$-$\MinIP$ program in the $\gamma$ case.  Let
  $T_{\min}:\N\times\N\to\R_{\geq0}$ be a semantic worst-case bound for that
  active Min-IP program.  Suppose
  there are constants $\eps>0$, $C>0$, $b\in\N$, and $n_0$ such that
  \[
    T_{\min}(n,\ell)\leq C n^{2-\eps}(1+\ell)^b
  \]
  for every $n\geq n_0$ and every allowed dimension $\ell$.

  Then the constructed OV program
  $R_{\MinIP\to\OV}[A_{\min}(\operatorname{source})]$ has the inherited
  runtime component
  \[
    \operatorname{MinIPDetectsOVRuntime}_R
      (\operatorname{source},T_{\min},C,\eps,b,n_0),
  \]
  which is exactly the certificate supplied by
  Lemma~\ref{lem:ov-from-minip-inherited-subquadratic-certificate}.  Eliminating
  it gives the concrete bound
  \[
    T_{\OV}^R(n,\ell)=
    \begin{cases}
      C_1, & n<2,\\
      T_{\min}(n,\ell)+C_0(1+\ell)+C_1, & n\geq2,
    \end{cases}
  \]
  a proof that $T_{\OV}^R$ bounds the actual first-halting time of the same
  constructed OV program, and constants
  $\eps'>0$, $C'>0$, $b'\in\N$, and $n_0'$ with
  \[
    \forall n\geq n_0',\ \forall\ell,\qquad
      T_{\OV}^R(n,\ell)\leq C'n^{2-\eps'}(1+\ell)^{b'}.
  \]
  This node has the explicit parameters $T_{\min}$, $C$, $\eps$, $b$, and
  $n_0$ for the active source only; the runtime component is not derivable
  from branch-local accounting alone and is not an assumed property of an
  arbitrary OV bound.  It must not require a runtime bound for the inactive
  source branch.
\end{lemma}
\begin{proof}
  \uses{lem:ov-from-minip-inherited-subquadratic-certificate}
  Apply Lemma~\ref{lem:ov-from-minip-inherited-subquadratic-certificate} to the
  active source branch, its semantic bound $T_{\min}$, and the displayed
  constants.  The resulting tuple is precisely the stated
  $\operatorname{MinIPDetectsOVRuntime}_R$ component.
\end{proof}

\begin{lemma}[Minimum inner product detects OV]
  \label{lem:min-ip-detects-ov}
  \uses{lem:min-ip-detects-ov-correctness-certificate, lem:min-ip-detects-ov-inherited-runtime-component}
  Fix $\gamma\geq1$, one active source branch
  \[
    \operatorname{source}\in
      \operatorname{ExactSource}(A_{\rm exact})
      \sqcup
      \operatorname{GammaSource}(A_{\gamma})
  \]
  from Definition~\ref{def:ov-from-minip-reduction-datum}.  Let
  $A_{\min}(\operatorname{source})$ denote the program carried by that branch:
  it is $A_{\rm exact}$ in the exact case and $A_{\gamma}$ in the $\gamma$
  case.  Let $T_{\min}$ be a semantic worst-case bound for this same active
  program and for no other source branch.  Suppose
  there are constants $\eps>0$, $C>0$, $b\in\N$, and $n_0$ such that
  \[
    T_{\min}(n,\ell)\leq C n^{2-\eps}(1+\ell)^b
  \]
  for every $n\geq n_0$ and every allowed dimension $\ell$.

  The concrete run-and-test program
  $R_{\MinIP\to\OV}[A_{\min}(\operatorname{source})]$ directly certifies an
  OV solver together with
  the inherited subquadratic running-time certificate.  Lean-facing, the
  conclusion of this node is exactly the ordered pair
  \[
    \bigl(
      \operatorname{MinIPDetectsOVCorrect}_R(\operatorname{source}),
      \operatorname{MinIPDetectsOVRuntime}_R
        (\operatorname{source},T_{\min},C,\eps,b,n_0)
    \bigr).
  \]
  The first component is the branch-specific correctness certificate of
  Lemma~\ref{lem:min-ip-detects-ov-correctness-certificate}: direct
  small-input behavior, active subroutine first-halting output, scanned active
  pair, OV yes/no correctness, and first-halting-time-plus-scan accounting.
  The second component is the inherited runtime component of
  Lemma~\ref{lem:min-ip-detects-ov-inherited-runtime-component}: the concrete
  piecewise bound $T_{\OV}^R$, the semantic worst-case proof for the actual
  constructed OV program, and the constants
  $\eps'>0$, $C'>0$, $b'\in\N$, and $n_0'$ proving
  \[
    T_{\OV}^R(n,\ell)\leq C'n^{2-\eps'}(1+\ell)^{b'}
  \]
  for all $n\geq n_0'$ and every allowed $\ell$.

  Thus the declaration for this label has exactly one active source parameter,
  one active semantic bound $T_{\min}$ for the program in that source, and
  hypotheses for $C$, $\eps$, $b$, and $n_0$ about that same bound.  In the
  exact case it may require the exact program's correctness and the exact
  bound; it must not require any $\gamma$-program correctness or runtime
  hypothesis.  In the $\gamma$ case it may require the $\gamma$ program's
  correctness and the $\gamma$ bound; it must not require any exact-program
  correctness or runtime hypothesis.  A statement with simultaneous
  hypotheses for both branches, such as both $h_{\rm Exact}$ and
  $h_{\rm Gamma}$ or both $T_{\rm exact}$ and $T_{\gamma}$, is strictly
  stronger than this node and is not a valid formalization of this label.
\end{lemma}
\begin{proof}
  \uses{lem:min-ip-detects-ov-correctness-certificate, lem:min-ip-detects-ov-inherited-runtime-component}
  The first component is
  Lemma~\ref{lem:min-ip-detects-ov-correctness-certificate}.  The second
  component is
  Lemma~\ref{lem:min-ip-detects-ov-inherited-runtime-component} applied to the
  same active source branch, the displayed semantic bound $T_{\min}$, and the
  constants $C,\eps,b,n_0$.  Pairing those two components gives exactly the
  stated certificate for the same constructed program
  $R_{\MinIP\to\OV}[A_{\min}(\operatorname{source})]$.  The proof is by the
  same case split on $\operatorname{source}$ in both components, so the exact
  branch never mentions approximate-source data and the approximate branch
  never mentions exact-source data.
\end{proof}

\begin{definition}[Balanced recursive split tree for Min-IP]
  \label{def:minip-balanced-recursive-split-tree}
  \uses{def:min-ip}
  For a monochromatic input
  $V:\{1,\ldots,n\}\to\{0,1\}^{\ell}$, the balanced split tree is the
  following concrete finite list of intervals; it is not an arbitrary
  tree-like structure.  An interval is encoded by endpoints $(L,R)$ with
  $1\leq L\leq R\leq n$, denoting $\{L,\ldots,R\}$.  For every such interval
  define
  \[
    M(L,R)=L+\lfloor(R-L)/2\rfloor .
  \]
  The recursive interval-list constructor
  $\operatorname{SplitIntervals}_n(L,R)$ is defined by primitive recursion on
  the length $R-L+1$.  Its Lean-facing defining equations are exactly the
  following equations, with no hidden tree parameter:
  \[
    \operatorname{SplitIntervals}_n(L,L)=[(L,L)],
  \]
  and, when $L<R$,
  \[
    \operatorname{SplitIntervals}_n(L,R)
      =
      [(L,R)]
      \mathbin{++}
      \operatorname{SplitIntervals}_n(L,M(L,R))
      \mathbin{++}
      \operatorname{SplitIntervals}_n(M(L,R)+1,R).
  \]
  Equivalently, if the recursion is implemented by well-founded recursion
  rather than by a primitive-recursive syntax form, the recursive calls in the
  defining body are precisely the two calls on the strict subintervals
  $(L,M(L,R))$ and $(M(L,R)+1,R)$, and the two displayed equations are
  projection lemmas for the declaration itself.
  The split tree for the input size $n$ is definitionally
  \[
    \mathcal T_n=
    \begin{cases}
      [\,], & n=0,\\
      \operatorname{SplitIntervals}_n(1,n), & n\geq1.
    \end{cases}
  \]
  It depends on $V$ only through the row values used by attached queries; the
  node set depends only on $n$.  The internal-node list is the filter
  \[
    \operatorname{InternalIntervals}_n
      =\{(L,R)\in\mathcal T_n : L<R\}
  \]
  with inherited list order, and the leaf list is the complementary filter
  $L=R$.  The child intervals of an internal node $(L,R)$ are exactly
  \[
    I_L=(L,M(L,R))
  \]
  and right child
  \[
    I_R=(M(L,R)+1,R).
  \]
  These two children are nonempty, disjoint, their union is $(L,R)$, and both
  children occur as the next recursively generated roots in $\mathcal T_n$.
  No other intervals are nodes: every node of $\mathcal T_n$ is produced by
  exactly one application of the displayed recursion from the root, and the
  list has no duplicate endpoint pairs.

  The parent relation, ancestor relation, depth, and internal-node predicate
  are derived from this concrete list and these child equations.  Consequently
  the following invariants are part of this definition, not later optional
  hypotheses: two generated intervals are either disjoint or one contains the
  other; every non-root node has a unique parent; every internal node has
  exactly the two displayed children; and every endpoint pair appears at most
  once in $\mathcal T_n$.

  The bichromatic query attached to an internal node $I=(L,R)$ is the
  rectangular instance
  \[
    Q_I(V)=\bigl((v_i)_{i\in I_L},(v_j)_{j\in I_R}\bigr)
  \]
  of Definition~\ref{def:bichromatic-min-ip}, with principal size
  $N_I=\max\{|I_L|,|I_R|\}\leq |I|$ and dimension $\ell$.
  The canonical query schedule is the finite list
  \[
    \operatorname{SplitQueries}_n(V)
      =
      [\,Q_I(V): I\in\operatorname{InternalIntervals}_n
        \text{ in inherited list order}\,].
  \]
  If $\operatorname{InternalIntervals}_n[s]=(L_s,R_s)$ and
  $M_s=M(L_s,R_s)$, then the $s$th query record is definitionally the
  rectangular bichromatic input
  \[
    \operatorname{SplitQueries}_n(V)[s]
      =
      \bigl(A_s,B_s\bigr),
  \]
  where
  \[
    A_s:\{L_s,\ldots,M_s\}\to\{0,1\}^{\ell},\qquad A_s(i)=v_i,
  \]
  and
  \[
    B_s:\{M_s+1,\ldots,R_s\}\to\{0,1\}^{\ell},\qquad B_s(j)=v_j.
  \]
  The principal size and dimension projections of this concrete query are
  \[
    N(\operatorname{SplitQueries}_n(V)[s])
      =\max\{M_s-L_s+1,R_s-M_s\},
    \qquad
    L(\operatorname{SplitQueries}_n(V)[s])=\ell.
  \]

  Lean-facing, the declaration for this node must define the interval list
  $\mathcal T_n$, the internal-node filtered list
  $\operatorname{InternalIntervals}_n$, the midpoint function $M$, the two
  child functions, the query function $Q_I(V)$, and the query-schedule list
  $\operatorname{SplitQueries}_n(V)$ by the displayed equations.
  The recursive constructor
  $\operatorname{SplitIntervals}_n(L,R)$ must be part of this declaration's
  exposed definitional data: callers can unfold it at $(L,L)$ and at
  $(L,R)$ with $L<R$ to obtain the two equations above.  It is not acceptable
  for this node to return only a wrapper around an unspecified helper such as
  $\operatorname{minipSplitIntervals}$ unless that helper is the same
  declaration's exposed recursive component and the base/step equations above
  are fields or projection lemmas of this node.
  It must not instead introduce a structure with fields such as
  $\operatorname{nodes}$, $\operatorname{children}$, $\operatorname{isTree}$,
  or $\operatorname{coversPairs}$ that can be instantiated by an arbitrary
  object.  It must also not define
  $\operatorname{InternalIntervals}_n$ or $\operatorname{SplitQueries}_n(V)$
  from an arbitrary list satisfying a tree predicate.  Extra nodes, duplicate
  endpoint pairs, missing recursive children, duplicate intervals with the
  same endpoints, and extra or missing query records are ruled out
  definitionally by the list recursion and the displayed map from internal
  intervals to attached rectangular queries.
\end{definition}

\begin{lemma}[Split tree covers each unordered pair once]
  \label{lem:minip-split-tree-pair-coverage}
  \uses{def:minip-balanced-recursive-split-tree}
  For every distinct pair of indices $p\neq q$ in $\{1,\ldots,n\}$, there is a
  unique endpoint pair $(L,R)$ in the concrete filtered list
  $\operatorname{InternalIntervals}_n$ of
  Definition~\ref{def:minip-balanced-recursive-split-tree} such that
  $L\leq p,q\leq R$ and the two indices lie in different displayed children:
  \[
    \bigl(p\leq M(L,R)<q\bigr)
    \quad\text{or}\quad
    \bigl(q\leq M(L,R)<p\bigr).
  \]
  Equivalently, the unordered pair $\{p,q\}$ appears as a cross-pair in
  exactly one attached bichromatic query $Q_{(L,R)}(V)$ from the canonical
  split-tree query schedule.  The uniqueness quantifier ranges only over the
  recursively generated list $\operatorname{InternalIntervals}_n$; there is no
  arbitrary split-tree value in the statement.

  The query-schedule form of the same uniqueness is also part of the lemma.
  For every monochromatic input
  $V:\{1,\ldots,n\}\to\{0,1\}^{\ell}$ and every distinct pair $p\neq q$, there
  is a unique query-list position
  \[
    s\in\{1,\ldots,
      \operatorname{length}(\operatorname{SplitQueries}_n(V))\}
  \]
  such that, writing
  $\operatorname{InternalIntervals}_n[s]=(L_s,R_s)$ and
  $M_s=M(L_s,R_s)$, the query at that position is exactly
  \[
    \operatorname{SplitQueries}_n(V)[s]
      =
      \bigl((v_i)_{L_s\leq i\leq M_s},
            (v_j)_{M_s<j\leq R_s}\bigr),
  \]
  and the pair $(p,q)$ is represented in this concrete query with one index in
  its left row list and one index in its right row list:
  \[
    (L_s\leq p\leq M_s\ \wedge\ M_s<q\leq R_s)
    \quad\text{or}\quad
    (L_s\leq q\leq M_s\ \wedge\ M_s<p\leq R_s).
  \]
  If a second query-list position $s'$ has the same property, then
  $s'=s$ and hence
  \[
    \operatorname{SplitQueries}_n(V)[s']
      =
      \operatorname{SplitQueries}_n(V)[s]
  \]
  as the same attached rectangular bichromatic query record, not merely as the
  same endpoint pair.  Conversely, if a query record in
  $\operatorname{SplitQueries}_n(V)$ contains $(p,q)$ as such a left-right
  cross-pair, then its position is this unique $s$.

  More explicitly, if $(L,R)$ and $(L',R')$ are both internal intervals from
  $\operatorname{InternalIntervals}_n$ satisfying the displayed separation
  predicate for the same pair $(p,q)$, then $L=L'$ and $R=R'$.  If
  $(L,R)$ is the unique interval, the pair occurs in the query
  \[
    Q_{(L,R)}(V)=
    \bigl((v_i)_{L\leq i\leq M(L,R)},
          (v_j)_{M(L,R)<j\leq R}\bigr)
  \]
  with one index on each side, and it occurs in no other query in the query
  list.

  Lean-facing, the formal statement for this label must expose both uniqueness
  results: uniqueness of the separating generated interval and uniqueness of
  the corresponding element of the concrete list
  $\operatorname{SplitQueries}_n(V)$.  It must not replace the query-list
  uniqueness conclusion by a second copy of endpoint-pair uniqueness, and it
  must not state only that some endpoint pair could be used to build a query.
  The query conclusion must identify the actual attached bichromatic query
  record, including its left vector list, right vector list, principal size,
  and dimension projections inherited from
  Definition~\ref{def:minip-balanced-recursive-split-tree}.
\end{lemma}
\begin{proof}
  \uses{def:minip-balanced-recursive-split-tree}
  Starting at the root, follow the unique child interval that contains both
  $p$ and $q$ as long as such a child exists.  This descent must stop before a
  singleton leaf because a singleton cannot contain two distinct indices.  At
  the stopping node, the two indices lie in different children, giving
  existence in the concrete list.  Uniqueness follows from the invariants built
  into Definition~\ref{def:minip-balanced-recursive-split-tree}: generated
  intervals are laminar and have no duplicate endpoint pairs.  Above the
  stopping node both indices are in the same child, below it no descendant
  contains both indices, and a disjoint interval contains neither both indices.
  Hence no second generated internal interval can satisfy the same separation
  predicate.
  The query-schedule statement follows because
  $\operatorname{SplitQueries}_n(V)$ is definitionally the list obtained by
  mapping each internal interval, in inherited order, to its attached
  rectangular query $Q_I(V)$.  The no-duplicate endpoint invariant gives a
  one-to-one correspondence between internal-interval positions and query-list
  positions.  Therefore the unique separating interval gives a unique actual
  query-list position, and unfolding $Q_I(V)$ gives exactly the displayed left
  and right row lists, size projection, and dimension projection.  Any other
  query record containing the same pair as a cross-pair would come from another
  separating internal interval, contradicting the interval uniqueness already
  proved.
\end{proof}

\begin{definition}[Rectangular bichromatic Min-IP oracle problem specs]
  \label{def:rectangular-bichromatic-minip-oracle-specs}
  \uses{def:bichromatic-min-ip, def:computational-problem-spec}
  From the rectangular input record of Definition~\ref{def:bichromatic-min-ip}
  we form three concrete computational-problem specifications.  Their input
  type is the same rectangular record
  \[
    X=(A:\{1,\ldots,n_A\}\to\{0,1\}^{\ell},
       B:\{1,\ldots,n_B\}\to\{0,1\}^{\ell}),
  \]
  with validity condition $n_A\geq1$ and $n_B\geq1$, principal size
  $N(X)=\max\{n_A,n_B\}$, and auxiliary dimension $L(X)=\ell$.

  The exact oracle specification
  $\operatorname{BiMinIPExactSpec}$ has output type
  $\{1,\ldots,n_A\}\times\{1,\ldots,n_B\}$.  Its correctness relation
  $\operatorname{exact\_correct}(X,(i^*,j^*))$ is exactly
  \[
    \forall i\in\{1,\ldots,n_A\},\ \forall j\in\{1,\ldots,n_B\},\qquad
    \ip{A_{i^*}}{B_{j^*}}\leq\ip{A_i}{B_j}.
  \]
  The $\gamma$-approximate oracle specification
  $\operatorname{BiMinIPGammaSpec}(\gamma)$, for $\gamma\geq1$, has the same
  output type.  Its correctness relation
  $\operatorname{gamma\_correct}(X,(i^*,j^*))$ is exactly
  \[
    \min_{i,j}\ip{A_i}{B_j}
      \leq \ip{A_{i^*}}{B_{j^*}}
      \leq \gamma\min_{i,j}\ip{A_i}{B_j}.
  \]
  The threshold oracle specification
  $\operatorname{BiMinIPThresholdSpec}(t)$, for $0\leq t\leq\ell$, has output
  type $\{0,1\}$.  Its correctness relation
  $\operatorname{threshold\_correct}(X,z)$ is exactly
  \[
    z=1\Longleftrightarrow
    \exists i\in\{1,\ldots,n_A\},\exists j\in\{1,\ldots,n_B\},\
      \ip{A_i}{B_j}\leq t,
  \]
  and $z=0$ in the negation of this condition.

  These are real computational-problem specifications in the sense of
  Definition~\ref{def:computational-problem-spec}.  They are not aliases for a
  family tag and they do not require any additional abstract correctness
  fields: the only output and correctness relations are the three displayed
  relations named $\operatorname{exact\_correct}$,
  $\operatorname{gamma\_correct}$, and
  $\operatorname{threshold\_correct}$.
\end{definition}

\begin{definition}[Promised monochromatic Min-IP source problem specs]
  \label{def:promised-monochromatic-minip-source-specs}
  \uses{def:min-ip, def:gamma-min-ip, def:min-ip-decision, def:computational-problem-spec}
  The monochromatic exact and $\gamma$-approximate Min-IP source specifications
  used in the recursive reduction have the input type
  $V:\{1,\ldots,n\}\to\{0,1\}^{\ell}$ but are valid only under the promise
  \[
    n\geq2.
  \]
  This promise is part of the valid-input predicate, because the output type
  for both optimization variants is the distinct-pair subtype
  \[
    \{(i,j):\{1,\ldots,n\}\times\{1,\ldots,n\}\mid i\neq j\}.
  \]
  The exact correctness relation requires the returned distinct pair
  $(i^*,j^*)$ to satisfy
  \[
    \forall i,j,\ i\neq j\Rightarrow
      \ip{v_{i^*}}{v_{j^*}}\leq\ip{v_i}{v_j}.
  \]
  The $\gamma$-approximate correctness relation, for $\gamma\geq1$, requires
  \[
    \min_{i\neq j}\ip{v_i}{v_j}
      \leq \ip{v_{i^*}}{v_{j^*}}
      \leq \gamma\min_{i\neq j}\ip{v_i}{v_j}.
  \]

  The threshold source specification
  $\operatorname{MonoMinIPThresholdSpec}(t)$ has Boolean output type and may be
  evaluated for all $n$; when $n<2$ its yes condition is false because there is
  no feasible distinct pair.  For $n\geq2$ it asks exactly whether some
  distinct pair has inner product at most $t$.
  Thus the optimization sources never quantify over an infeasible distinct
  pair type, while the decision source keeps the paper's constant-time
  $n<2$ no branch.
\end{definition}

\begin{lemma}[Split-tree query views are valid rectangular queries]
  \label{lem:minip-split-query-view-valid}
  \uses{def:minip-balanced-recursive-split-tree, def:bichromatic-min-ip, def:rectangular-bichromatic-minip-oracle-specs}
  Let $V:\{1,\ldots,n\}\to\{0,1\}^{\ell}$ and let
  $I=(L,R)$ be an internal interval of
  $\operatorname{InternalIntervals}_n$.  With
  $M=L+\lfloor(R-L)/2\rfloor$, the left child range
  $\{L,\ldots,M\}$ and the right child range $\{M+1,\ldots,R\}$ are both
  nonempty.  The concrete query view
  \[
    Q_I(V)=
      \bigl((v_i)_{L\leq i\leq M},(v_j)_{M<j\leq R}\bigr)
  \]
  is therefore a valid rectangular input for each of
  $\operatorname{BiMinIPExactSpec}$,
  $\operatorname{BiMinIPGammaSpec}(\gamma)$, and
  $\operatorname{BiMinIPThresholdSpec}(t)$.  Its principal size and dimension
  are exactly
  \[
    N(Q_I(V))=\max\{M-L+1,R-M\},\qquad L(Q_I(V))=\ell.
  \]
  The query uses row views into the original input with the displayed index
  maps; it does not pad, copy to an arbitrary list, or hide validity behind an
  existential rectangular-input witness.
\end{lemma}
\begin{proof}
  \uses{def:minip-balanced-recursive-split-tree, def:bichromatic-min-ip, def:rectangular-bichromatic-minip-oracle-specs}
  Since $I$ is internal, $L<R$.  The midpoint satisfies $L\leq M<R$, so both
  ranges $\{L,\ldots,M\}$ and $\{M+1,\ldots,R\}$ contain at least one index.
  Substituting these nonempty ranges into
  Definition~\ref{def:bichromatic-min-ip} gives a valid rectangular record.
  The displayed equations for principal size and dimension are the projections
  of that record.
\end{proof}

\begin{definition}[Affine local index maps for split-query answers]
  \label{def:minip-split-answer-affine-index-maps}
  \uses{lem:minip-split-query-view-valid, def:minip-balanced-recursive-split-tree}
  For an internal interval $I=(L,R)$ of
  $\operatorname{InternalIntervals}_n$, write
  $M=L+\lfloor(R-L)/2\rfloor$.  The local left index type of the attached
  rectangular query is
  \[
    U_I=\{1,\ldots,M-L+1\},
  \]
  and the local right index type is
  \[
    W_I=\{1,\ldots,R-M\}.
  \]
  Define the two concrete affine inclusions
  \[
    \operatorname{leftIndex}_I:U_I\to\{1,\ldots,n\},
    \qquad
    \operatorname{rightIndex}_I:W_I\to\{1,\ldots,n\}
  \]
  by the equations
  \[
    \operatorname{leftIndex}_I(u)=L+u-1,
    \qquad
    \operatorname{rightIndex}_I(v)=M+v.
  \]
  If the rectangular query is implemented using the child-range subtype
  indices $\{L,\ldots,M\}$ and $\{M+1,\ldots,R\}$ instead of local ordinals,
  these same maps are the coercions to the underlying original row indices.
  This node defines only the two index maps and their displayed equations.
\end{definition}

\begin{lemma}[Affine split-answer indices are valid and distinct]
  \label{lem:minip-split-answer-affine-index-valid}
  \uses{def:minip-split-answer-affine-index-maps, lem:minip-split-query-view-valid}
  For every internal interval $I=(L,R)$ with midpoint $M$, every
  $u\in\{1,\ldots,M-L+1\}$, and every
  $v\in\{1,\ldots,R-M\}$,
  \[
    L\leq \operatorname{leftIndex}_I(u)\leq M,
    \qquad
    M<\operatorname{rightIndex}_I(v)\leq R.
  \]
  Consequently both indices are original input indices in $\{1,\ldots,n\}$ and
  \[
    \operatorname{leftIndex}_I(u)\neq
    \operatorname{rightIndex}_I(v).
  \]
\end{lemma}
\begin{proof}
  \uses{def:minip-split-answer-affine-index-maps, lem:minip-split-query-view-valid}
  The bounds on $u$ give $0\leq u-1\leq M-L$, hence
  $L\leq L+u-1\leq M$.  The bounds on $v$ give $1\leq v\leq R-M$, hence
  $M< M+v\leq R$.  Since the first index is at most $M$ and the second is
  strictly greater than $M$, they are distinct.
\end{proof}

\begin{definition}[Local row functions projected from a split-query view]
  \label{def:minip-split-query-local-row-functions}
  \uses{def:minip-split-answer-affine-index-maps, lem:minip-split-query-view-valid}
  Let $V:\{1,\ldots,n\}\to\{0,1\}^{\ell}$ and let $I=(L,R)$ be an internal
  interval with midpoint $M$.  Let
  \[
    Q_I(V)=
      \bigl((v_i)_{L\leq i\leq M},(v_j)_{M<j\leq R}\bigr)
  \]
  be the actual rectangular query view from
  Lemma~\ref{lem:minip-split-query-view-valid}.  Define
  \[
    \operatorname{SplitLeftRows}_I(V):
      \{1,\ldots,M-L+1\}\to\{0,1\}^{\ell}
  \]
  and
  \[
    \operatorname{SplitRightRows}_I(V):
      \{1,\ldots,R-M\}\to\{0,1\}^{\ell}
  \]
  to be the first and second row functions obtained by projecting this same
  query record $Q_I(V)$ and reindexing its child-range domains by the canonical
  local ordinal maps
  \[
    u\mapsto L+u-1,\qquad v\mapsto M+v.
  \]
  Equivalently,
  \[
    \bigl(\operatorname{SplitLeftRows}_I(V),
          \operatorname{SplitRightRows}_I(V)\bigr)
  \]
  is the local-ordinal row view of the concrete rectangular query
  $Q_I(V)$.  This node defines only these two row functions as projections of
  the actual split-query view.  It does not define them by an independent
  direct lookup into $V$, does not accept caller-supplied row functions, and
  does not add any oracle-validity hypothesis beyond the fact that $I$ is an
  internal interval whose query view is the one displayed above.
\end{definition}

\begin{lemma}[Left split-query row projection has the affine index]
  \label{lem:minip-split-query-left-row-affine}
  \uses{def:minip-split-query-local-row-functions, def:minip-split-answer-affine-index-maps}
  Let $V:\{1,\ldots,n\}\to\{0,1\}^{\ell}$ and let $I=(L,R)$ be an internal
  interval with midpoint $M$.  For every
  $u\in\{1,\ldots,M-L+1\}$,
  \[
    \operatorname{SplitLeftRows}_I(V)(u)
      =
    v_{\operatorname{leftIndex}_I(u)}.
  \]
\end{lemma}
\begin{proof}
  \uses{def:minip-split-query-local-row-functions, def:minip-split-answer-affine-index-maps}
  Unfold the local row-function projection from
  Definition~\ref{def:minip-split-query-local-row-functions}.  The left row
  function is the first component of the actual query record $Q_I(V)$,
  reindexed by $u\mapsto L+u-1$.  Unfolding
  Definition~\ref{def:minip-split-answer-affine-index-maps} identifies this
  original row index with $\operatorname{leftIndex}_I(u)$.
\end{proof}

\begin{lemma}[Right split-query row projection has the affine index]
  \label{lem:minip-split-query-right-row-affine}
  \uses{def:minip-split-query-local-row-functions, def:minip-split-answer-affine-index-maps}
  Let $V:\{1,\ldots,n\}\to\{0,1\}^{\ell}$ and let $I=(L,R)$ be an internal
  interval with midpoint $M$.  For every
  $w\in\{1,\ldots,R-M\}$,
  \[
    \operatorname{SplitRightRows}_I(V)(w)
      =
    v_{\operatorname{rightIndex}_I(w)}.
  \]
\end{lemma}
\begin{proof}
  \uses{def:minip-split-query-local-row-functions, def:minip-split-answer-affine-index-maps}
  Unfold the local row-function projection from
  Definition~\ref{def:minip-split-query-local-row-functions}.  The right row
  function is the second component of the actual query record $Q_I(V)$,
  reindexed by $w\mapsto M+w$.  Unfolding
  Definition~\ref{def:minip-split-answer-affine-index-maps} identifies this
  original row index with $\operatorname{rightIndex}_I(w)$.
\end{proof}

\begin{lemma}[Affine split-answer indices preserve query rows]
  \label{lem:minip-split-answer-row-preservation}
  \uses{def:minip-split-query-local-row-functions, lem:minip-split-query-left-row-affine, lem:minip-split-query-right-row-affine}
  Let $V:\{1,\ldots,n\}\to\{0,1\}^{\ell}$ and let $I=(L,R)$ be an internal
  interval with midpoint $M$.  Let
  \[
    U_I=\{1,\ldots,M-L+1\},
    \qquad
    W_I=\{1,\ldots,R-M\}
  \]
  be exactly the local left and right answer-index types introduced in
  Definition~\ref{def:minip-split-answer-affine-index-maps}.  In the query
  \[
    Q_I(V)=
      \bigl((v_i)_{L\leq i\leq M},(v_j)_{M<j\leq R}\bigr),
  \]
  let the public row functions be exactly the local-ordinal projections of
  this query record:
  \[
    A=\operatorname{SplitLeftRows}_I(V),
    \qquad
    B=\operatorname{SplitRightRows}_I(V).
  \]
  Thus $A$ and $B$ are the rows produced by the split-query view, not
  caller-supplied functions and not independent direct lookups into $V$.  For
  every
  \[
    u:U_I,\qquad w:W_I,
  \]
  \[
    A(u)=v_{\operatorname{leftIndex}_I(u)},
    \qquad
    B(w)=v_{\operatorname{rightIndex}_I(w)}.
  \]
  Therefore
  \[
    \ip{A(u)}{B(w)}
      =
    \ip{v_{\operatorname{leftIndex}_I(u)}}
       {v_{\operatorname{rightIndex}_I(w)}}.
  \]
  Lean-facing, the declaration
  $\operatorname{lem\_minip\_split\_answer\_row\_preservation}$ must quantify
  directly over these displayed local ordinal types:
  \[
    \forall\,u:U_I,\ \forall\,w:W_I,\quad
      A(u)=v_{\operatorname{leftIndex}_I(u)}
      \ \wedge\
      B(w)=v_{\operatorname{rightIndex}_I(w)}
      \ \wedge\
      \ip{A(u)}{B(w)}
        =
      \ip{v_{\operatorname{leftIndex}_I(u)}}
         {v_{\operatorname{rightIndex}_I(w)}}.
  \]
  Here $A$ and $B$ are fixed by the preceding displayed equations as
  $\operatorname{SplitLeftRows}_I(V)$ and
  $\operatorname{SplitRightRows}_I(V)$, the two projections of the actual
  query $Q_I(V)$.  The declaration has no premise asserting validity of an
  unrelated oracle problem instance and no parameter supplying alternate row
  functions.
  It is not enough to quantify over implementation-list indices such as
  $\operatorname{Fin}(\operatorname{leftRows.length})$ and
  $\operatorname{Fin}(\operatorname{rightRows.length})$ with extra
  caller-supplied bound hypotheses.  If the rectangular query is internally
  stored as lists, the implementation must first expose the definitional
  identification of those list index types with $U_I$ and $W_I$, and the
  preservation lemma for this label must still have $u:U_I$ and $v:W_I$ as its
  public quantified variables.  No hypotheses named like $h_u$, $h_v$,
  $\operatorname{hu}$, or $\operatorname{hv}$ may be required to turn an
  implementation-list index into a valid local answer index.
\end{lemma}
\begin{proof}
  \uses{lem:minip-split-query-left-row-affine, lem:minip-split-query-right-row-affine}
  Fix arbitrary $u:U_I$ and $w:W_I$.  Unfold the rectangular query view
  $Q_I(V)$ only through the projected row functions
  $\operatorname{SplitLeftRows}_I(V)$ and
  $\operatorname{SplitRightRows}_I(V)$.  Lemma~\ref{lem:minip-split-query-left-row-affine}
  gives
  $A(u)=v_{\operatorname{leftIndex}_I(u)}$, and
  Lemma~\ref{lem:minip-split-query-right-row-affine} gives
  $B(w)=v_{\operatorname{rightIndex}_I(w)}$.  Substituting these two
  projection equalities into the inner product gives the displayed equality,
  still for the same quantified variables $u:U_I$ and $w:W_I$.
\end{proof}

\begin{definition}[Local-to-global decoding of split-query answers]
  \label{def:minip-split-answer-local-to-global-decoding}
  \uses{def:minip-split-answer-affine-index-maps, lem:minip-split-answer-affine-index-valid, lem:minip-split-answer-row-preservation}
  For an internal interval $I=(L,R)$ with midpoint $M$, define the deterministic
  decoding map
  \[
    \operatorname{DecodePair}_I:
      \{1,\ldots,M-L+1\}\times\{1,\ldots,R-M\}
      \to
      \{(a,b):\{1,\ldots,n\}\times\{1,\ldots,n\}\mid a\neq b\}
  \]
  by
  \[
    \operatorname{DecodePair}_I(u,v)
      =
    \bigl(\operatorname{leftIndex}_I(u),
          \operatorname{rightIndex}_I(v)\bigr),
  \]
  where membership in the distinct-pair subtype is certified by
  Lemma~\ref{lem:minip-split-answer-affine-index-valid}.  The row-value and
  inner-product preservation facts for this decoded pair are exactly
  Lemma~\ref{lem:minip-split-answer-row-preservation}.

  Lean-facing, the declaration
  \[
    \operatorname{def\_minip\_split\_answer\_local\_to\_global\_decoding}
      (n,\ell,V,I)
  \]
  returns this function $\operatorname{DecodePair}_I$ itself.  It is not a
  proposition, interface, or structure with a caller-supplied decoding field.
  Its only computational body is the displayed pair of affine index maps; the
  bounds, distinctness, row-value preservation, and inner-product preservation
  obligations are separate projection lemmas in the \uses list.
\end{definition}

\begin{definition}[Decoded optimization candidates from a split history]
  \label{def:minip-split-history-decoded-candidates}
  \uses{def:minip-split-answer-local-to-global-decoding, def:rectangular-bichromatic-minip-oracle-specs}
  For a full optimization history
  \[
    H=((Q_{I_1}(V),z_1),\ldots,(Q_{I_m}(V),z_m))
  \]
  in the same list order as $\operatorname{InternalIntervals}_n$, with
  optimization oracle answers $z_s=(u_s,v_s)$, define the decoded candidate at
  history position $s$ to be
  \[
    \operatorname{CandPair}(V,H,s)
      =
    (a_s,b_s)=\operatorname{DecodePair}_{I_s}(u_s,v_s).
  \]
  Define its score by
  \[
    \operatorname{CandBeta}(V,H,s)=\ip{v_{a_s}}{v_{b_s}}.
  \]
  This node defines only the finite indexed candidate-pair and beta-value
  functions obtained by decoding each history position.
\end{definition}

\begin{definition}[First minimum beta position in split history]
  \label{def:minip-split-history-first-min-beta}
  \uses{def:minip-split-history-decoded-candidates}
  For a nonempty optimization history $H$ of length $m>0$, define
  \[
    \operatorname{MinBeta}(V,H)
      =
    \min_{1\leq r\leq m}\operatorname{CandBeta}(V,H,r).
  \]
  The selected history position is the natural-order first position attaining
  this value:
  \[
    \operatorname{FirstMinBetaPos}(V,H)
      =
    \min\{s\in\{1,\ldots,m\}:
      \operatorname{CandBeta}(V,H,s)=\operatorname{MinBeta}(V,H)\}.
  \]
  Tie-breaking is by the finite history position $s$, not by the decoded pair
  value and not by a list search on pair values.  If two different history
  positions decode to the same global pair, they remain separate positions in
  this minimum scan.

  Lean-facing, this declaration is not just two equations assigning
  $\operatorname{MinBeta}$ to the score at some implementation-selected index.
  It returns the concrete finite-minimum package
  \[
    \operatorname{def\_minip\_split\_history\_first\_min\_beta}(V,H)
    =
    \bigl(\operatorname{MinBeta}(V,H),
          \operatorname{FirstMinBetaPos}(V,H),
          h_{\rm in},h_{\rm att},h_{\rm le},h_{\rm first}\bigr)
  \]
  with exactly the following eliminable fields.
  First,
  \[
    h_{\rm in}:\quad
    1\leq\operatorname{FirstMinBetaPos}(V,H)\leq m .
  \]
  Second, the selected position attains the displayed finite minimum:
  \[
    h_{\rm att}:\quad
    \operatorname{CandBeta}
      \bigl(V,H,\operatorname{FirstMinBetaPos}(V,H)\bigr)
      =
    \operatorname{MinBeta}(V,H).
  \]
  Third, the selected value is no larger than every history-position score:
  \[
    h_{\rm le}:\quad
    \forall r\in\{1,\ldots,m\},\qquad
      \operatorname{MinBeta}(V,H)
      \leq
      \operatorname{CandBeta}(V,H,r).
  \]
  Fourth, the selected position is the natural-order first minimizer:
  \[
    h_{\rm first}:\quad
    \forall r\in\{1,\ldots,m\},\quad
      r<\operatorname{FirstMinBetaPos}(V,H)
      \Rightarrow
      \operatorname{CandBeta}(V,H,r)
      \neq
      \operatorname{MinBeta}(V,H).
  \]
  Equivalently, the pair $(h_{\rm att},h_{\rm first})$ states that
  $\operatorname{FirstMinBetaPos}(V,H)$ is the least natural-number position
  among all positions attaining $\operatorname{MinBeta}(V,H)$.

  These four fields are part of this node's formal content.  A declaration for
  this label is invalid if it merely chooses an index such as
  $\operatorname{minIPSplitFirstIndexByNat}$ and then defines
  $\operatorname{MinBeta}$ to be
  $\operatorname{CandBeta}$ at that index without proving $h_{\rm le}$ and
  $h_{\rm first}$.  The minimum property must quantify over the actual finite
  position type $\{1,\ldots,m\}$ of the supplied history $H$, and the
  firstness property must use the natural order on those history positions.
  No decoded-pair tie-breaking, pair-value search, existential minimizer, or
  caller-supplied proof of optimality may replace these fields.
\end{definition}

\begin{definition}[Pair-valued split-history finish scan]
  \label{def:minip-split-finish-pair-scan}
  \uses{def:minip-split-history-decoded-candidates, def:minip-split-history-first-min-beta}
  For a nonempty optimization history $H$, define the concrete pair-valued
  finish function by
  \[
    \operatorname{def\_minip\_split\_finish\_pair}(V,H)
      =
    \operatorname{CandPair}
      \bigl(V,H,\operatorname{FirstMinBetaPos}(V,H)\bigr).
  \]
  If a total formal function is required on empty histories, its empty-history
  branch is an arbitrary fixed distinct-pair default supplied only after a
  proof that the optimization source is invalid there; the defining equation
  above is the only equation used for valid optimization histories.
\end{definition}

\begin{definition}[Threshold split-history Boolean OR scan]
  \label{def:minip-split-finish-threshold-scan}
  \uses{def:rectangular-bichromatic-minip-oracle-specs}
  For a threshold history
  \[
    H=((Q_{I_1}(V),z_1),\ldots,(Q_{I_m}(V),z_m)),
    \qquad z_s\in\{0,1\},
  \]
  define the concrete Boolean finish function by the finite OR
  \[
    \operatorname{def\_minip\_split\_finish\_threshold}(H)
      =
    \bigvee_{s=1}^{m} z_s.
  \]
  For $m=0$ this OR is the empty finite OR and has value $0$.
\end{definition}

\begin{definition}[Deterministic split-history postprocessing]
  \label{def:minip-split-history-postprocessing}
  \uses{def:minip-split-finish-pair-scan, def:minip-split-finish-threshold-scan}
  The deterministic split-history postprocessors are exactly the two concrete
  functions
  \[
    \operatorname{def\_minip\_split\_finish\_pair}
    \qquad\text{and}\qquad
    \operatorname{def\_minip\_split\_finish\_threshold}
  \]
  defined in Definitions~\ref{def:minip-split-finish-pair-scan} and
  \ref{def:minip-split-finish-threshold-scan}.  The exact and
  $\gamma$-approximate variants use the pair-valued first-minimum beta scan;
  the threshold variant uses the Boolean OR scan.

  Lean-facing, this node is not a proposition stating equations about
  externally supplied functions, and it is not a record accepting
  $\operatorname{FinishPair}$ or $\operatorname{FinishThreshold}$ fields from
  the caller.  It exposes these two already-defined functions as the
  postprocessing functions for the split-tree reduction.  The pair-valued
  function's nonempty-history equation and the threshold function's empty and
  nonempty OR equations are inherited from the two helper definitions.
\end{definition}

\begin{definition}[Threshold-indexed Min-IP split reduction variant]
  \label{def:minip-threshold-indexed-reduction-variant}
  \uses{def:promised-monochromatic-minip-source-specs, def:rectangular-bichromatic-minip-oracle-specs}
  For a fixed dimension $\ell$, the split-tree Min-IP reduction variant is the
  following concrete tagged parameter:
  \[
    \operatorname{MinIPSplitReductionVariant}(\ell)
      =
      \operatorname{exact}
      \;\sqcup\;
      \operatorname{gamma}(\gamma,h_\gamma)
      \;\sqcup\;
      \operatorname{threshold}(t,h_{0},h_{\ell}).
  \]
  Here $h_\gamma$ is the proof $\gamma\geq1$, $h_0$ is the proof
  $0\leq t$, and $h_\ell$ is the proof $t\leq\ell$.  The threshold tag is not
  a nullary constructor: it carries the actual natural threshold $t$ and both
  displayed bounds.  Define
  \[
    \operatorname{thresholdOf}
      (\operatorname{threshold}(t,h_0,h_\ell))=t
  \]
  with projection lemmas returning $h_0$ and $h_\ell$.

  The variant determines the source and target problem specifications by the
  following equations:
  \[
    P_{\operatorname{exact}}=\operatorname{MonoMinIPExactSpec},\qquad
    Q_{\operatorname{exact}}=\operatorname{BiMinIPExactSpec},
  \]
  \[
    P_{\operatorname{gamma}(\gamma,h_\gamma)}
      =\operatorname{MonoMinIPGammaSpec}(\gamma),\qquad
    Q_{\operatorname{gamma}(\gamma,h_\gamma)}
      =\operatorname{BiMinIPGammaSpec}(\gamma),
  \]
  and
  \[
    P_{\operatorname{threshold}(t,h_0,h_\ell)}
      =\operatorname{MonoMinIPThresholdSpec}(t),\qquad
    Q_{\operatorname{threshold}(t,h_0,h_\ell)}
      =\operatorname{BiMinIPThresholdSpec}(t).
  \]

  Lean-facing, the declaration
  \[
    \operatorname{def\_minip\_threshold\_indexed\_reduction\_variant}(\ell)
  \]
  is the dependent package consisting of this parameter type and the two
  specification selectors
  \[
    \operatorname{sourceSpec}_{\ell}:
      \operatorname{MinIPSplitReductionVariant}(\ell)
      \to\operatorname{ComputationalProblemSpec},
  \]
  \[
    \operatorname{oracleSpec}_{\ell}:
      \operatorname{MinIPSplitReductionVariant}(\ell)
      \to\operatorname{ComputationalProblemSpec},
  \]
  defined by case analysis on the displayed constructors.  The threshold
  branch equations of these returned selector functions are definitional
  projection equations:
  \[
    \operatorname{sourceSpec}_{\ell}
      (\operatorname{threshold}(t,h_0,h_\ell))
      =
    \operatorname{MonoMinIPThresholdSpec}(t),
  \]
  \[
    \operatorname{oracleSpec}_{\ell}
      (\operatorname{threshold}(t,h_0,h_\ell))
      =
    \operatorname{BiMinIPThresholdSpec}(t).
  \]
  The same carried threshold variable $t$ appears on both right-hand sides,
  before either specification is formed.

  A formalization must not use a frozen constructor $\operatorname{threshold}$
  with no threshold field, must not default the threshold to $0$, and must not
  recover $t$ from a source input after the target oracle specification has
  already been chosen.  It must also not implement
  $\operatorname{sourceSpec}_{\ell}(\operatorname{threshold}(t,h_0,h_\ell))$
  by returning a tag-only or tuple-projection specification such as an
  unapplied threshold component of
  Definition~\ref{def:promised-monochromatic-minip-source-specs}.  The
  threshold selector must apply that source specification to the carried
  $t$, so eliminating the threshold case yields the exact equality
  $P_{\operatorname{threshold}(t,h_0,h_\ell)}
  =\operatorname{MonoMinIPThresholdSpec}(t)$, not merely a threshold-family
  tag.  The fixed threshold is part of the variant datum before both
  $\operatorname{MonoMinIPThresholdSpec}(t)$ and
  $\operatorname{BiMinIPThresholdSpec}(t)$ are formed.
\end{definition}

\begin{definition}[Fixed-threshold source and oracle specification alignment]
  \label{def:minip-fixed-threshold-source-oracle-spec}
  \uses{def:minip-threshold-indexed-reduction-variant, def:min-ip-decision, def:bichromatic-min-ip}
  Given a threshold variant
  \[
    \theta=\operatorname{threshold}(t,h_0,h_\ell)
    \in\operatorname{MinIPSplitReductionVariant}(\ell),
  \]
  define the threshold split reduction's paired problem specifications
  directly from this constructor pattern, not by reusing a possibly
  threshold-erasing selector:
  \[
    \operatorname{FixedThresholdSourceSpec}_{\ell}(\theta)
      :=\operatorname{MonoMinIPThresholdSpec}(t),
    \qquad
    \operatorname{FixedThresholdOracleSpec}_{\ell}(\theta)
      :=\operatorname{BiMinIPThresholdSpec}(t).
  \]
  The projection equations
  \[
    P_\theta
      =
    \operatorname{FixedThresholdSourceSpec}_{\ell}(\theta)
      =
    \operatorname{MonoMinIPThresholdSpec}(t),
  \]
  and
  \[
    Q_\theta
      =
    \operatorname{FixedThresholdOracleSpec}_{\ell}(\theta)
      =
    \operatorname{BiMinIPThresholdSpec}(t)
  \]
  are part of this node's returned data.
  For a monochromatic source input
  $V:\{1,\ldots,n\}\to\{0,1\}^{\ell}$, the source yes condition is
  \[
    \exists i,j,\ i\neq j\wedge \ip{v_i}{v_j}\leq t.
  \]
  For every split-tree query
  $Q_I(V)=((v_i)_{L\leq i\leq M},(v_j)_{M<j\leq R})$, the target oracle yes
  condition for the same fixed $\theta$ is
  \[
    \exists u,w,\quad
      \ip{\operatorname{SplitLeftRows}_I(V)(u)}
         {\operatorname{SplitRightRows}_I(V)(w)}
      \leq t.
  \]
  The threshold appearing in these two displayed predicates is definitionally
  the same value $\operatorname{thresholdOf}(\theta)=t$.

  This node is the threshold-source specification restricted to the same fixed
  $t$ as the rectangular bichromatic oracle specification.  It supplies the
  declaration
  \[
    \operatorname{def\_minip\_fixed\_threshold\_source\_oracle\_spec}
      (\ell,\theta)
  \]
  whose output is the concrete tuple
  \[
    \bigl(
      \operatorname{MonoMinIPThresholdSpec}(t),
      \operatorname{BiMinIPThresholdSpec}(t),
      \operatorname{sourceYes}_t,
      \operatorname{queryYes}_t,
      h_P,h_Q
    \bigr),
  \]
  where $h_P$ and $h_Q$ are exactly the two displayed specification equalities.
  The source yes predicate and query yes predicate are projections from these
  two concrete threshold-$t$ specifications.  They are not independent generic
  predicates whose threshold happens to be written as
  $\operatorname{thresholdOf}(\theta)$ while the source specification itself
  remains untied to $t$.

  It is not an oracle-reduction datum and it is not a proof of correctness of
  the whole reduction.  It only fixes the matching threshold-bearing source
  and target specifications so later reduction data cannot instantiate the
  source as an unapplied threshold family, instantiate the target as
  $\operatorname{BiMinIPThresholdSpec}(0)$, as a tag-only threshold problem, or
  as a specification whose threshold depends on a query answer or on an
  unrelated source-input field.
\end{definition}

\begin{definition}[Oracle history to full optimization split history adapter]
  \label{def:minip-oracle-history-to-full-optimization-history}
  \uses{def:oracle-reduction-answer-history, def:minip-split-history-decoded-candidates, def:minip-split-finish-pair-scan}
  For the exact and $\gamma$ split-tree variants, let
  $H_{\rm raw}\in\Hist_R(V,r(V))$ be a full oracle-answer history of the
  oracle reduction datum whose queries are the split queries
  $Q_{\operatorname{InternalIntervals}_n[s]}(V)$ and whose answers are
  optimization pairs $z_s=(u_s,w_s)$.  Define
  \[
    \operatorname{OptHist}(V,H_{\rm raw})
      :\operatorname{MinIPSplitFullOptimizationHistory}(V)
  \]
  to be the concrete full optimization history
  \[
    ((Q_{I_1}(V),z_1),\ldots,(Q_{I_m}(V),z_m)),
    \qquad
    I_s=\operatorname{InternalIntervals}_n[s],
  \]
  in the same list order as $\operatorname{InternalIntervals}_n$, where each
  query component and answer component is copied from the raw query-answer
  entry at position $s$.

  Its defining projection equations are:
  \[
    \operatorname{length}(\operatorname{OptHist}(V,H_{\rm raw}))=r(V),
  \]
  \[
    \operatorname{OptHist}(V,H_{\rm raw})[s].\operatorname{query}
      =
    Q_{\operatorname{InternalIntervals}_n[s]}(V),
  \]
  and
  \[
    \operatorname{OptHist}(V,H_{\rm raw})[s].\operatorname{answer}=z_s
  \]
  for every $1\leq s\leq r(V)$.  Consequently the decoded candidates and beta
  values used by
  Definition~\ref{def:minip-split-history-decoded-candidates} are computed
  from the actual oracle answers in $H_{\rm raw}$, at the same positions.

  Lean-facing, this node is the function
  \[
    \operatorname{def\_minip\_oracle\_history\_to\_full\_optimization\_history}
      (V,H_{\rm raw})
      =
    \operatorname{OptHist}(V,H_{\rm raw}).
  \]
  It is the only adapter from the generic oracle-answer history type to
  $\operatorname{MinIPSplitFullOptimizationHistory}(V)$.  The exact and
  $\gamma$ finish functions must therefore be stated as
  \[
    \operatorname{def\_minip\_split\_finish\_pair}
      \bigl(V,\operatorname{OptHist}(V,H_{\rm raw})\bigr),
  \]
  not as a fixed fallback pair, an empty-history default, or a call to
  $\operatorname{def\_minip\_split\_finish\_pair}$ on the raw history type
  itself.
\end{definition}

\begin{definition}[Split-tree local phase counter schedule]
  \label{def:minip-bi-to-mono-local-phase-counter-schedule}
  \uses{def:oracle-reduction-local-phases, def:minip-balanced-recursive-split-tree, lem:minip-split-query-view-valid, def:minip-split-history-postprocessing}
  For the split-tree reduction on input $V$, define
  \[
    m=\operatorname{length}(\operatorname{InternalIntervals}_n),
    \qquad
    \State_R(V)=\{0,\ldots,m\}.
  \]
  The initial state is the counter value
  \[
    \init_R(V)=0.
  \]
  At oracle-call phase $s$, where $1\leq s\leq m$, the previous compatible
  state is the counter value $s-1$, and the update after any well-typed query
  answer is exactly
  \[
    \update_R(V,s,\sigma_{s-1},H_s,q_s,z_s)=s.
  \]
  This state contains no query answer and no history; the full answer history
  remains the separate history argument of
  Definition~\ref{def:oracle-reduction-answer-history}.
\end{definition}

\begin{definition}[Split-tree local phase query schedule]
  \label{def:minip-bi-to-mono-local-phase-query-schedule}
  \uses{def:minip-bi-to-mono-local-phase-counter-schedule, lem:minip-split-query-view-valid}
  At oracle-call phase $s$ with $1\leq s\leq m$, the query constructor reads
  the interval
  \[
    I_s=\operatorname{InternalIntervals}_n[s]
  \]
  and returns exactly the valid rectangular query view
  \[
    \query_R(V,s,\sigma_{s-1},H_s)=Q_{\operatorname{InternalIntervals}_n[s]}(V),
  \]
  independent of the previous history argument $H_s$ and independent of the
  previous counter value except for its phase-compatibility role.  The query
  validity and its principal-size and dimension projections are those of
  Lemma~\ref{lem:minip-split-query-view-valid}.
\end{definition}

\begin{definition}[Split-tree local phase variant finish functions]
  \label{def:minip-bi-to-mono-local-phase-variant-finish}
  \uses{def:minip-bi-to-mono-local-phase-counter-schedule, def:minip-split-history-postprocessing, def:minip-oracle-history-to-full-optimization-history}
  The final phase of the local presentation is variant tagged.  For the exact
  optimization variant, the output type is the exact promised monochromatic
  Min-IP distinct-pair type and
  \[
    \finish_R^{\rm exact}(V,\sigma_m,H_{\rm raw})
      =
    \operatorname{def\_minip\_split\_finish\_pair}
      \bigl(V,\operatorname{OptHist}(V,H_{\rm raw})\bigr).
  \]
  For the $\gamma$-approximate optimization variant, the output type is the
  $\gamma$-approximate promised monochromatic Min-IP distinct-pair type and
  \[
    \finish_R^{\gamma}(V,\sigma_m,H_{\rm raw})
      =
    \operatorname{def\_minip\_split\_finish\_pair}
      \bigl(V,\operatorname{OptHist}(V,H_{\rm raw})\bigr).
  \]
  For the threshold decision variant, the output type is Boolean and
  \[
    \finish_R^{\rm thresh}(V,\sigma_m,H)
      =
    \operatorname{def\_minip\_split\_finish\_threshold}(H).
  \]
  There is no common output value containing both pair and Boolean components,
  and no later projection chooses a component.  Each variant has its own
  finish function with the displayed output type.  In the two optimization
  variants the raw full oracle-answer history is first adapted by
  Definition~\ref{def:minip-oracle-history-to-full-optimization-history};
  therefore the pair scan receives the concrete
  $\operatorname{MinIPSplitFullOptimizationHistory}(V)$ required by
  Definition~\ref{def:minip-split-finish-pair-scan}.
\end{definition}

\begin{definition}[Split-tree Min-IP local-phase presentation]
  \label{def:minip-bi-to-mono-local-phase-presentation}
  \uses{def:minip-threshold-indexed-reduction-variant, def:minip-bi-to-mono-local-phase-counter-schedule, def:minip-bi-to-mono-local-phase-query-schedule, def:minip-bi-to-mono-local-phase-variant-finish}
  For the split-tree reduction on input $V$, the local-phase presentation is
  the concrete variant-tagged presentation whose components are exactly the
  counter schedule, query schedule, and finish functions of
  Definitions~\ref{def:minip-bi-to-mono-local-phase-counter-schedule},
  \ref{def:minip-bi-to-mono-local-phase-query-schedule}, and
  \ref{def:minip-bi-to-mono-local-phase-variant-finish}.  Its oracle-call
  count is
  \[
    r(V)=\operatorname{length}(\operatorname{InternalIntervals}_n),
  \]
  its state space is $\{0,\ldots,r(V)\}$, its initial state is $0$, its
  phase-$s$ query is the $s$th split query view, and its update after phase
  $s$ is the counter value $s$.

  The finish component is selected by the variant tag before the presentation
  is formed: exact and $\gamma$-approximate variants use the concrete
  pair-valued beta-minimum scan
  $\operatorname{def\_minip\_split\_finish\_pair}
  (V,\operatorname{OptHist}(V,H_{\rm raw}))$, while the threshold variant uses
  the concrete Boolean OR scan
  $\operatorname{def\_minip\_split\_finish\_threshold}(H)$.  The local work is
  precisely split-tree construction, query-view construction, counter updates,
  and the final scan attached to the selected variant.

  Lean-facing, the declaration
  \[
    \operatorname{def\_minip\_bi\_to\_mono\_local\_phase\_presentation}
      (\operatorname{variant},V)
  \]
  returns this concrete local-phase presentation.  Its component equations are
  definitional equations of the returned presentation:
  \[
    m=\operatorname{length}(\operatorname{InternalIntervals}_n),\quad
    \State_R(V)=\{0,\ldots,m\},\quad
    \init_R(V)=0,
  \]
  the query at phase $s$ is exactly the $s$th split query view, the update is
  exactly the counter value $s$, and the finish function is the
  variant-specific function displayed in
  Definition~\ref{def:minip-bi-to-mono-local-phase-variant-finish}.  The
  declaration must not be an abstract tuple of components that can be
  instantiated independently of these formulas, and it must not return a
  product of unrelated query, update, and finish components.  The returned
  presentation is connected to the oracle-reduction datum by these displayed
  equations for all three variants.  The threshold case uses only a variant
  value $\operatorname{threshold}(t,h_0,h_\ell)$ from
  Definition~\ref{def:minip-threshold-indexed-reduction-variant}; hence the
  same fixed threshold $t$ is available before the threshold source and target
  specifications are selected.
\end{definition}

\begin{definition}[Bichromatic-to-monochromatic Min-IP recursive reduction datum]
  \label{def:minip-bi-to-mono-reduction-datum}
  \uses{def:promised-monochromatic-minip-source-specs, def:rectangular-bichromatic-minip-oracle-specs, def:minip-threshold-indexed-reduction-variant, def:minip-fixed-threshold-source-oracle-spec, lem:minip-split-query-view-valid, def:minip-split-answer-local-to-global-decoding, def:minip-split-history-postprocessing, def:minip-oracle-history-to-full-optimization-history, def:minip-bi-to-mono-local-phase-presentation, def:oracle-reduction-data}
  Fix one of the three Min-IP variants: exact optimization, $\gamma$-
  approximation for a fixed $\gamma\geq1$, or threshold decision represented
  by a fixed threshold-indexed parameter
  $\operatorname{threshold}(t,h_0,h_\ell)$ with $0\leq t\leq\ell$.  The source
  problem is the corresponding promised
  monochromatic specification from
  Definition~\ref{def:promised-monochromatic-minip-source-specs}, and the
  target oracle problem is the corresponding rectangular bichromatic
  specification from
  Definition~\ref{def:rectangular-bichromatic-minip-oracle-specs}.  Given a
  correct rectangular bichromatic solver $A_{\rm bi}$ for that same variant,
  the concrete oracle reduction datum
  $R_{\rm biMinIP\to monoMinIP}$ is assembled from the helper declarations as
  follows.

  For a valid optimization source input, the promise gives $n\geq2$.  For a
  threshold source input with $n<2$, the datum has zero oracle calls and final
  output $0$.  Otherwise it uses the local-phase presentation of
  Definition~\ref{def:minip-bi-to-mono-local-phase-presentation}; hence its
  oracle-call count is
  \[
    r(V)=\operatorname{length}(\operatorname{InternalIntervals}_n),
  \]
  and its $s$th query is exactly the valid rectangular query view
  \[
    q_s(V)=Q_{\operatorname{InternalIntervals}_n[s]}(V).
  \]
  Query validity is provided by
  Lemma~\ref{lem:minip-split-query-view-valid}.  Optimization oracle answers
  are decoded into original distinct indices by
  Definition~\ref{def:minip-split-answer-local-to-global-decoding}; final
  exact and $\gamma$ outputs are computed by the deterministic beta-minimum
  rule of Definition~\ref{def:minip-split-history-postprocessing} after
  converting the raw full oracle history with
  Definition~\ref{def:minip-oracle-history-to-full-optimization-history}; and
  final threshold outputs are computed by its deterministic Boolean OR rule.

  Lean-facing, the declaration
  \[
    \operatorname{def\_minip\_bi\_to\_mono\_reduction\_datum}
      (\operatorname{variant})
  \]
  returns a value whose codomain is exactly the oracle-reduction-datum type
  \[
    \operatorname{def\_oracle\_reduction\_data}
      \bigl(P_{\operatorname{variant}},Q_{\operatorname{variant}}\bigr)
  \]
  from Definition~\ref{def:oracle-reduction-data}, where
  $P_{\operatorname{variant}}$ is the selected promised monochromatic source
  specification and $Q_{\operatorname{variant}}$ is the selected rectangular
  bichromatic oracle specification.  In the threshold case these two
  specifications are the aligned pair
  \[
    P_{\operatorname{threshold}(t,h_0,h_\ell)}
      =\operatorname{MonoMinIPThresholdSpec}(t),
    \qquad
    Q_{\operatorname{threshold}(t,h_0,h_\ell)}
      =\operatorname{BiMinIPThresholdSpec}(t)
  \]
  from Definition~\ref{def:minip-fixed-threshold-source-oracle-spec}.  Thus
  the target problem is not chosen from a threshold-erasing tag and does not
  default to $t=0$.  It is a concrete datum value, not a
  proposition, predicate, certificate, or structure saying that some datum
  exists.  In particular, there is no declaration, field, or helper proposition
  named
  $\operatorname{uses\_oracle\_reduction\_data}$,
  $\operatorname{exists\_oracle\_reduction\_data}$,
  $\operatorname{is\_oracle\_reduction\_data}$, or any synonym.  The
  implementation must not prove this node by introducing
  \[
    \exists R,\ \operatorname{def\_oracle\_reduction\_data}
      (P_{\operatorname{variant}},Q_{\operatorname{variant}})
  \]
  or by returning a wrapper whose only connection to
  Definition~\ref{def:oracle-reduction-data} is such an existential
  proposition.

  The returned datum has the following projection equations as part of this
  definition.  Its input map, size functions, output types, and correctness
  relations are exactly those of
  Definitions~\ref{def:promised-monochromatic-minip-source-specs} and
  \ref{def:rectangular-bichromatic-minip-oracle-specs}.  For an optimization
  source input $V$, the valid-input projection includes the promise $n\geq2$.
  For a threshold source input with $n<2$, the oracle-call-count projection is
  $r(V)=0$ and the final-output projection is $0$.  In every other valid case,
  the oracle-call-count projection is
  \[
    r(V)=\operatorname{length}(\operatorname{InternalIntervals}_n),
  \]
  the query projection is the total function satisfying, for every
  $1\leq s\leq r(V)$,
  \[
    q_s(V)=
      Q_{\operatorname{InternalIntervals}_n[s]}(V),
  \]
  and the query-size projections are
  \[
    N_Q(q_s(V))=
      N\!\left(Q_{\operatorname{InternalIntervals}_n[s]}(V)\right),
    \qquad
    L_Q(q_s(V))=\ell.
  \]
  The state-space, initial-state, query-construction, update, and finish
  projections of the returned datum are exactly the corresponding projections
  of
  $\operatorname{def\_minip\_bi\_to\_mono\_local\_phase\_presentation}
  (\operatorname{variant},V)$, not caller-supplied functions.  In the exact
  and $\gamma$ variants the final postprocessing projection is precisely
  \[
    \operatorname{def\_minip\_split\_finish\_pair}
      \bigl(V,\operatorname{OptHist}(V,H_{\rm raw})\bigr),
  \]
  where $H_{\rm raw}$ is the generic full oracle-answer history received by
  the oracle-reduction postprocess field and
  $\operatorname{OptHist}(V,H_{\rm raw})$ is the adapter of
  Definition~\ref{def:minip-oracle-history-to-full-optimization-history}.  In
  the threshold variant it is precisely
  $\operatorname{def\_minip\_split\_finish\_threshold}(H)$.  Local-to-global
  decoding of optimization answers is exactly
  $\operatorname{DecodePair}_I$ from
  Definition~\ref{def:minip-split-answer-local-to-global-decoding}.  These
  equations are eliminable by unfolding the returned datum itself; they are
  not separate assumptions and not partial constraints on an abstract $R$.

  This node therefore only assembles the concrete components into the oracle
  reduction datum required by Definition~\ref{def:oracle-reduction-data}.  It
  does not introduce new problem specifications, split-interval validity
  fields, local-to-global decoding fields, postprocessing fields, or
  local-phase compatibility fields.  All such mathematical content is supplied
  by the helper nodes in the \uses list, and the projections of the resulting
  datum are exactly their displayed equations.
\end{definition}

\begin{definition}[Exact split-history oracle-answer contract]
  \label{def:minip-exact-split-history-oracle-answer-contract}
  \uses{def:minip-bi-to-mono-reduction-datum, def:rectangular-bichromatic-minip-oracle-specs, def:minip-split-history-decoded-candidates}
  For a valid exact-optimization split-tree run on
  $V:\{1,\ldots,n\}\to\{0,1\}^{\ell}$, a full oracle-answer history
  \[
    H=((Q_{I_1}(V),z_1),\ldots,(Q_{I_m}(V),z_m))
  \]
  is exact-correct when every answer $z_s=(u_s,w_s)$ is a correct exact
  rectangular Min-IP answer for the actual query at the same position:
  \[
    \forall\,1\leq s\leq m,\quad
      \operatorname{exact\_correct}
        \bigl(Q_{I_s}(V),(u_s,w_s)\bigr).
  \]
  Equivalently, writing
  $A_s=\operatorname{SplitLeftRows}_{I_s}(V)$ and
  $B_s=\operatorname{SplitRightRows}_{I_s}(V)$, this condition is the concrete
  family of inequalities
  \[
    \forall\,1\leq s\leq m,\ \forall u,\ \forall w,\qquad
      \ip{A_s(u_s)}{B_s(w_s)}
      \leq
      \ip{A_s(u)}{B_s(w)}.
  \]
  The decoded candidate
  $\operatorname{CandPair}(V,H,s)$ and score
  $\operatorname{CandBeta}(V,H,s)$ are the functions of
  Definition~\ref{def:minip-split-history-decoded-candidates} applied to this
  same answer $z_s$.  This contract is not a global correctness tag on $H$:
  it exposes one exact optimality inequality for each concrete query-list
  position and each local cross-pair in that query.
\end{definition}

\begin{lemma}[Exact split answer is optimal after decoding]
  \label{lem:minip-exact-decoded-answer-cross-optimal}
  \uses{def:minip-exact-split-history-oracle-answer-contract, lem:minip-split-answer-row-preservation, def:minip-split-history-decoded-candidates}
  Let $H$ satisfy
  Definition~\ref{def:minip-exact-split-history-oracle-answer-contract}.  Fix
  a query position $s$ with interval $I_s$ and answer
  $z_s=(u_s,w_s)$.  For every local cross-pair $u,w$ in the same query,
  let
  \[
    (a_s,b_s)=\operatorname{CandPair}(V,H,s)
  \]
  be the decoded global pair, and let
  \[
    (p,q)=\operatorname{DecodePair}_{I_s}(u,w)
  \]
  be the decoded global pair for $(u,w)$.  Then
  \[
    \ip{v_{a_s}}{v_{b_s}}\leq \ip{v_p}{v_q}.
  \]
  In particular, if a distinct original pair $(p,q)$ occurs as a cross-pair in
  the query at position $s$, the decoded exact oracle answer at position $s$
  has score no larger than $\ip{v_p}{v_q}$.
\end{lemma}
\begin{proof}
  \uses{def:minip-exact-split-history-oracle-answer-contract, lem:minip-split-answer-row-preservation, def:minip-split-history-decoded-candidates}
  The exact-answer contract at position $s$ gives the local inequality
  \[
    \ip{A_s(u_s)}{B_s(w_s)}\leq\ip{A_s(u)}{B_s(w)}.
  \]
  Lemma~\ref{lem:minip-split-answer-row-preservation} rewrites the left side
  as the inner product of the decoded candidate
  $\operatorname{CandPair}(V,H,s)=(a_s,b_s)$ and rewrites the right side as the
  inner product of the decoded pair
  $\operatorname{DecodePair}_{I_s}(u,w)=(p,q)$.  Substitution gives the
  displayed global inequality.
\end{proof}

\begin{lemma}[Pair-valued finish scan selects a minimum candidate]
  \label{lem:minip-split-finish-pair-minimal}
  \uses{def:minip-split-finish-pair-scan, def:minip-split-history-first-min-beta, def:minip-split-history-decoded-candidates}
  For every nonempty optimization history $H$ of length $m>0$, let
  \[
    s^*=\operatorname{FirstMinBetaPos}(V,H)
  \]
  and let
  \[
    (a^*,b^*)=
      \operatorname{def\_minip\_split\_finish\_pair}(V,H)
      =
      \operatorname{CandPair}(V,H,s^*).
  \]
  Then for every history position $1\leq s\leq m$,
  \[
    \ip{v_{a^*}}{v_{b^*}}
    =
    \operatorname{CandBeta}(V,H,s^*)
    =
    \operatorname{MinBeta}(V,H)
    \leq
    \operatorname{CandBeta}(V,H,s).
  \]
  Thus the concrete finish function returns a decoded candidate whose score is
  no larger than the score of any decoded oracle-answer candidate in the same
  history.
\end{lemma}
\begin{proof}
  \uses{def:minip-split-finish-pair-scan, def:minip-split-history-first-min-beta, def:minip-split-history-decoded-candidates}
  By Definition~\ref{def:minip-split-history-first-min-beta},
  $s^*$ is the first position attaining the finite minimum
  $\operatorname{MinBeta}(V,H)$, so
  $\operatorname{CandBeta}(V,H,s^*)=\operatorname{MinBeta}(V,H)$ and this value
  is at most every $\operatorname{CandBeta}(V,H,s)$.  Definition~\ref{def:minip-split-finish-pair-scan}
  identifies the returned pair with
  $\operatorname{CandPair}(V,H,s^*)$, whose score is by definition
  $\operatorname{CandBeta}(V,H,s^*)$.
\end{proof}

\begin{lemma}[Exact recursive Min-IP answer combination]
  \label{lem:minip-recursive-exact-answer-combination}
  \uses{def:minip-bi-to-mono-reduction-datum, def:minip-exact-split-history-oracle-answer-contract, lem:minip-split-tree-pair-coverage, lem:minip-exact-decoded-answer-cross-optimal, lem:minip-split-finish-pair-minimal, def:min-ip}
  Let $V:\{1,\ldots,n\}\to\{0,1\}^{\ell}$ be a valid exact-optimization
  monochromatic Min-IP source input, so $n\geq2$.  Let $H$ be the full
  exact-optimization oracle-answer history generated by the concrete reduction
  datum of Definition~\ref{def:minip-bi-to-mono-reduction-datum}, in the same
  order as $\operatorname{InternalIntervals}_n$, and assume $H$ satisfies the
  exact split-history oracle-answer contract of
  Definition~\ref{def:minip-exact-split-history-oracle-answer-contract}.  Let
  \[
    (a^*,b^*)=
      \operatorname{def\_minip\_split\_finish\_pair}(V,H)
  \]
  be the exact variant's final postprocessing output.  Then $(a^*,b^*)$ is a
  valid distinct pair and
  \[
    \forall i,j,\ i\neq j\Rightarrow
      \ip{v_{a^*}}{v_{b^*}}\leq\ip{v_i}{v_j}.
  \]
  Equivalently, the exact split-tree reduction's deterministic finish scan
  returns a global minimizer for monochromatic $\MinIP$.
\end{lemma}
\begin{proof}
  \uses{def:minip-bi-to-mono-reduction-datum, def:minip-exact-split-history-oracle-answer-contract, lem:minip-split-tree-pair-coverage, lem:minip-exact-decoded-answer-cross-optimal, lem:minip-split-finish-pair-minimal, def:min-ip}
  Let $(p,q)$ be any distinct pair.  By
  Lemma~\ref{lem:minip-split-tree-pair-coverage}, it appears as a cross-pair
  in a unique concrete query-list position $s$ with interval $I_s$.  Let
  $(u,w)$ be the corresponding local indices in that query.  Lemma~\ref{lem:minip-exact-decoded-answer-cross-optimal}
  applied to the exact-answer contract for $H$ gives
  \[
    \operatorname{CandBeta}(V,H,s)\leq\ip{v_p}{v_q}.
  \]
  Since $n\geq2$, the split tree has at least one internal query, so the
  exact-optimization history is nonempty.  Lemma~\ref{lem:minip-split-finish-pair-minimal}
  says that the finish pair $(a^*,b^*)$ has score no larger than
  $\operatorname{CandBeta}(V,H,s)$.  Combining the two inequalities gives
  $\ip{v_{a^*}}{v_{b^*}}\leq\ip{v_p}{v_q}$.  The decoded candidate returned by
  the finish scan is a distinct pair because every
  $\operatorname{DecodePair}_{I_s}$ lands in the distinct-pair subtype by
  Definition~\ref{def:minip-split-answer-local-to-global-decoding}.  Since
  $(p,q)$ was arbitrary, the returned pair is a global monochromatic Min-IP
  minimizer.
\end{proof}

\begin{lemma}[Approximate recursive Min-IP answer combination]
  \label{lem:minip-recursive-approx-answer-combination}
  \uses{def:minip-bi-to-mono-reduction-datum, lem:minip-split-tree-pair-coverage, def:gamma-min-ip}
  Assume that every internal query $Q_I(V)$ is answered by a
  $\gamma$-approximate cross minimizer and let $(p^*,q^*)$ be a global
  monochromatic minimizer.  If $I_0$ is the unique internal node whose query
  contains $(p^*,q^*)$ as a cross-pair, then the optimum value of $Q_{I_0}(V)$
  equals the global optimum value.  Consequently the final postprocessing rule
  of Definition~\ref{def:minip-bi-to-mono-reduction-datum} returns a distinct
  pair whose inner product is at most $\gamma$ times the global optimum and at
  least the global optimum.
\end{lemma}
\begin{proof}
  \uses{def:minip-bi-to-mono-reduction-datum, lem:minip-split-tree-pair-coverage, def:gamma-min-ip}
  The query $Q_{I_0}(V)$ contains the global minimizer, so its cross optimum is
  at most the global optimum.  Since its feasible pairs are a subset of all
  distinct monochromatic pairs, its cross optimum is also at least the global
  optimum; hence the two optima are equal.  The oracle answer at $I_0$ has
  value at most $\gamma$ times this common optimum.  The final answer is the
  smallest value among all oracle-returned pairs, so it is no larger than the
  $I_0$ returned value.  Every returned pair is a valid monochromatic pair, so
  its value is at least the global optimum.
\end{proof}

\begin{lemma}[Threshold recursive Min-IP answer combination]
  \label{lem:minip-recursive-threshold-answer-combination}
  \uses{def:minip-bi-to-mono-reduction-datum, lem:minip-split-tree-pair-coverage, def:min-ip-decision}
  In the threshold variant of
  Definition~\ref{def:minip-bi-to-mono-reduction-datum}, the final Boolean OR
  of the internal-node oracle answers is $1$ if and only if the original
  monochromatic input has a distinct pair $i\neq j$ with
  $\ip{v_i}{v_j}\leq t$.
\end{lemma}
\begin{proof}
  \uses{def:minip-bi-to-mono-reduction-datum, lem:minip-split-tree-pair-coverage, def:min-ip-decision}
  If the final OR is $1$, then some internal query has a correct yes answer,
  so that query contains a cross-pair with inner product at most $t$; this is a
  distinct pair in the original input.  Conversely, if the original input has a
  pair $(p,q)$ with value at most $t$, then
  Lemma~\ref{lem:minip-split-tree-pair-coverage} places it as a cross-pair in
  a unique internal query.  The correct threshold answer for that query is
  $1$, so the final OR is $1$.
\end{proof}

\begin{lemma}[Recursive Min-IP query-size accounting]
  \label{lem:minip-recursive-query-size-accounting}
  \uses{def:minip-balanced-recursive-split-tree, def:bichromatic-min-ip}
  At every depth $h$ of the balanced split tree, the internal-node intervals
  are disjoint and have total size at most $n$.  Each query at depth $h$ has
  principal size at most the size of its parent interval.  Therefore, for any
  $\eps\in(0,1]$,
  \[
    \sum_{I\text{ internal at depth }h} N_I^{2-\eps}
    \leq
    n^{2-\eps}2^{-h(1-\eps)}
  \]
  up to an absolute constant accounting for the floor/ceiling imbalance.
\end{lemma}
\begin{proof}
  \uses{def:minip-balanced-recursive-split-tree, def:bichromatic-min-ip}
  Balanced splitting makes every depth-$h$ interval have size at most
  $\lceil n/2^h\rceil$, and the intervals at the same depth are disjoint.
  Since each query size $N_I$ is at most its parent interval size, the sum of
  $N_I^{2-\eps}$ over that depth is bounded by the number of nonempty intervals
  times $(\lceil n/2^h\rceil)^{2-\eps}$.  Absorbing the ceiling factor into an
  absolute constant gives the displayed geometric bound.
\end{proof}

\begin{lemma}[Recursive Min-IP local and oracle cost]
  \label{lem:minip-recursive-cost-transfer}
  \uses{def:minip-bi-to-mono-reduction-datum, lem:minip-recursive-query-size-accounting, def:subquadratic-oracle-query-budget, def:oracle-reduction-cost, def:truly-subquadratic}
  Suppose the rectangular bichromatic solver has running-time bound
  \[
    T_{\rm bi}(N,\ell)\leq C N^{2-\eps}(1+\ell)^b
  \]
  for some fixed $\eps>0$.  Then every locally generated trace of
  $R_{\rm biMinIP\to monoMinIP}[A_{\rm bi}]$ has total time
  \[
    L+\sum_{I\text{ internal in }\mathcal T(V)}T_{\rm bi}(N_I,\ell)
    \leq C'n^{2-\eps'}(1+\ell)^{b'}
  \]
  for constants $\eps'>0$, $C'>0$, and $b'\in\N$ depending only on
  $\eps,C,b$ and the fixed reduction, not on $V$ or on oracle answers.
\end{lemma}
\begin{proof}
  \uses{def:minip-bi-to-mono-reduction-datum, lem:minip-recursive-query-size-accounting, def:oracle-reduction-cost, def:truly-subquadratic}
  The local construction and postprocessing inspect or write only a constant
  number of records per tree node and at most $O(\ell)$ coordinates for each
  returned pair or copied row.  Across all depths this is bounded by
  $C_0n(1+\ell)\lceil\log_2(n+1)\rceil$ for an absolute constant $C_0$; in a
  view-based implementation the same bound holds a fortiori.  For oracle time,
  sum the assumed bound over depths and apply
  Lemma~\ref{lem:minip-recursive-query-size-accounting}.  The geometric series
  in $2^{-h(1-\eps)}$ is bounded by a constant depending only on $\eps$, giving
  $O(n^{2-\eps}(1+\ell)^b)$ oracle time when $0<\eps<1$; if $\eps\geq1$, reduce
  to any fixed smaller positive exponent.  The
  $n(1+\ell)\log(n+1)$ local term is absorbed by decreasing the exponent to
  some $\eps'>0$ and increasing $C'$ and $b'$.
\end{proof}

\begin{lemma}[Bichromatic Min-IP is harder than monochromatic Min-IP, Lemma A.8]
  \label{lem:minip-bi-to-mono}
  \uses{def:minip-bi-to-mono-reduction-datum, lem:minip-recursive-exact-answer-combination, lem:minip-recursive-approx-answer-combination, lem:minip-recursive-threshold-answer-combination, lem:minip-recursive-cost-transfer, def:fine-grained-reduction}
  For each of the exact, $\gamma$-approximate, and threshold Min-IP variants,
  any truly subquadratic rectangular bichromatic solver induces the concrete
  monochromatic solver
  $R_{\rm biMinIP\to monoMinIP}[A_{\rm bi}]$ of
  Definition~\ref{def:minip-bi-to-mono-reduction-datum}.  The induced solver's
  oracle queries are exactly the split-tree queries $Q_I(V)$, every distinct
  original pair is covered as a cross-pair in exactly one such query, the final
  answer is combined by the displayed minimum or Boolean-OR rule, and its total
  running time remains truly subquadratic up to polynomial factors in $\ell$.
\end{lemma}
\begin{proof}
  \uses{def:minip-bi-to-mono-reduction-datum, lem:minip-split-tree-pair-coverage, lem:minip-recursive-exact-answer-combination, lem:minip-recursive-approx-answer-combination, lem:minip-recursive-threshold-answer-combination, lem:minip-recursive-cost-transfer, def:fine-grained-reduction}
  Definition~\ref{def:minip-bi-to-mono-reduction-datum} gives the actual
  composed algorithm and its concrete query schedule.  Pair coverage is
  Lemma~\ref{lem:minip-split-tree-pair-coverage}.  The three answer-combination
  obligations are respectively
  Lemmas~\ref{lem:minip-recursive-exact-answer-combination},
  \ref{lem:minip-recursive-approx-answer-combination}, and
  \ref{lem:minip-recursive-threshold-answer-combination}.  The running-time
  transfer, including local tree construction and the sum of all oracle
  queries, is Lemma~\ref{lem:minip-recursive-cost-transfer}.  These are exactly
  the correctness and subquadratic-budget obligations required by
  Definition~\ref{def:fine-grained-reduction} for the concrete composed
  monochromatic algorithm.
\end{proof}

\begin{definition}[Maximum Inner Product, Definition A.9]
  \label{def:max-ip}
  Given binary vectors $v_1,\ldots,v_n\in\{0,1\}^{\ell}$,
  $\MaxIP_{n,\ell}$ asks to find a pair $i\neq j$ maximizing $\ip{v_i}{v_j}$.
\end{definition}

\begin{definition}[Approximate Maximum Inner Product, Definition A.10]
  \label{def:gamma-max-ip}
  \uses{def:max-ip}
  For $\gamma\geq 1$, $\gamma$-$\MaxIP_{n,\ell}$ asks for a pair whose inner
  product is a $\gamma$-approximation of the maximum inner product.
\end{definition}

\begin{definition}[Maximum Inner Product decision version, Definition A.11]
  \label{def:max-ip-decision}
  \uses{def:max-ip}
  Given a threshold $t$, $\MaxIP_{n,\ell,t}$ asks whether some pair $i\neq j$
  has $\ip{v_i}{v_j}\geq t$.
\end{definition}

\begin{definition}[Bichromatic Maximum Inner Product]
  \label{def:bichromatic-max-ip}
  \uses{def:max-ip}
  Given $A,B\subseteq\{0,1\}^{\ell}$ of size $n$, the bichromatic versions of
  $\MaxIP$, $\gamma$-$\MaxIP$, and $\MaxIP_{n,\ell,t}$ optimize, approximate, or
  decide over pairs in $A\times B$.
\end{definition}

\begin{definition}[Bichromatic additive Maximum Inner Product]
  \label{def:additive-bichromatic-max-ip}
  \uses{def:bichromatic-max-ip}
  Given sets $A,B\subseteq\{0,1\}^{\ell}$ of size $n$, an integer threshold
  $\alpha$, and an additive gap $g>0$, bichromatic $g$-Additive-$\MaxIP$
  distinguishes the following promised cases:
  \[
    \text{yes: }\exists (a,b)\in A\times B,\ \ip{a}{b}\geq\alpha,
  \]
  and
  \[
    \text{no: }\forall (a,b)\in A\times B,\ \ip{a}{b}<\alpha-g.
  \]
\end{definition}

\begin{lemma}[Bichromatic Max-IP is harder than monochromatic Max-IP, Lemma A.12]
  \label{lem:maxip-bi-to-mono}
  \uses{def:max-ip, def:gamma-max-ip, def:max-ip-decision, def:bichromatic-max-ip}
  If any bichromatic Max-IP variant has a truly subquadratic algorithm, then the
  corresponding monochromatic variant has a truly subquadratic algorithm, up to
  polynomial factors in $\ell$.
\end{lemma}
\begin{proof}
  \uses{lem:minip-bi-to-mono, def:bichromatic-max-ip}
  Use the same divide-and-recurse enumeration as in
  Lemma~\ref{lem:minip-bi-to-mono}.  Each unordered pair of input vectors is
  presented to exactly one bichromatic subproblem, so taking the best answer
  among all recursion nodes gives the monochromatic optimum, approximation, or
  threshold decision.
\end{proof}

\begin{lemma}[Bichromatic Max-IP is OV-hard, Lemma A.13]
  \label{lem:bichromatic-maxip-hard-ov}
  \uses{def:ov, def:bichromatic-max-ip, lem:binary-complement-inner}
  A truly subquadratic algorithm for bichromatic $\MaxIP_{n,\ell}$ implies a
  truly subquadratic algorithm for $\OV_{n,\ell}$, up to polynomial factors in
  $\ell$.
\end{lemma}
\begin{proof}
  \uses{def:ov, def:bichromatic-max-ip, lem:binary-complement-inner}
  Partition the OV input vectors into $S_0,\ldots,S_\ell$, where $S_i$ contains
  the vectors with exactly $i$ ones.  For each pair $(i,j)$, form
  $\bar S_j$ by complementing every vector in $S_j$ and run bichromatic Max-IP on
  $S_i,\bar S_j$.  For $v\in S_i$ and $w\in S_j$,
  Lemma~\ref{lem:binary-complement-inner} gives
  $\ip{v}{\bar w}=i-\ip{v}{w}$.  Hence $\ip{v}{w}=0$ exactly when
  $\ip{v}{\bar w}=i$, the largest possible value for a vector in $S_i$.  Trying
  all $\ell^2$ pairs of weight classes therefore decides whether any
  orthogonal pair exists.
\end{proof}

\begin{theorem}[Approximate Max-IP hardness from Karthik--Manurangsi]
  \label{thm:km-approx-maxip-hardness}
  \uses{def:gamma-max-ip, def:ov, conj:ovc}
  \notready
  Approximate $\MaxIP$ in the parameter ranges used by Appendix B is hard under
  OVC/SETH.
\end{theorem}
% TODO: The paper invokes Karthik--Manurangsi (2020) for this stronger
% approximate Max-IP hardness and does not reproduce its proof.

\begin{definition}[Most Similar Documents, Definition 2.10]
  \label{def:msd}
  \uses{def:cosine-similarity}
  Given nonzero binary document embeddings $v_1,\ldots,v_n\in\{0,1\}^{\ell}$,
  $\MSD_{n,\ell}$ asks to find a pair $i\neq j$ maximizing
  $\ip{v_i}{v_j}/(\|v_i\|\|v_j\|)$.
\end{definition}

\begin{definition}[Approximate Most Similar Documents, Definition 2.11]
  \label{def:gamma-msd}
  \uses{def:msd}
  For $\gamma\geq 1$, $\gamma$-$\MSD_{n,\ell}$ asks to find a pair
  $(i^*,j^*)$ whose cosine similarity is between $1/\gamma$ times the optimum
  and the optimum.
\end{definition}

\begin{definition}[Most Similar Documents decision version, Definition 2.12]
  \label{def:msd-decision}
  \uses{def:msd}
  Given $t\in[0,1]$, $\MSD_{n,\ell,t}$ asks whether some pair $i\neq j$ has
  cosine similarity at least $t$.
\end{definition}

\begin{definition}[Bichromatic Most Similar Documents]
  \label{def:bichromatic-msd}
  \uses{def:msd}
  Bichromatic $\MSD$ and its approximate and threshold variants are defined over
  pairs drawn from two input sets $A,B\subseteq\{0,1\}^{\ell}$.
\end{definition}

\begin{lemma}[Decision MSD suffices for exact MSD]
  \label{lem:msd-decision-search}
  \uses{def:msd, def:msd-decision}
  If $\MSD_{n,\ell,t}$ is truly subquadratic for every threshold $t$, then
  $\MSD_{n,\ell}$ is truly subquadratic.
\end{lemma}
\begin{proof}
  \uses{def:msd, def:msd-decision}
  For binary vectors, the value $\ip{v_i}{v_j}/(\|v_i\|\|v_j\|)$ is determined by
  three integers: $\ip{v_i}{v_j}$, $\|v_i\|_1$, and $\|v_j\|_1$, each between
  $0$ and $\ell$.  Hence there are only $O(\ell^3)$ possible similarity values.
  Binary search or enumeration over this discrete set using the decision oracle
  finds the optimum with only a polynomial-in-$\ell$ overhead.
\end{proof}

\begin{definition}[Least Similar Documents, Definition A.14]
  \label{def:lsd}
  \uses{def:cosine-similarity}
  Given nonzero binary document embeddings $v_1,\ldots,v_n\in\{0,1\}^{\ell}$,
  $\LSD_{n,\ell}$ asks to find a pair $i\neq j$ minimizing
  $\ip{v_i}{v_j}/(\|v_i\|\|v_j\|)$.
\end{definition}

\begin{definition}[Approximate Least Similar Documents, Definition A.15]
  \label{def:gamma-lsd}
  \uses{def:lsd}
  For $\gamma\geq 1$, $\gamma$-$\LSD_{n,\ell}$ asks to find a pair
  $(i^*,j^*)$ whose cosine similarity is at least the optimum and at most
  $\gamma$ times the optimum.
\end{definition}

\begin{definition}[Least Similar Documents decision version, Definition A.16]
  \label{def:lsd-decision}
  \uses{def:lsd}
  Given $t\in[0,1]$, $\LSD_{n,\ell,t}$ asks whether some pair $i\neq j$ has
  cosine similarity at most $t$.
\end{definition}

\begin{definition}[Bichromatic Least Similar Documents]
  \label{def:bichromatic-lsd}
  \uses{def:lsd}
  Bichromatic $\LSD$ and its approximate and threshold variants are defined over
  pairs drawn from two input sets $A,B\subseteq\{0,1\}^{\ell}$.
\end{definition}

\begin{lemma}[Decision LSD suffices for exact LSD]
  \label{lem:lsd-decision-search}
  \uses{def:lsd, def:lsd-decision}
  If $\LSD_{n,\ell,t}$ is truly subquadratic for every threshold $t$, then
  $\LSD_{n,\ell}$ is truly subquadratic.
\end{lemma}
\begin{proof}
  \uses{def:lsd, def:lsd-decision, lem:msd-decision-search}
  The same $O(\ell^3)$ discretization used in
  Lemma~\ref{lem:msd-decision-search} applies to cosine similarities for LSD.
  Searching those threshold values with an LSD decision oracle returns the
  minimum.
\end{proof}

\begin{lemma}[Bichromatic similarity variants are harder]
  \label{lem:bichromatic-similarity-harder}
  \uses{def:bichromatic-msd, def:bichromatic-lsd, lem:minip-bi-to-mono}
  Truly subquadratic algorithms for bichromatic MSD or LSD variants imply truly
  subquadratic algorithms for the corresponding monochromatic variants.
\end{lemma}
\begin{proof}
  \uses{lem:minip-bi-to-mono, def:bichromatic-msd, def:bichromatic-lsd}
  The divide-and-recurse reduction of Lemma~\ref{lem:minip-bi-to-mono} depends
  only on the fact that every unordered pair appears as a cross-pair in exactly
  one recursive subproblem.  It therefore applies verbatim to cosine-similarity
  maximization, minimization, approximation, and threshold decision.
\end{proof}

\chapter{Hardness of document similarity}

\begin{lemma}[Zero cosine is exactly orthogonality]
  \label{lem:lsd-zero-equivalence}
  \uses{def:ov, def:cosine-similarity}
  For nonzero binary vectors $v,w$, their cosine similarity is $0$ if and only
  if $\ip{v}{w}=0$.
\end{lemma}
\begin{proof}
  \uses{def:cosine-similarity, def:inner-norm}
  Since $v$ and $w$ are nonzero, $\|v\|\|w\|>0$.  Dividing by this positive
  denominator preserves whether the numerator is zero, so
  $\ip{v}{w}/(\|v\|\|w\|)=0$ exactly when $\ip{v}{w}=0$.
\end{proof}

\begin{theorem}[LSD hardness, Theorem 3.1]
  \label{thm:lsd-hardness}
  \uses{conj:seth, conj:ovc, def:gamma-lsd, def:ov}
  Assuming SETH or OVC, for every $\eps>0$ there exists a constant $c>0$ such
  that for every $\gamma\geq1$, $\gamma$-$\LSD_{n,\ell}$ cannot be solved in
  $O(n^{2-\eps})$ time when $\ell=c\log n$.
\end{theorem}
\begin{proof}
  \uses{lem:lsd-zero-equivalence, conj:ovc, thm:seth-implies-ovc}
  Suppose, toward contradiction, that an algorithm $A$ solves
  $\gamma$-$\LSD_{n,c\log n}$ in $O(n^{2-\eps})$ time for the constants needed
  by OVC.  Given an OV instance, first check in $O(n\ell)$ time whether any
  vector is zero; if so, answer yes.  Otherwise all vectors are nonzero and the
  LSD objective is well-defined.  Run $A$ and inspect the returned pair.  By
  Lemma~\ref{lem:lsd-zero-equivalence}, the optimum LSD value is $0$ exactly
  when the OV instance has an orthogonal pair.  Any multiplicative
  $\gamma$-approximation to an optimum of $0$ must also return value $0$, while
  if no orthogonal pair exists every cosine similarity is positive.  Thus this
  decides $\OV_{n,c\log n}$ in $O(n\ell+n^{2-\eps})=O(n^{2-\eps'})$ time,
  contradicting OVC, and hence SETH through Theorem~\ref{thm:seth-implies-ovc}.
\end{proof}

\begin{corollary}[LSD variant hardness, Corollary 3.2]
  \label{cor:lsd-hardness-variants}
  \uses{thm:lsd-hardness, lem:lsd-decision-search, lem:bichromatic-similarity-harder}
  Assuming SETH or OVC, for every $\eps>0$ there exists $c>0$ such that
  $\LSD_{n,\ell}$ cannot be solved in $O(n^{2-\eps})$ time for $\ell\geq
  c\log n$.  The same lower bound holds for bichromatic $\gamma$-$\LSD$ for all
  $\gamma\geq1$ and for $\LSD_{n,\ell,t}$ for some $t\in[0,1]$.
\end{corollary}
\begin{proof}
  \uses{thm:lsd-hardness, lem:lsd-decision-search, lem:bichromatic-similarity-harder}
  Exact LSD is at least as hard as approximate LSD because an exact optimum is a
  valid approximation.  Bichromatic LSD is harder than monochromatic LSD by
  Lemma~\ref{lem:bichromatic-similarity-harder}.  If every threshold problem
  $\LSD_{n,\ell,t}$ were truly subquadratic, then
  Lemma~\ref{lem:lsd-decision-search} would give a truly subquadratic exact LSD
  algorithm, contradicting Theorem~\ref{thm:lsd-hardness}.  Therefore at least
  one threshold $t$ is hard.
\end{proof}

\begin{definition}[Bichromatic additive MSD]
  \label{def:additive-bichromatic-msd}
  \uses{def:bichromatic-msd}
  Given $A,B\subseteq\{0,1\}^{\ell}$, $\alpha\in[0,1]$, and an additive gap
  $\delta>0$, bichromatic $\delta$-Additive-MSD distinguishes yes instances,
  where some $(a,b)\in A\times B$ has cosine similarity at least $\alpha$, from
  no instances, where every $(a,b)$ has cosine similarity less than
  $\alpha-\delta$.
\end{definition}

\begin{lemma}[Kronecker powers amplify cosine similarity]
  \label{lem:kron-cosine-amplification}
  \uses{lem:kron-inner, lem:kron-norm, def:cosine-similarity}
  For nonzero vectors $v,w$ and positive integer $q$,
  \[
    \frac{\ip{v^{\otimes q}}{w^{\otimes q}}}
         {\|v^{\otimes q}\|\|w^{\otimes q}\|}
    =
    \left(\frac{\ip{v}{w}}{\|v\|\|w\|}\right)^q .
  \]
\end{lemma}
\begin{proof}
  \uses{lem:kron-inner, lem:kron-norm, def:cosine-similarity}
  Substitute Lemma~\ref{lem:kron-inner} into the numerator and
  Lemma~\ref{lem:kron-norm} into both denominator factors.  The resulting ratio
  is $\ip{v}{w}^q/(\|v\|^q\|w\|^q)$, which is the $q$th power of the original
  cosine similarity.
\end{proof}

\begin{lemma}[Tensor boosting for approximate MSD, Lemma B.2]
  \label{lem:msd-tensor-boost}
  \uses{def:gamma-msd, lem:kron-cosine-amplification}
  If $\gamma$-$\MSD_{n,\ell}$ in dimension
  $\ell=(\log n)^{c\log n/(\log\log n)^2}$ is solvable in $O(n^{2-\eps})$ time
  for all constants $c>0$ and
  $\gamma\leq(1+1/\log\log n)^{\log n/(\log\log n)^2}$, then
  $(1+1/\log\log n)$-$\MSD_{n,(\log n)^k}$ is solvable in truly subquadratic
  time for every constant $k>0$.
\end{lemma}
\begin{proof}
  \uses{lem:kron-cosine-amplification, def:gamma-msd}
  Given vectors in dimension $(\log n)^k$, set
  $q=\log n/(\log\log n)^2$ and replace every vector $v_i$ by
  $v_i^{\otimes q}$.  The new dimension is
  $(\log n)^{k\log n/(\log\log n)^2}$, which has the form required by the
  assumed algorithm.  By Lemma~\ref{lem:kron-cosine-amplification}, every cosine
  similarity is raised to the $q$th power.  Therefore a $\gamma$-approximation
  after tensoring becomes a $\gamma^{1/q}$-approximation before tensoring, and
  the stated bound on $\gamma$ makes this at most $1+1/\log\log n$.
\end{proof}

\begin{theorem}[Karthik--Manurangsi reduction to additive Max-IP, Theorem B.4]
  \label{thm:km-msd-to-additive-maxip}
  \uses{def:gamma-msd, def:additive-bichromatic-msd, def:additive-bichromatic-max-ip}
  \notready
  If $(1+1/\log\log n)$-$\MSD_{n,(\log n)^k}$ is truly subquadratic for every
  constant $k>0$, then bichromatic $(\log n)$-Additive-Max-IP in dimension
  $c\log n$ is truly subquadratic for every constant $c>0$.
\end{theorem}
% TODO: The paper quotes this as Theorem 6.2 of Karthik--Manurangsi (2020)
% rather than proving the graph construction.

\begin{lemma}[Equal-weight padding preserves cross inner products]
  \label{lem:equal-weight-padding-inner}
  \uses{def:binary-vector-ops, def:inner-norm, def:kron-concat}
  Let $a,b\in\{0,1\}^{\ell}$.  Define
  \[
    a'=(a,1^{\ell-\|a\|_1},0^{\ell+\|a\|_1}),\qquad
    b'=(b,0^{\ell+\|b\|_1},1^{\ell-\|b\|_1}).
  \]
  Then $a',b'\in\{0,1\}^{3\ell}$, each has exactly $\ell$ ones, and
  $\ip{a'}{b'}=\ip{a}{b}$.
\end{lemma}
\begin{proof}
  \uses{def:binary-vector-ops, def:inner-norm, def:kron-concat}
  The appended block of $a'$ containing ones is paired with a zero block of
  $b'$, and the appended block of $b'$ containing ones is paired with a zero
  block of $a'$.  Hence all appended coordinates contribute zero to the cross
  inner product, while the first $\ell$ coordinates contribute exactly
  $\ip{a}{b}$.  The number of ones in $a'$ is
  $\|a\|_1+(\ell-\|a\|_1)=\ell$, and the same calculation gives $\ell$ ones
  for $b'$.
\end{proof}

\begin{lemma}[Equal-weight padding converts inner product gaps to cosine gaps]
  \label{lem:equal-weight-padding-cosine}
  \uses{lem:equal-weight-padding-inner, def:cosine-similarity}
  With $a'$ and $b'$ as in Lemma~\ref{lem:equal-weight-padding-inner},
  \[
    \frac{\ip{a'}{b'}}{\|a'\|\|b'\|}=\frac{\ip{a}{b}}{\ell}.
  \]
  Consequently an additive Max-IP threshold $\alpha$ and gap $g$ become the
  additive MSD threshold $\alpha/\ell$ and gap $g/\ell$.
\end{lemma}
\begin{proof}
  \uses{lem:equal-weight-padding-inner, def:cosine-similarity, def:inner-norm}
  By Lemma~\ref{lem:equal-weight-padding-inner}, both padded vectors have
  exactly $\ell$ binary one-coordinates, so their Euclidean norms are
  $\sqrt{\ell}$.  The same lemma gives $\ip{a'}{b'}=\ip{a}{b}$, and substituting
  these identities into the cosine-similarity formula gives
  $\ip{a}{b}/(\sqrt{\ell}\sqrt{\ell})=\ip{a}{b}/\ell$.  Scaling both the
  threshold and the additive gap by this common factor preserves the promised
  yes/no separation.
\end{proof}

\begin{lemma}[Additive MSD and additive Max-IP equivalence, Lemma B.5]
  \label{lem:additive-msd-maxip-equivalence}
  \uses{def:additive-bichromatic-msd, def:additive-bichromatic-max-ip, lem:equal-weight-padding-inner, lem:equal-weight-padding-cosine}
  Bichromatic $(\log n)/\ell$-Additive-MSD in dimension $\ell=O(\log n)$ and
  bichromatic $(\log n)$-Additive-Max-IP in dimension $O(\log n)$ are
  subquadratic equivalent.
\end{lemma}
\begin{proof}
  \uses{def:additive-bichromatic-msd, def:additive-bichromatic-max-ip, lem:equal-weight-padding-inner, lem:equal-weight-padding-cosine}
  Given a bichromatic additive Max-IP instance $A,B\subseteq\{0,1\}^{\ell}$ with
  threshold $\alpha$, replace each $a\in A$ by
  $a'=(a,1^{\ell-\|a\|_1},0^{\ell+\|a\|_1})$ and each $b\in B$ by
  $b'=(b,0^{\ell+\|b\|_1},1^{\ell-\|b\|_1})$.  Then every new vector has exactly
  $\ell$ ones, and Lemma~\ref{lem:equal-weight-padding-cosine} says that the
  cosine similarity of $a',b'$ is $\ip{a}{b}/\ell$.  Thus the additive Max-IP
  threshold $\alpha$ becomes the additive MSD threshold $\alpha/\ell$, with the
  additive gap scaled by $1/\ell$.

  Conversely, given an additive MSD instance, apply the same padding so all new
  vectors have equal norm $\sqrt{\ell}$ while preserving cross inner products.
  A cosine threshold $\alpha$ becomes the Max-IP threshold $\alpha\ell$, and an
  additive cosine gap of $(\log n)/\ell$ becomes an additive inner-product gap
  of $\log n$.  The dimension only triples.
\end{proof}

\begin{theorem}[Chen additive Max-IP hardness, Lemma B.6]
  \label{thm:chen-additive-maxip-hard}
  \uses{def:additive-bichromatic-max-ip, def:ov, conj:ovc}
  \notready
  If bichromatic $(\log n)$-Additive-Max-IP in dimension $c\log n$ is truly
  subquadratic for every constant $c>0$, then $\OV_{n,c'\log n}$ is truly
  subquadratic for every constant $c'>0$, contradicting OVC and SETH.
\end{theorem}
% TODO: The paper cites Chen (2020) for this additive Max-IP lower bound.

\begin{lemma}[MSD hardness chain]
  \label{lem:msd-hardness-chain}
  \uses{lem:msd-tensor-boost, thm:km-msd-to-additive-maxip, lem:additive-msd-maxip-equivalence, thm:chen-additive-maxip-hard}
  A truly subquadratic algorithm for the approximate MSD regime of
  Theorem~\ref{thm:msd-hardness} would imply a truly subquadratic OV algorithm
  in logarithmic dimension.
\end{lemma}
\begin{proof}
  \uses{lem:msd-tensor-boost, thm:km-msd-to-additive-maxip, lem:additive-msd-maxip-equivalence, thm:chen-additive-maxip-hard}
  Lemma~\ref{lem:msd-tensor-boost} converts the assumed high-dimensional
  approximate MSD algorithm into a truly subquadratic
  $(1+1/\log\log n)$-approximate MSD algorithm in every polylogarithmic
  dimension.  Theorem~\ref{thm:km-msd-to-additive-maxip} then yields a truly
  subquadratic additive Max-IP algorithm.  Lemma~\ref{lem:additive-msd-maxip-equivalence}
  transfers between the additive MSD and additive Max-IP formulations with only
  logarithmic dimension blowup.  Finally
  Theorem~\ref{thm:chen-additive-maxip-hard} converts this into a truly
  subquadratic OV algorithm, contradicting OVC.
\end{proof}

\begin{theorem}[MSD hardness, Theorem 3.3/B.1]
  \label{thm:msd-hardness}
  \uses{conj:seth, conj:ovc, def:gamma-msd}
  Assuming SETH or OVC, for every $\eps>0$ there is a constant $c>0$ such that
  $\gamma$-$\MSD_{n,\ell}$ cannot be solved in $O(n^{2-\eps})$ time when
  \[
    \ell \geq (\log n)^{c\log n/(\log\log n)^2}
    \quad\text{and}\quad
    \gamma \leq \left(1+\frac1{\log\log n}\right)^{\log n/(\log\log n)^2}.
  \]
\end{theorem}
\begin{proof}
  \uses{lem:msd-hardness-chain, conj:ovc, thm:seth-implies-ovc}
  If such an MSD algorithm existed for all constants required by OVC, then
  Lemma~\ref{lem:msd-hardness-chain} would imply a truly subquadratic algorithm
  for $\OV_{n,c'\log n}$ for every constant $c'>0$.  This contradicts OVC, and
  contradicts SETH through Theorem~\ref{thm:seth-implies-ovc}.
\end{proof}

\begin{corollary}[MSD variant hardness, Corollary 3.4]
  \label{cor:msd-hardness-variants}
  \uses{thm:msd-hardness, lem:msd-decision-search, lem:bichromatic-similarity-harder}
  Assuming SETH or OVC, for every $\eps>0$ there is a constant $c>0$ such that
  $\MSD_{n,\ell}$ cannot be solved in $O(n^{2-\eps})$ time when
  $\ell\geq(\log n)^{c\log n/(\log\log n)^2}$.  The same lower bound holds for
  bichromatic $\gamma$-$\MSD$ in the same approximation range and for
  $\MSD_{n,\ell,t}$ for some $t\in[0,1]$.
\end{corollary}
\begin{proof}
  \uses{thm:msd-hardness, lem:msd-decision-search, lem:bichromatic-similarity-harder}
  Exact MSD is at least as hard as approximate MSD.  Bichromatic variants are
  harder by Lemma~\ref{lem:bichromatic-similarity-harder}.  If every decision
  threshold were truly subquadratic, Lemma~\ref{lem:msd-decision-search} would
  make exact MSD truly subquadratic, contradicting
  Theorem~\ref{thm:msd-hardness}.
\end{proof}

\begin{theorem}[Introductory LSD lower bound, Theorem 1.1]
  \label{thm:intro-lsd-hardness}
  \uses{thm:lsd-hardness, cor:lsd-hardness-variants}
  The least-similar-document problem and its listed variants require quadratic
  time under SETH in the logarithmic-dimensional regime stated in the
  introduction.
\end{theorem}
\begin{proof}
  \uses{thm:lsd-hardness, cor:lsd-hardness-variants}
  The exact statement is the combination of Theorem~\ref{thm:lsd-hardness} for
  approximate LSD and Corollary~\ref{cor:lsd-hardness-variants} for exact,
  threshold, and bichromatic variants.
\end{proof}

\begin{theorem}[Introductory MSD lower bound, Theorem 1.2]
  \label{thm:intro-msd-hardness}
  \uses{thm:msd-hardness, cor:msd-hardness-variants}
  The most-similar-document problem and its threshold and approximation variants
  require quadratic time under SETH in the quasipolylogarithmic-dimensional
  regime stated in the introduction.
\end{theorem}
\begin{proof}
  \uses{thm:msd-hardness, cor:msd-hardness-variants}
  This is precisely Theorem~\ref{thm:msd-hardness} together with
  Corollary~\ref{cor:msd-hardness-variants}.
\end{proof}

\begin{theorem}[Bichromatic MSD logarithmic lower bound, Theorem 1.3]
  \label{thm:intro-bichromatic-msd-hardness}
  \uses{lem:bichromatic-maxip-hard-ov, lem:additive-msd-maxip-equivalence, thm:chen-additive-maxip-hard}
  Assuming SETH, for every $\eps>0$ there exists $c>0$ such that bichromatic
  $\MSD_{n,\ell}$ cannot be solved in $O(n^{2-\eps})$ time when $\ell=c\log n$.
\end{theorem}
\begin{proof}
  \uses{lem:bichromatic-maxip-hard-ov, lem:additive-msd-maxip-equivalence, thm:chen-additive-maxip-hard}
  The bichromatic setting avoids the monochromatic tensor-power loss.  The
  paper routes logarithmic-dimensional bichromatic MSD through the corresponding
  bichromatic Max-IP hardness: equal-norm padding preserves cross inner products
  as cosine similarities, and the cited additive/Max-IP hardness supplies the
  OV contradiction in logarithmic dimension.
\end{proof}

\chapter{Representational strength of transformers}

\begin{lemma}[Continuous threshold OR, Lemma D.1]
  \label{lem:continuous-threshold-or}
  \uses{def:mlp}
  For $a<b$ there is a continuous function $f:\R^\ell\to\R$ such that
  $f(x)=1$ if some coordinate $x[i]\geq b$, and $f(x)=0$ if all coordinates
  satisfy $x[i]<a$.
\end{lemma}
\begin{proof}
  \uses{def:mlp}
  Define
  \[
    g(y)=
    \begin{cases}
      1, & y\geq b,\\
      (y-a)/(b-a), & a\leq y<b,\\
      0, & y<a.
    \end{cases}
  \]
  This piecewise-linear function is continuous.  Set
  $f(x)=1-\prod_{i=1}^{\ell}(1-g(x[i]))$.  If some $x[i]\geq b$, then one
  factor is zero and $f(x)=1$.  If every $x[i]<a$, every $g(x[i])$ is zero and
  $f(x)=0$.
\end{proof}

\begin{lemma}[Finite row prescriptions can be implemented by an MLP]
  \label{lem:finite-row-prescription-mlp}
  \uses{def:mlp, def:finite-real-arrays}
  Let $S\subset\R^a$ be a finite set of valid row inputs, and assign a target
  vector $y_s\in\R^b$ to every $s\in S$.  There exists a continuous function
  $f:\R^a\to\R^b$ such that $f(s)=y_s$ for every $s\in S$.
\end{lemma}
\begin{proof}
  \uses{def:mlp, def:finite-real-arrays}
  Because $S$ is finite, choose pairwise disjoint open balls centered at the
  points of $S$.  On the ball around $s$, use a continuous bump function that
  is $1$ at $s$ and $0$ on the boundary, multiply it by $y_s$, and set the
  function to $0$ outside the union of the balls.  The balls are disjoint, so
  these local definitions agree on overlaps, and the resulting piecewise
  function is continuous and has the prescribed values on $S$.
\end{proof}

\begin{lemma}[MLP appending an end token, Lemma D.2]
  \label{lem:append-end-token-mlp}
  \uses{def:mlp, def:kron-concat, lem:finite-row-prescription-mlp}
  Fix an integer threshold $1\leq t\leq\ell$.  Let the raw sentinel row be
  \[
    s_t=(0,\ldots,0,t+1)\in\R^\ell,
  \]
  and let the target attention-space sentinel row be
  \[
    \tau_t=(0,\ldots,0,t+1)\in\R^{\ell+1},
  \]
  with $\ell$ leading zero coordinates and last coordinate $t+1$.  On the
  finite valid row set consisting of all binary rows in $\{0,1\}^{\ell}$ and
  the sentinel row $s_t$, there exists a continuous
  $f:\R^\ell\to\R^{\ell+1}$ such that
  \[
    f(x)=(x,1)\quad\text{for every }x\in\{0,1\}^{\ell},
    \qquad
    f(s_t)=\tau_t.
  \]
\end{lemma}
\begin{proof}
  \uses{lem:finite-row-prescription-mlp}
  Since $t\geq1$, the sentinel row has last coordinate at least $2$ and is not
  a binary row.  The valid row set is therefore finite with no conflicting
  prescriptions.  Assign $(x,1)$ to each binary row $x$ and assign $\tau_t$ to
  $s_t$, then apply Lemma~\ref{lem:finite-row-prescription-mlp}.
\end{proof}

\begin{lemma}[MLP normalizing and appending an end token, Lemma D.3]
  \label{lem:normalize-end-token-mlp}
  \uses{def:mlp, def:kron-concat, def:inner-norm, lem:finite-row-prescription-mlp}
  Fix a threshold $t\in[0,1]$.  Let the raw sentinel row be
  \[
    s_t=(0,\ldots,0,t+2)\in\R^\ell,
  \]
  and let the target attention-space sentinel row be
  \[
    \tau_t=(0,\ldots,0,t+1)\in\R^{\ell+1},
  \]
  with $\ell$ leading zero coordinates and last coordinate $t+1$.  On the
  finite set consisting of all nonzero binary document rows and $s_t$, there
  exists a continuous $f:\R^\ell\to\R^{\ell+1}$ such that
  \[
    f(x)=(x/\|x\|,1)\quad\text{for every nonzero }x\in\{0,1\}^{\ell},
  \]
  and
  \[
    f(s_t)=\tau_t.
  \]
\end{lemma}
\begin{proof}
  \uses{lem:finite-row-prescription-mlp, def:inner-norm}
  The set of valid rows is finite: it contains the $2^\ell-1$ nonzero binary
  vectors and the one special vector $s_t$.  Since $t+2>1$, the sentinel row is
  not binary, so there is no conflict with the document-row prescriptions.
  Prescribe the displayed value for each valid row and apply
  Lemma~\ref{lem:finite-row-prescription-mlp}.  The document-row prescription
  uses Euclidean normalization, so the inner product of two prescribed document
  rows is exactly their cosine similarity.
\end{proof}

\begin{lemma}[Low-rank attention score matrices can realize scalar identities]
  \label{lem:attention-factorization}
  \uses{def:attention}
  If $m\geq d_{\rm in}$, then for any scalar $\lambda$ there exist
  $Q,K\in\R^{d_{\rm in}\times m}$ with $QK^\top=\lambda I_{d_{\rm in}}$.
\end{lemma}
\begin{proof}
  \uses{def:attention}
  Put $\lambda$ on the first $d_{\rm in}$ diagonal entries of $Q$ and put $1$
  on the corresponding diagonal entries of $K$, with all other entries zero.
  Then the matrix product has diagonal entries $\lambda$ and off-diagonal
  entries zero.
\end{proof}

\begin{lemma}[OV attention output formula]
  \label{lem:ov-attention-output}
  \uses{def:attention, lem:attention-factorization, def:ov, def:sentinel-augmented-input}
  In the OV construction of Theorem~\ref{thm:transformer-solves-ov}, with
  raw sentinel row $s=0^\ell$, augmented attention input
  $\Aug_s(E)\in\R^{(n+1)\times\ell}$, $QK^\top=-\lambda I_\ell$ for a concrete
  scale $\lambda>0$, and $V$ the all-one column, the $i$th attention output for
  every document row
  $1\leq i\leq n$ equals
  \[
    \frac{\sum_{j=1}^{n}\exp(-\lambda\ip{v_i}{v_j})\|v_j\|_1}
         {\sum_{k=1}^{n}\exp(-\lambda\ip{v_i}{v_k})+1}.
  \]
\end{lemma}
\begin{proof}
  \uses{def:attention, def:softmax-matrix, def:inner-norm, def:sentinel-augmented-input}
  The attention score from row $i$ to row $j$ is
  $-\lambda\ip{v_i}{v_j}$, so exponentiating it inside softmax gives
  $\exp(-\lambda\ip{v_i}{v_j})$.
  Definition~\ref{def:sentinel-augmented-input} makes the last row exactly
  $0^\ell$, so it contributes score $0$ and hence denominator contribution
  $1$.  Multiplication by the all-one value column returns $\|v_j\|_1$ for
  document rows and $0$ for the sentinel row, giving the displayed formula for
  the first $n$ rows.
\end{proof}

\begin{lemma}[OV yes-instance attention gap]
  \label{lem:ov-yes-gap}
  \uses{lem:ov-attention-output, def:ov}
  If the OV instance has an orthogonal pair, then some attention output entry in
  the OV construction is at least $1/(n+1)$.
\end{lemma}
\begin{proof}
  \uses{lem:ov-attention-output}
  Let $\ip{v_{i^*}}{v_{j^*}}=0$.  The numerator of row $i^*$ contains the term
  $\exp(0)\|v_{j^*}\|_1$.  After the preprocessing used in the proof, zero
  vectors have already been handled, so $\|v_{j^*}\|_1\geq1$.  The denominator
  is at most $n+1$ because every factor
  $\exp(-\lambda\ip{v_{i^*}}{v_k})$ is at most $1$.  Therefore the row output
  is at least $1/(n+1)$.
\end{proof}

\begin{lemma}[OV no-instance attention gap]
  \label{lem:ov-no-gap}
  \uses{lem:ov-attention-output, def:ov}
  If the OV instance has no orthogonal pair, then every attention output entry
  in the OV construction is at most $n\ell\exp(-\lambda)$.  In particular, if
  $\lambda$ is chosen so that
  $n\ell\exp(-\lambda)<1/(2(n+1))$, then every no-instance output entry is
  below $1/(2(n+1))$.
\end{lemma}
\begin{proof}
  \uses{lem:ov-attention-output}
  In a no instance, every document pair has inner product at least $1$.  Hence
  every numerator weight is at most $\exp(-\lambda)$, and every $\|v_j\|_1$ is
  at most $\ell$.  The numerator is at most $n\ell\exp(-\lambda)$, while the
  denominator is at least $1$ because of the zero sentinel row.  The final
  displayed implication is exactly the assumed upper bound on
  $n\ell\exp(-\lambda)$.
\end{proof}

\begin{theorem}[A single attention unit solves OV, Theorem 4.1]
  \label{thm:transformer-solves-ov}
  \uses{def:sentinel-transformer-solves, def:ov, lem:continuous-threshold-or}
  An attention unit with input and output MLPs and parameters
  $d=\ell$, $d_{\rm in}=\ell$, $d_{\rm out}=1$, $m\geq\ell+1$, and one fixed
  raw sentinel row $0^\ell$ solves $\OV_{n,\ell}$ in the
  sentinel-augmented sense of Definition~\ref{def:sentinel-transformer-solves}.
\end{theorem}
\begin{proof}
  \uses{def:sentinel-augmented-input, lem:attention-factorization, lem:ov-attention-output, lem:ov-yes-gap, lem:ov-no-gap, lem:continuous-threshold-or}
  Evaluate the transformer on $\Aug_{0^\ell}(E)$, whose last row is the fixed
  zero sentinel and whose first $n$ rows are exactly the OV instance vectors.
  Choose $Q,K$ so that
  $QK^\top=-\lambda I_\ell$ for a scale $\lambda>0$ satisfying
  $n\ell\exp(-\lambda)<1/(2(n+1))$, and choose $V$ to be the all-one column.
  Lemma~\ref{lem:ov-attention-output} gives the row outputs.  By
  Lemma~\ref{lem:ov-yes-gap}, yes instances produce an entry at least
  $1/(n+1)$; by Lemma~\ref{lem:ov-no-gap}, no instances produce only entries
  below $1/(2(n+1))$.  The output MLP from
  Lemma~\ref{lem:continuous-threshold-or}, with
  $a=1/(2(n+1))$ and $b=1/(n+1)$, separates these two ranges and outputs $1$ or
  $0$ accordingly.
\end{proof}

\begin{lemma}[Max-IP and Min-IP attention gap]
  \label{lem:maxip-minip-attention-gap}
  \uses{def:attention, def:max-ip-decision, def:min-ip-decision, lem:append-end-token-mlp, def:sentinel-augmented-input}
  For $1\leq t\leq\ell$, evaluate the transformer on
  $\Aug_{s_t}(E)$, where $s_t=(0,\ldots,0,t+1)\in\R^\ell$ is the raw sentinel
  row from Lemma~\ref{lem:append-end-token-mlp}.  After the input MLP, the
  first $n$ rows are $x_i=(v_i,1)$ and the last row is
  $\tau_t=(0,\ldots,0,t+1)\in\R^{\ell+1}$.  With the attention parameters of
  Appendix C, choose a scale $\lambda>0$ with
  $n\ell\exp(-\lambda)<1/(2(n+1))$.  The attention output has a gap of at least
  $1/(n+1)$ versus below $1/(2(n+1))$ between yes and no instances of
  $\MaxIP_{n,\ell,t}$ and $\MinIP_{n,\ell,t}$.
\end{lemma}
\begin{proof}
  \uses{def:attention, lem:append-end-token-mlp, lem:attention-factorization, def:sentinel-augmented-input}
  Lemma~\ref{lem:append-end-token-mlp} maps each binary document row to
  $x_i=(v_i,1)$ and maps the fixed raw sentinel row to
  $\tau_t=(0,\ldots,0,t+1)$.  Definition~\ref{def:sentinel-augmented-input}
  ensures that the sentinel is present as the unique $(n+1)$st row and is not a
  candidate pair member.

  For Max-IP, choose $QK^\top=\lambda I$ and value vector
  $(1,\ldots,1,0)$.  If some document pair has inner product at least $t$, then
  the shifted score $\ip{x_i}{x_j}$ is at least $t+1$, matching or exceeding the
  sentinel score $\ip{x_i}{\tau_t}=t+1$.  The best document contribution in
  that row therefore carries at least a $1/(n+1)$ softmax fraction and has
  value at least $1$.  If all document pairs have inner product at most $t-1$,
  every document score is lower than the sentinel by at least one, hence by a
  factor $\exp(\lambda)$ after exponentiation.  The total document value is at
  most $n\ell\exp(-\lambda)$, which is below $1/(2(n+1))$ by the chosen scale.

  For Min-IP, choose $QK^\top=-\lambda I$ and use the same positive target
  sentinel row $\tau_t$.  A yes witness with $\ip{v_i}{v_j}\leq t$ has shifted
  score at most $t+1$, so after negation its exponent is at least the sentinel
  exponent; a no instance has all shifted document scores at least $t+2$, so
  every document exponent is smaller than the sentinel exponent by a factor
  $\exp(\lambda)$.  The same numerator bound gives the stated gap.
\end{proof}

\begin{theorem}[A single attention unit solves Max-IP and Min-IP decisions, Theorem C.1]
  \label{thm:transformer-solves-maxip-minip}
  \uses{def:sentinel-transformer-solves, def:max-ip-decision, def:min-ip-decision}
  An attention unit with input and output MLPs and parameters
  $d=\ell$, $d_{\rm in}=\ell+1$, $d_{\rm out}=1$, $m\geq\ell+1$, and the raw
  sentinel row $s_t=(0,\ldots,0,t+1)\in\R^\ell$ solves
  $\MaxIP_{n,\ell,t}$ and $\MinIP_{n,\ell,t}$ for $1\leq t\leq\ell$ in the
  sentinel-augmented sense of Definition~\ref{def:sentinel-transformer-solves}.
\end{theorem}
\begin{proof}
  \uses{lem:maxip-minip-attention-gap, lem:continuous-threshold-or, def:sentinel-transformer-solves}
  Lemma~\ref{lem:maxip-minip-attention-gap} constructs attention parameters and
  input MLPs that separate yes and no instances by the numerical gap
  $1/(2(n+1))<1/(n+1)$.  Applying the continuous threshold OR of
  Lemma~\ref{lem:continuous-threshold-or} to the attention outputs gives the
  required decision value.
\end{proof}

\begin{lemma}[Finite cosine thresholds have a positive no-instance gap]
  \label{lem:finite-cosine-threshold-gap}
  \uses{def:cosine-similarity, def:msd-decision, def:lsd-decision}
  Fix $n,\ell$ and $t\in[0,1]$.  Among nonzero binary vectors in
  $\{0,1\}^{\ell}$, the set of possible cosine similarities is finite.
  Therefore, for $\MSD_{n,\ell,t}$, either a pair has similarity at least $t$,
  or there is a number $\Delta>0$ such that every pair has similarity at most
  $t-\Delta$.  Similarly, for $\LSD_{n,\ell,t}$, either a pair has similarity
  at most $t$, or there is a number $\Delta>0$ such that every pair has
  similarity at least $t+\Delta$.
\end{lemma}
\begin{proof}
  \uses{def:cosine-similarity}
  There are only finitely many nonzero binary vectors of length $\ell$, hence
  only finitely many pairwise cosine-similarity values.  If no value lies on the
  accepting side of the threshold, the finite set of distances from the
  remaining values to $t$ has a positive minimum.  This minimum is the desired
  $\Delta$ for the corresponding maximization or minimization decision.
\end{proof}

\begin{lemma}[Normalized document rows realize cosine scores]
  \label{lem:normalized-document-inner-cosine}
  \uses{lem:normalize-end-token-mlp, def:cosine-similarity}
  For nonzero binary document embeddings $v,w$, if the input MLP maps document
  rows to $x_v=(v/\|v\|,1)$ and $x_w=(w/\|w\|,1)$, then
  \[
    \ip{x_v}{x_w}=1+\frac{\ip{v}{w}}{\|v\|\|w\|}.
  \]
  If the raw sentinel row is mapped to the target attention-space row
  $\tau_t=(0,\ldots,0,t+1)$ with $\ell$ leading zeros, then
  $\ip{x_v}{\tau_t}=t+1$.
\end{lemma}
\begin{proof}
  \uses{def:cosine-similarity, def:inner-norm}
  Expanding the concatenated inner product gives
  $\ip{v/\|v\|}{w/\|w\|}+1$, which is
  $\ip{v}{w}/(\|v\|\|w\|)+1$.  The target sentinel has zeros in the document
  coordinates and last coordinate $t+1$, while $x_v$ has last coordinate $1$,
  so their inner product is $t+1$.
\end{proof}

\begin{lemma}[MSD normalized attention gap]
  \label{lem:msd-normalized-attention-gap}
  \uses{def:attention, def:msd-decision, lem:normalized-document-inner-cosine, lem:finite-cosine-threshold-gap, def:sentinel-augmented-input}
  For fixed $n,\ell,t$, evaluate the transformer on
  $\Aug_{s_t}(E)$, where $s_t=(0,\ldots,0,t+2)\in\R^\ell$ is the raw sentinel
  row from Lemma~\ref{lem:normalize-end-token-mlp}.  The input MLP maps the
  first $n$ nonzero binary rows to $(v_i/\|v_i\|,1)$ and maps the sentinel to
  $\tau_t=(0,\ldots,0,t+1)\in\R^{\ell+1}$.  Choose $Q,K$ so that
  $QK^\top=\lambda I_{\ell+1}$ for a sufficiently large $\lambda>0$, and choose
  the value vector $(1,\ldots,1,0)$.  Then yes instances of
  $\MSD_{n,\ell,t}$ produce some attention output at least $1/(n+1)$, while no
  instances produce all attention outputs below $1/(2(n+1))$.
\end{lemma}
\begin{proof}
  \uses{lem:normalize-end-token-mlp, lem:normalized-document-inner-cosine, lem:finite-cosine-threshold-gap, lem:attention-factorization, def:sentinel-augmented-input}
  Lemma~\ref{lem:normalize-end-token-mlp} gives exactly the displayed row
  images, and Definition~\ref{def:sentinel-augmented-input} places the raw
  sentinel as the unique last row.  In a yes instance choose a row $i$ with
  some witness row $j$ whose cosine similarity is at least $t$, and then choose
  a document row $j'$ maximizing the attention score against row $i$.  Its
  score is at least the sentinel score $t+1$ by
  Lemma~\ref{lem:normalized-document-inner-cosine}.  The value of a nonzero
  normalized binary row under $(1,\ldots,1,0)$ is positive and at least $1$.
  Since the chosen document row has maximal score among the $n+1$ rows up to
  ties with the sentinel, its softmax mass is at least $1/(n+1)$, so the output
  is at least $1/(n+1)$.

  In a no instance Lemma~\ref{lem:finite-cosine-threshold-gap} gives
  $\Delta>0$ such that every document score is at most $t+1-\Delta$, whereas
  the target sentinel score is $t+1$.  The numerator is at most
  $n\sqrt{\ell}\exp(\lambda(t+1-\Delta))$, and the denominator is at least
  $\exp(\lambda(t+1))$.  Choosing $\lambda$ so that
  $n\sqrt{\ell}\exp(-\lambda\Delta)<1/(2(n+1))$ gives the claimed no-case
  bound.
\end{proof}

\begin{lemma}[LSD normalized attention gap]
  \label{lem:lsd-normalized-attention-gap}
  \uses{def:attention, def:lsd-decision, lem:normalized-document-inner-cosine, lem:finite-cosine-threshold-gap, def:sentinel-augmented-input}
  For fixed $n,\ell,t$, use the same augmented raw sentinel row
  $s_t=(0,\ldots,0,t+2)$, normalized document rows, and positive target
  sentinel row $\tau_t=(0,\ldots,0,t+1)$ as in
  Lemma~\ref{lem:msd-normalized-attention-gap}.  Choose $Q,K$ so that
  $QK^\top=-\lambda I_{\ell+1}$ for a sufficiently large $\lambda>0$, and
  choose the value vector $(1,\ldots,1,0)$.  Then yes instances of
  $\LSD_{n,\ell,t}$ produce some attention output at least $1/(n+1)$, while no
  instances produce all attention outputs below $1/(2(n+1))$.
\end{lemma}
\begin{proof}
  \uses{lem:normalize-end-token-mlp, lem:normalized-document-inner-cosine, lem:finite-cosine-threshold-gap, lem:attention-factorization, def:sentinel-augmented-input}
  Multiplication by $-\lambda I_{\ell+1}$ reverses the score order.  In a yes
  instance a witness pair has cosine similarity at most $t$, so its negated
  document score is at least the negated target-sentinel score.  Choosing the
  best such document score in the row gives softmax mass at least $1/(n+1)$
  and value at least $1$, as in
  Lemma~\ref{lem:msd-normalized-attention-gap}.

  In a no instance every cosine similarity is at least $t+\Delta$ for the
  positive gap from Lemma~\ref{lem:finite-cosine-threshold-gap}.  Thus every
  negated document score is at most the negated target-sentinel score minus
  $\lambda\Delta$.  Taking $\lambda$ large enough makes the total document
  contribution smaller than $1/(2(n+1))$, exactly as in the MSD no-case.
\end{proof}

\begin{theorem}[A single attention unit solves MSD and LSD decisions, Theorem C.2]
  \label{thm:transformer-solves-msd-lsd}
  \uses{def:sentinel-transformer-solves, def:msd-decision, def:lsd-decision}
  An attention unit with input and output MLPs and parameters
  $d=\ell$, $d_{\rm in}=\ell+1$, $d_{\rm out}=1$, $m\geq\ell+1$, and the raw
  sentinel row $s_t=(0,\ldots,0,t+2)\in\R^\ell$ solves $\MSD_{n,\ell,t}$ and
  $\LSD_{n,\ell,t}$ for every $t\in[0,1]$ in the sentinel-augmented sense of
  Definition~\ref{def:sentinel-transformer-solves}.
\end{theorem}
\begin{proof}
  \uses{lem:normalize-end-token-mlp, lem:msd-normalized-attention-gap, lem:lsd-normalized-attention-gap, lem:continuous-threshold-or, def:sentinel-transformer-solves}
  For $t=0$, MSD is trivial on instances with at least two nonzero document
  rows, so the output MLP may constantly output $1$.  For all other MSD cases,
  Lemma~\ref{lem:msd-normalized-attention-gap} constructs attention parameters
  whose outputs have a gap between yes and no instances.  For LSD, including
  the endpoint $t=0$, Lemma~\ref{lem:lsd-normalized-attention-gap} gives the
  analogous gap with the sign of $QK^\top$ reversed.  Applying the continuous
  threshold OR of Lemma~\ref{lem:continuous-threshold-or} to the resulting
  attention-output vector gives the scalar decision value.
\end{proof}

\begin{theorem}[Introductory transformer representation theorem, Theorem 1.4]
  \label{thm:intro-transformer-strength}
  \uses{thm:transformer-solves-ov, thm:transformer-solves-msd-lsd, def:sentinel-transformer-solves}
  A single unit of standard attention with input and output MLPs and embedding
  dimension $\ell+1$ solves $\OV_{n,\ell}$, $\MSD_{n,\ell,t}$, and
  $\LSD_{n,\ell,t}$ for every $0\leq t\leq1$ under the explicit
  sentinel-augmented input encoding of
  Definition~\ref{def:sentinel-transformer-solves}.
\end{theorem}
\begin{proof}
  \uses{thm:transformer-solves-ov, thm:transformer-solves-msd-lsd, def:sentinel-transformer-solves}
  The OV part is Theorem~\ref{thm:transformer-solves-ov}.  The MSD and LSD
  decision parts are Theorem~\ref{thm:transformer-solves-msd-lsd}.  These
  constructions use one attention unit plus simple continuous MLPs, with the
  fixed sentinel row supplied by the input encoding rather than counted as a
  problem-instance document row.
\end{proof}
